
\Data\ID{1103018801}
\Updated{2021-04-07 03:04:38}
\Level{A}
\SubArea{Lines and planes: distances and angles}
\IMG{1103018801.pdf}
\LANGen{
\Question{
Choose  the correct verbal description  of the angle shown in the picture:}
{The angle between two space diagonals of a cube.}{The angle between a space diagonal of a cube and its edge.}{The angle between two face diagonals of a cube.}{The angle between a space diagonal of a cube and a face diagonal.}{}{}}
\LANGpl{
\Question{
Wybierz prawidłowy opis kąta pokazanego na rysunku.}
{Kąt miedzy dwoma przekątnymi   sześcianu.}{Kąt miedzy dwoma przekątnymi  sześcianu, a jego  krawędzią.}{Kąt miedzy dwoma przekątnymi  ścian  sześcianu.}{Kat między przekątną sześcianu, a przekątną jednej z jego ścian.}{}{}}
\LANGcs{
\Question{
Úhel vyznačený na obrázku znázorňuje:}
{Odchylku dvou tělesových úhlopříček.}{Odchylku tělesové úhlopříčky a hrany krychle.}{Odchylku dvou stěnových úhlopříček.}{Odchylku tělesové úhlopříčky a stěnové úhlopříčky.}{}{}}
\LANGsk{
\Question{
Uhol vyznačený na obrázku znázorňuje:}
{Odchýlku dvoch telesových uhlopriečok.}{Odchýlku telesovej uhlopriečky a hrany kocky.}{Odchýlku dvoch stenových uhlopriečok.}{Odchýlku telesovej uhlopriečky a stenovej uhlopriečky.}{}{}}
\LANGes{
\Question{
El ángulo mostrado en la figura muestra:}
{El ángulo entre dos diagonales espaciales de un cubo.}{El ángulo entre una diagonal espacial de un cubo y su arista.}{El ángulo entre dos diagonales de cara de un cubo.}{El ángulo entre una diagonal espacial de un cubo y una diagonal de cara.}{}{}}
\End

\Data\ID{1103018802}
\Updated{2020-07-15 03:07:12}
\Level{A}
\SubArea{Lines and planes: distances and angles}
\IMG{1103018802.pdf}
\LANGen{
\Question{
Choose  the correct verbal description  of the angle shown in the picture.}
{The angle between a space diagonal of a cube and its face diagonal.}{The angle between a space diagonal of a cube and its edge.}{The angle between two face diagonals of a cube.}{The angle between a face diagonal of a cube and its edge.}{}{}}
\LANGpl{
\Question{
Wybierz prawidłowy opis kąta pokazanego na rysunku.}
{Kat między przekątną sześcianu, a przekątną jednej z jego ścian.}{Kąt miedzy  przekątną  sześcianu, a jego  krawędzią.}{Kąt miedzy dwoma przekątnymi  ścian  sześcianu.}{Kąt miedzy  przekątną  ściany  sześcianu, a jego  krawędzią.}{}{}}
\LANGcs{
\Question{
Úhel vyznačený na obrázku znázorňuje:}
{Odchylku tělesové úhlopříčky a stěnové úhlopříčky.}{Odchylku tělesové úhlopříčky a hrany krychle.}{Odchylku dvou stěnových úhlopříček.}{Odchylku stěnové úhlopříčky a hrany krychle.}{}{}}
\LANGsk{
\Question{
Uhol vyznačený na obrázku znázorňuje:}
{Odchýlku telesovej uhlopriečky a stenovej uhlopriečky.}{Odchýlku telesovej uhlopriečky a hrany kocky.}{Odchýlku dvoch stenových uhlopriečok.}{Odchýlku stenovej uhlopriečky a hrany kocky.}{}{}}
\LANGes{
\Question{
El ángulo mostrado en la figura muestra:}
{El ángulo entre una diagonal espacial de un cubo y una diagonal de cara.}{El ángulo entre una diagonal espacial de un cubo y su arista.}{El ángulo entre dos diagonales de cara de un cubo.}{El ángulo entre una arista y una diagonal de cara.}{}{}}
\End

\Data\ID{1103018803}
\Updated{2020-07-15 03:07:12}
\Level{A}
\SubArea{Lines and planes: distances and angles}
\IMG{1103018803.pdf}
\LANGen{
\Question{
Choose  the correct verbal description  of the angle shown in the picture.}
{The angle between an edge and a face diagonal.}{The angle between two edges of a cube.}{The angle between a space diagonal of a cube and its edge.}{The angle between two face diagonals of a cube.}{}{}}
\LANGpl{
\Question{
Wybierz prawidłowy opis kąta pokazanego na rysunku.}
{Kąt między krawędzią, a przekątną ściany bocznej.}{Kąt między dwoma krawędziami sześcianu.}{Kąt między przekątną sześcianu, a jego krawędzią.}{Kąt między dwoma przekątnymi ścian bocznych sześcianu.}{}{}}
\LANGcs{
\Question{
Úhel vyznačený na obrázku znázorňuje:}
{Odchylku hrany krychle a stěnové úhlopříčky.}{Odchylku dvou sousedních hran krychle.}{Odchylku tělesové úhlopříčky a hrany krychle.}{Odchylku dvou stěnových úhlopříček.}{}{}}
\LANGsk{
\Question{
Uhol vyznačený na obrázku znázorňuje:}
{Odchýlku hrany kocky a stenovej uhlopriečky.}{Odchýlku dvoch susedených hrán kocky.}{Odchýlku telesovej uhlopriečky a hrany kocky.}{Odchýlku dvoch stenových uhlopriečok.}{}{}}
\LANGes{
\Question{
El ángulo mostrado en la figura muestra:}
{El ángulo entre una arista y una diagonal de cara.}{El ángulo entre dos aristas de un cubo.}{El ángulo entre una diagonal espacial de un cubo y su arista.}{El ángulo entre dos diagonales de cara de un cubo.}{}{}}
\End

\Data\ID{1103018804}
\Updated{2020-07-15 03:07:12}
\Level{A}
\SubArea{Lines and planes: distances and angles}
\IMG{1103018804.pdf}
\LANGen{
\Question{
Choose the correct verbal description of the angle shown in the picture, where the point \(S_{EF}\) is the midpoint of the edge \(EF\).}
{The angle between the line \(AS_{EF}\) and the plane \(BCG\) (right side face).}{The angle between the line \(AS_{EF}\) and the plane \(EFG\) (top face).}{The angle between the line \(AS_{EF}\) and the plane \(DCG\) (back face).}{The angle between the line \(AS_{EF}\) and the plane \(ABF\) (front face).}{}{}}
\LANGpl{
\Question{
Wybierz prawidłowy opis kąta pokazanego na rysunku, gdzie punkt \(S_{EF}\) jest  środkiem krawędzi  \(EF\).}
{Kąt między prostą  \(AS_{EF}\), a płaszczyzną  \(BCG\) (prawa ściana).}{Kąt między prostą \(AS_{EF}\), a płaszczyzną \(EFG\) (górna podstawa).}{Kąt między prostą \(AS_{EF}\), a płaszczyzną \(DCG\) (tylna ściana).}{Kąt między prostą \(AS_{EF}\), a płaszczyzną \(ABF\) (ściana frontowa).}{}{}}
\LANGcs{
\Question{
Úhel vyznačený na obrázku znázorňuje (bod \(S_{EF}\) je střed úsečky \(EF\)):}
{Odchylku přímky \(AS_{EF}\) a roviny \(BCG\) (boční stěny).}{Odchylku přímky \(AS_{EF}\) a roviny \(EFG\) (horní podstavy).}{Odchylku přímky \(AS_{EF}\) a roviny \(DCG\) (zadní stěna).}{Odchylku přímky \(AS_{EF}\) a roviny \(ABF\) (přední stěna).}{}{}}
\LANGsk{
\Question{
Uhol vyznačený na obrázku znázorňuje (bod \(S_{EF}\) je stred úsečky \(EF\)):}
{Odchýlku priamky \(AS_{EF}\) a roviny \(BCG\) (bočnej steny).}{Odchýlku priamky \(AS_{EF}\) a roviny \(EFG\) (hornej podstavy).}{Odchýlku priamky \(AS_{EF}\) and roviny \(DCG\) (zadnej steny).}{Odchýlku priamky \(AS_{EF}\) a roviny \(ABF\) (prednej steny).}{}{}}
\LANGes{
\Question{
El ángulo mostrado en la figura muestra (el punto \(S_{EF}\) es el centro del segmento\(EF\)):}
{El ángulo entre la recta \(AS_{EF}\) y el plano \(BCG\) (cara lateral derecha).}{El ángulo entre la recta \(AS_{EF}\) y el plano \(EFG\) (cara superior).}{El ángulo entre la recta \(AS_{EF}\) y el plano \(DCG\) (cara posterior).}{El ángulo entre la recta \(AS_{EF}\) y el plano \(ABF\) (cara frontal).}{}{}}
\End

\Data\ID{1103018901}
\Updated{2020-07-15 03:07:12}
\Level{A}
\SubArea{Lines and planes: distances and angles}
\IMG{1103018901.pdf}
\LANGen{
\Question{
In the cube \( ABCDEFGH \), find the angle between the space diagonals \( AG \) and \( BH \). Round the result to two decimal places.}
{\( 70.53^{\circ} \)}{\( 45^{\circ} \)}{\( 54.74^{\circ} \)}{\( 35.27^{\circ} \)}{}{}}
\LANGpl{
\Question{
Dany jest sześcian \( ABCDEFGH \), wyznacz kąt pomiędzy jego przekątnymi \( AG \) i \( BH \). Zaokrągli wynik do dwóch miejsc po przecinku.}
{\( 70{,}53^{\circ} \)}{\( 45^{\circ} \)}{\( 54{,}74^{\circ} \)}{\( 35{,}27^{\circ} \)}{}{}}
\LANGcs{
\Question{
V krychli \( ABCDEFGH \) určete odchylku tělesových úhlopříček \( AG \) a \( BH \). Výsledek zaokrouhlete na dvě desetinná místa.}
{\( 70{,}53^{\circ} \)}{\( 45^{\circ} \)}{\( 54{,}74^{\circ} \)}{\( 35{,}27^{\circ} \)}{}{}}
\LANGsk{
\Question{
V kocke \( ABCDEFGH \) určte odchýlku telesových uhlopriečok \( AG \) a \( BH \). Výsledok zaokrúhlite na dve desatinné miesta.}
{\( 70{,}53^{\circ} \)}{\( 45^{\circ} \)}{\( 54{,}74^{\circ} \)}{\( 35{,}27^{\circ} \)}{}{}}
\LANGes{
\Question{
En el cubo\( ABCDEFGH \), determina el ángulo entre las diagonales espaciales \( AG \) y \( BH \). Redondea a dos posiciones decimales.}
{\( 70.53^{\circ} \)}{\( 45^{\circ} \)}{\( 54.74^{\circ} \)}{\( 35.27^{\circ} \)}{}{}}
\End

\Data\ID{1103018902}
\Updated{2020-07-15 03:07:12}
\Level{A}
\SubArea{Lines and planes: distances and angles}
\IMG{1103018902.pdf}
\LANGen{
\Question{
Let \( \varphi \) bet the angle between a space diagonal of a cube and its face diagonal.  Choose  the correct expression for \( \varphi \).}
{\( \mathrm{tg}\,\varphi = \frac{\sqrt2}2 \)}{\( \mathrm{sin}\,\varphi = \frac{\sqrt2}2 \)}{\( \mathrm{tg}\,\varphi = \frac{\sqrt3}3 \)}{\( \mathrm{tg}\,\varphi = \frac{\sqrt6}3 \)}{}{}}
\LANGpl{
\Question{
Dany jest kąt  \( \varphi \) między  przekątną  sześcianu, a przekątną ściany bocznej.   Wybierz poprawne wyrażenie dla \( \varphi \).}
{\( \mathrm{tg}\,\varphi = \frac{\sqrt2}2 \)}{\( \mathrm{sin}\,\varphi = \frac{\sqrt2}2 \)}{\( \mathrm{tg}\,\varphi = \frac{\sqrt3}3 \)}{\( \mathrm{tg}\,\varphi = \frac{\sqrt6}3 \)}{}{}}
\LANGcs{
\Question{
Je dána krychle \( ABCDEFGH \) s hranou délky \( a \). Vyberte vztah, který platí pro odchylku \( \varphi \) tělesové a stěnové úhlopříčky.}
{\( \mathrm{tg}\,\varphi = \frac{\sqrt2}2 \)}{\( \mathrm{sin}\,\varphi = \frac{\sqrt2}2 \)}{\( \mathrm{tg}\,\varphi = \frac{\sqrt3}3 \)}{\( \mathrm{tg}\,\varphi = \frac{\sqrt6}3 \)}{}{}}
\LANGsk{
\Question{
Je daná kocka \( ABCDEFGH \) s hranou dĺžky \( a \). Vyberte vzťah, ktorý platí pre odchýlku \( \varphi \) telesovej a stenovej uhlopriečky.}
{\( \mathrm{tg}\,\varphi = \frac{\sqrt2}2 \)}{\( \mathrm{sin}\,\varphi = \frac{\sqrt2}2 \)}{\( \mathrm{tg}\,\varphi = \frac{\sqrt3}3 \)}{\( \mathrm{tg}\,\varphi = \frac{\sqrt6}3 \)}{}{}}
\LANGes{
\Question{
Dado el cubo \( ABCDEFGH\) con una arista \( a\). Elige la expresión correcta para el ángulo \(\varphi\) que se encuentra entre una diagonal espacial y una diagonal de cara.}
{\( \mathrm{tg}\,\varphi = \frac{\sqrt2}2 \)}{\( \mathrm{sin}\,\varphi = \frac{\sqrt2}2 \)}{\( \mathrm{tg}\,\varphi = \frac{\sqrt3}3 \)}{\( \mathrm{tg}\,\varphi = \frac{\sqrt6}3 \)}{}{}}
\End

\Data\ID{1103018903}
\Updated{2020-07-15 03:07:12}
\Level{A}
\SubArea{Lines and planes: distances and angles}
\IMG{1103018903.pdf}
\LANGen{
\Question{
In the cube \( ABCDEFGH \) with \( S_{AC} \) being the midpoint of the diagonal \( AC \), let \( \varphi \) be the angle between the line \( ES_{AC} \) and the bottom face \( ABCD \). Choose the correct expression for \( \varphi \).}
{\( \mathrm{tg}\,\varphi = \sqrt2  \)}{\( \mathrm{sin}\,\varphi = \frac{\sqrt2}3  \)}{\( \mathrm{tg}\,\varphi = \frac{\sqrt2}2  \)}{\( \mathrm{cos}\,\varphi = \frac{\sqrt6}3  \)}{}{}}
\LANGpl{
\Question{
Dany jest sześcian \( ABCDEFGH \), gdzie  \( S_{AC} \) jest środkiem przekątnej \( AC \), kąt \( \varphi \) to kąt między prostą \( ES_{AC} \), a podstawą sześcianu \( ABCD \).  Wybierz poprawne wyrażenie dla \( \varphi \).}
{\( \mathrm{tg}\,\varphi = \sqrt2  \)}{\( \mathrm{sin}\,\varphi = \frac{\sqrt2}3  \)}{\( \mathrm{tg}\,\varphi = \frac{\sqrt2}2  \)}{\( \mathrm{cos}\,\varphi = \frac{\sqrt6}3  \)}{}{}}
\LANGcs{
\Question{
Je dána krychle \( ABCDEFGH \) s hranou délky \( a \) a bod \( S_{AC} \) střed úhlopříčky \( AC \). Vyberte vztah, který platí pro odchylku \( \varphi \) přímky \( ES_{AC} \) a roviny dolní podstavy \( ABCD \).}
{\( \mathrm{tg}\,\varphi = \sqrt2  \)}{\( \mathrm{sin}\,\varphi = \frac{\sqrt2}3  \)}{\( \mathrm{tg}\,\varphi = \frac{\sqrt2}2  \)}{\( \mathrm{cos}\,\varphi = \frac{\sqrt6}3  \)}{}{}}
\LANGsk{
\Question{
Je daná kocka \( ABCDEFGH \) s hranou dĺžky \( a \) a bod \( S_{AC} \) stred uhlopriečky \( AC \). Vyberte vzťah, ktorý platí pre odchýlku \( \varphi \) priamky \( ES_{AC} \) a roviny dolnej podstavy \( ABCD \).}
{\( \mathrm{tg}\,\varphi = \sqrt2  \)}{\( \mathrm{sin}\,\varphi = \frac{\sqrt2}3  \)}{\( \mathrm{tg}\,\varphi = \frac{\sqrt2}2  \)}{\( \mathrm{cos}\,\varphi = \frac{\sqrt6}3  \)}{}{}}
\LANGes{
\Question{
Dado el cubo \( ABCDEFGH \) con el punto \( S_{AC} \) y el centro de diagonal \( AC \). Elige la expresión correcta para el ángulo \( \varphi \) que se encuentra entre una recta \( ES_{AC} \)  y la cara inferior  \(ABCD \).}
{\( \mathrm{tg}\,\varphi = \sqrt2  \)}{\( \mathrm{sin}\,\varphi = \frac{\sqrt2}3  \)}{\( \mathrm{tg}\,\varphi = \frac{\sqrt2}2  \)}{\( \mathrm{cos}\,\varphi = \frac{\sqrt6}3  \)}{}{}}
\End

\Data\ID{1103018904}
\Updated{2020-07-15 03:07:12}
\Level{A}
\SubArea{Lines and planes: distances and angles}
\IMG{1103018904_0.pdf}
\LANGen{
\Question{
In the cube \( ABCDEFGH \), let \( S_{FG} \) be the midpoint of the edge \( FG \). Find the angle between the lines \( BS_{FG} \) and \( BF \). Round the result to two decimal places.}
{\( 26.57^{\circ} \)}{\( 22.5^{\circ} \)}{\( 45^{\circ} \)}{\( 54.74^{\circ} \)}{}{}}
\LANGpl{
\Question{
Dany jest sześcian \( ABCDEFGH \), gdzie \( S_{FG} \) to środek krawędzi  \( FG \). Wskaż kąt między prostymi \( BS_{FG} \) i \( BF \). Zaokrągli wynik do dwóch miejsc po przecinku.}
{\( 26{,}57^{\circ} \)}{\( 22{,}5^{\circ} \)}{\( 45^{\circ} \)}{\( 54{,}74^{\circ} \)}{}{}}
\LANGcs{
\Question{
Je dána krychle \( ABCDEFGH \) s hranou délky \( a \) a bod \( S_{FG} \) střed hrany \( FG \). Určete odchylku \( \varphi \) přímek \( BS_{FG} \) a \( BF \). Výsledek zaokrouhlete na dvě desetinná místa.}
{\( 26{,}57^{\circ} \)}{\( 22{,}5^{\circ} \)}{\( 45^{\circ} \)}{\( 54{,}74^{\circ} \)}{}{}}
\LANGsk{
\Question{
Je daná kocka \( ABCDEFGH \) s hranou dĺžky \( a \) a bod \( S_{FG} \) stred hrany \( FG \). Určte odchýlku \( \varphi \) priamok \( BS_{FG} \) a \( BF \). Výsledok zaokrúhlite na dve desatinné miesta.}
{\( 26{,}57^{\circ} \)}{\( 22{,}5^{\circ} \)}{\( 45^{\circ} \)}{\( 54{,}74^{\circ} \)}{}{}}
\LANGes{
\Question{
Dado el cubo \( ABCDEFGH \) con con el punto \( S_{FG} \) y el centro de la arista  \( FG \). Determina el ángulo entre la recta \( BS_{FG} \) y la recta \( BF \). Redondea a dos posiciones decimales.}
{\( 26.57^{\circ} \)}{\( 22.5^{\circ} \)}{\( 45^{\circ} \)}{\( 54.74^{\circ} \)}{}{}}
\End

\Data\ID{1103018905}
\Updated{2020-07-15 03:07:12}
\Level{A}
\SubArea{Lines and planes: distances and angles}
\IMG{1103018905.pdf}
\LANGen{
\Question{
In the cube \( ABCDEFGH \) with \( S_{AC} \) being the midpoint of the diagonal \( AC \), let \( \varphi \) be the angle between the line \( EG \) and the line \( GS_{AC} \). Choose the correct expression for \( \varphi \):}
{\( \mathrm{tg}\,\varphi = \sqrt2 \)}{\( \mathrm{sin}\,\varphi = \frac{\sqrt3}3 \)}{\( \mathrm{tg}\,\varphi = \frac{\sqrt2}2 \)}{\( \mathrm{cos}\,\varphi = \frac{\sqrt6}3 \)}{}{}}
\LANGpl{
\Question{
Dany jest sześcian \( ABCDEFGH \), gdzie \( S_{AC} \) jest środkiem przekątnej \( AC \), kąt  \( \varphi \) to kąt między prostą \( EG \), a prostą  \( GS_{AC} \). Wybierz poprawne wyrażenie dla  \( \varphi \):}
{\( \mathrm{tg}\,\varphi = \sqrt2 \)}{\( \mathrm{sin}\,\varphi = \frac{\sqrt3}3 \)}{\( \mathrm{tg}\,\varphi = \frac{\sqrt2}2 \)}{\( \mathrm{cos}\,\varphi = \frac{\sqrt6}3 \)}{}{}}
\LANGcs{
\Question{
Je dána krychle \( ABCDEFGH \) s hranou délky \( a \) a bod \( S_{AC} \) střed úhlopříčky \( AC \). Vyberte vztah, který platí pro odchylku \( \varphi \) přímek \( EG \) a \( GS_{AC} \):}
{\( \mathrm{tg}\,\varphi = \sqrt2 \)}{\( \mathrm{sin}\,\varphi = \frac{\sqrt3}3 \)}{\( \mathrm{tg}\,\varphi = \frac{\sqrt2}2 \)}{\( \mathrm{cos}\,\varphi = \frac{\sqrt6}3 \)}{}{}}
\LANGsk{
\Question{
Je daná kocka \( ABCDEFGH \) s hranou dĺžky \( a \) a bod \( S_{AC} \) stred uhlopriečky \( AC \). Vyberte vzťah, ktorý platí pre odchýlku \( \varphi \) priamok \( EG \) a \( GS_{AC} \):}
{\( \mathrm{tg}\,\varphi = \sqrt2 \)}{\( \mathrm{sin}\,\varphi = \frac{\sqrt3}3 \)}{\( \mathrm{tg}\,\varphi = \frac{\sqrt2}2 \)}{\( \mathrm{cos}\,\varphi = \frac{\sqrt6}3 \)}{}{}}
\LANGes{
\Question{
Dado el cubo \( ABCDEFGH \) con el punto \( S_{AC} \) y  el centro de la diagonal \( AC \). Elige la expresión correcta para el ángulo \( \varphi \) que se encuentra entre las rectas \( EG \) y  \( GS_{AC} \).}
{\( \mathrm{tg}\,\varphi = \sqrt2 \)}{\( \mathrm{sin}\,\varphi = \frac{\sqrt3}3 \)}{\( \mathrm{tg}\,\varphi = \frac{\sqrt2}2 \)}{\( \mathrm{cos}\,\varphi = \frac{\sqrt6}3 \)}{}{}}
\End

\Data\ID{1103018906}
\Updated{2020-07-15 03:07:12}
\Level{A}
\SubArea{Lines and planes: distances and angles}
\IMG{1103018906.pdf}
\LANGen{
\Question{
In the cube \( ABCDEFGH \),  find the angle between two face diagonals \( EG \) and \( EB \).}
{\( 60^{\circ} \)}{\( 90^{\circ} \)}{\( 45^{\circ} \)}{\( 30^{\circ} \)}{}{}}
\LANGpl{
\Question{
Dany jest sześcian \( ABCDEFGH \), wyznacz kąt między przekątnymi \( EG \) i  \( EB \).}
{\( 60^{\circ} \)}{\( 90^{\circ} \)}{\( 45^{\circ} \)}{\( 30^{\circ} \)}{}{}}
\LANGcs{
\Question{
Je dána krychle \( ABCDEFGH \) s hranou délky \( a \). Určete odchylku \( \varphi \) stěnových úhlopříček \( EG \) a \( EB \).}
{\( 60^{\circ} \)}{\( 90^{\circ} \)}{\( 45^{\circ} \)}{\( 30^{\circ} \)}{}{}}
\LANGsk{
\Question{
Je daná kocka \( ABCDEFGH \) s hranou dĺžky \( a \). Určte odchýlku \( \varphi \) stenových uhlopriečok \( EG \) a \( EB \).}
{\( 60^{\circ} \)}{\( 90^{\circ} \)}{\( 45^{\circ} \)}{\( 30^{\circ} \)}{}{}}
\LANGes{
\Question{
Dado el cubo \( ABCDEFGH \). Determina el ángulo de las diagonales de cara \( EG \) y \( EB \).}
{\( 60^{\circ} \)}{\( 90^{\circ} \)}{\( 45^{\circ} \)}{\( 30^{\circ} \)}{}{}}
\End

\Data\ID{1103055801}
\Updated{2020-07-15 03:07:12}
\Level{A}
\SubArea{Lines and planes: distances and angles}
\IMG{1103055801.pdf}
\LANGen{
\Question{
The rectangular box \( ABCDEFGH \) has edges of lengths \( |AB|=6\,\mathrm{cm} \), \( |BC|=4\,\mathrm{cm} \) and \( |AE|=8\,\mathrm{cm} \). The point \( S \) is the center  of the base \( ABCD \) as seen in the picture. Find the distance between the point \( E \) and the point \( S \).}
{\( \sqrt{77}\,\mathrm{cm} \)}{\( 3\sqrt{10}\,\mathrm{cm} \)}{\( 2\sqrt{26}\,\mathrm{cm} \)}{\( 10\,\mathrm{cm} \)}{}{}}
\LANGpl{
\Question{
Prostopadłościan  \( ABCDEFGH \) ma krawędzie o długości \( |AB|=6\,\mathrm{cm} \), \( |BC|=4\,\mathrm{cm} \) i \( |AE|=8\,\mathrm{cm} \). Punkt \( S \) to środek podstawy \( ABCD \).  Oblicz odległość pomiędzy punktem \( E \) i \( S \).}
{\( \sqrt{77}\,\mathrm{cm} \)}{\( 3\sqrt{10}\,\mathrm{cm} \)}{\( 2\sqrt{26}\,\mathrm{cm} \)}{\( 10\,\mathrm{cm} \)}{}{}}
\LANGcs{
\Question{
V kvádru \( ABCDEFGH \) platí: \( |AB|=6\,\mathrm{cm} \), \( |BC|=4\,\mathrm{cm} \), \( |AE|=8\,\mathrm{cm} \). Určete vzdálenost bodů  \( E \), \( S \), kde bod \( S \) je střed podstavy  \( ABCD \), viz obrázek.}
{\( \sqrt{77}\,\mathrm{cm} \)}{\( 3\sqrt{10}\,\mathrm{cm} \)}{\( 2\sqrt{26}\,\mathrm{cm} \)}{\( 10\,\mathrm{cm} \)}{}{}}
\LANGsk{
\Question{
V kvádri \( ABCDEFGH \) platí: \( |AB|=6\,\mathrm{cm} \), \( |BC|=4\,\mathrm{cm} \), \( |AE|=8\,\mathrm{cm} \). Určte vzdialenosť bodov  \( E \), \( S \), kde bod \( S \) je stred podstavy  \( ABCD \), viď obrázok.}
{\( \sqrt{77}\,\mathrm{cm} \)}{\( 3\sqrt{10}\,\mathrm{cm} \)}{\( 2\sqrt{26}\,\mathrm{cm} \)}{\( 10\,\mathrm{cm} \)}{}{}}
\LANGes{
\Question{
El ortoedro \( ABCDEFGH \) tiene las longitudes de las aristas  \( |AB|=6\,\mathrm{cm} \), \( |BC|=4\,\mathrm{cm} \) y \( |AE|=8\,\mathrm{cm} \). El punto \( S \) es el centro de la base \( ABCD \). Determina la distancia entre el punto  \( E \) y el punto \( S \).}
{\( \sqrt{77}\,\mathrm{cm} \)}{\( 3\sqrt{10}\,\mathrm{cm} \)}{\( 2\sqrt{26}\,\mathrm{cm} \)}{\( 10\,\mathrm{cm} \)}{}{}}
\End

\Data\ID{1103055802}
\Updated{2020-07-15 03:07:12}
\Level{A}
\SubArea{Lines and planes: distances and angles}
\IMG{1103055802.pdf}
\LANGen{
\Question{
The rectangular box \( ABCDEFGH \) has edges of lengths \( |AB|=6\,\mathrm{cm} \), \( |BC|=4\,\mathrm{cm} \) and \( |AE|=8\,\mathrm{cm} \). The point \( S \) is the center  of the left side face \( ADHE \) as seen in the picture. Find the distance between the point \( F \) and the point \( S \).}
{\( 2\sqrt{14}\,\mathrm{cm} \)}{\( 2\sqrt{29}\,\mathrm{cm} \)}{\( 2\sqrt{17}\,\mathrm{cm} \)}{\( 4\sqrt{2}\,\mathrm{cm} \)}{}{}}
\LANGpl{
\Question{
Prostopadłościan \( ABCDEFGH \) ma krawędzie o długości \( |AB|=6\,\mathrm{cm} \), \( |BC|=4\,\mathrm{cm} \) i \( |AE|=8\,\mathrm{cm} \). Punkt \( S \) to środek  ściany bocznej \( ADHE \).  Oblicz odległość pomiędzy punktem \( F \) i punktem \( S \).}
{\( 2\sqrt{14}\,\mathrm{cm} \)}{\( 2\sqrt{29}\,\mathrm{cm} \)}{\( 2\sqrt{17}\,\mathrm{cm} \)}{\( 4\sqrt{2}\,\mathrm{cm} \)}{}{}}
\LANGcs{
\Question{
V kvádru \( ABCDEFGH \) platí:  \( |AB|=6\,\mathrm{cm} \), \( |BC|=4\,\mathrm{cm} \), \( |AE|=8\,\mathrm{cm} \). Určete vzdálenost bodů \( F \), \( S \), kde bod  \( S \) je střed levé boční stěny \( ADHE \), viz obrázek.}
{\( 2\sqrt{14}\,\mathrm{cm} \)}{\( 2\sqrt{29}\,\mathrm{cm} \)}{\( 2\sqrt{17}\,\mathrm{cm} \)}{\( 4\sqrt{2}\,\mathrm{cm} \)}{}{}}
\LANGsk{
\Question{
V kvádri \( ABCDEFGH \) platí:  \( |AB|=6\,\mathrm{cm} \), \( |BC|=4\,\mathrm{cm} \), \( |AE|=8\,\mathrm{cm} \). Určte vzdialenosť bodov \( F \), \( S \), kde bod  \( S \) je stred ľavej bočnej steny \( ADHE \), viď obrázok.}
{\( 2\sqrt{14}\,\mathrm{cm} \)}{\( 2\sqrt{29}\,\mathrm{cm} \)}{\( 2\sqrt{17}\,\mathrm{cm} \)}{\( 4\sqrt{2}\,\mathrm{cm} \)}{}{}}
\LANGes{
\Question{
El ortoedro \( ABCDEFGH \) tiene las longitudes de las aristas \( |AB|=6\,\mathrm{cm} \), \( |BC|=4\,\mathrm{cm} \) y \( |AE|=8\,\mathrm{cm} \). El punto \( S \) es el centro de la cara izquierda \( ADHE \). Determina la distancia entre el punto \( F \) y el punto \( S \).}
{\( 2\sqrt{14}\,\mathrm{cm} \)}{\( 2\sqrt{29}\,\mathrm{cm} \)}{\( 2\sqrt{17}\,\mathrm{cm} \)}{\( 4\sqrt{2}\,\mathrm{cm} \)}{}{}}
\End

\Data\ID{1103055803}
\Updated{2020-07-15 03:07:12}
\Level{A}
\SubArea{Lines and planes: distances and angles}
\IMG{1103055803.pdf}
\LANGen{
\Question{
The rectangular box \( ABCDEFGH \) has edges of lengths \( |AB|=6\,\mathrm{cm} \), \( |BC|=3\,\mathrm{cm} \) and the face diagonal of the length \( |BG|=5\,\mathrm{cm} \). Find the distance between the center  of the top face \( EFGH \) and the center of the bottom base \( ABCD \).}
{\( 4\,\mathrm{cm} \)}{\( 6\,\mathrm{cm} \)}{\( 7.5\,\mathrm{cm} \)}{\( \sqrt{61}\,\mathrm{cm} \)}{}{}}
\LANGpl{
\Question{
Prostopadłościan \( ABCDEFGH \) ma krawędzie o długości \( |AB|=6\,\mathrm{cm} \), \( |BC|=3\,\mathrm{cm} \), długość przekątnej ściany bocznej jest równa \( |BG|=5\,\mathrm{cm} \). 
Oblicz odległość między środkiem podstawy górnej  \( EFGH \) i środkiem podstawy dolnej \( ABCD \).}
{\( 4\,\mathrm{cm} \)}{\( 6\,\mathrm{cm} \)}{\( 7{,}5\,\mathrm{cm} \)}{\( \sqrt{61}\,\mathrm{cm} \)}{}{}}
\LANGcs{
\Question{
V kvádru \( ABCDEFGH \) platí: \( |AB|=6\,\mathrm{cm} \), \( |BC|=3\,\mathrm{cm} \) a stěnová úhlopříčka \( |BG|=5\,\mathrm{cm} \). Určete vzdálenost středů horní a dolní podstavy (\( EFGH \) a \( ABCD \), viz obrázek).}
{\( 4\,\mathrm{cm} \)}{\( 6\,\mathrm{cm} \)}{\( 7{,}5\,\mathrm{cm} \)}{\( \sqrt{61}\,\mathrm{cm} \)}{}{}}
\LANGsk{
\Question{
V kvádri \( ABCDEFGH \) platí: \( |AB|=6\,\mathrm{cm} \), \( |BC|=3\,\mathrm{cm} \) a stenová uhlopriečka \( |BG|=5\,\mathrm{cm} \). Určte vzdialenosť stredov hornej a dolnej podstavy (\( EFGH \) a \( ABCD \), viď obrázok).}
{\( 4\,\mathrm{cm} \)}{\( 6\,\mathrm{cm} \)}{\( 7{,}5\,\mathrm{cm} \)}{\( \sqrt{61}\,\mathrm{cm} \)}{}{}}
\LANGes{
\Question{
El ortoedro \( ABCDEFGH \) tiene las longitudes de las aristas\( |AB|=6\,\mathrm{cm} \), \( |BC|=3\,\mathrm{cm} \) y la longitud de diagonal de cara \( |BG|=5\,\mathrm{cm} \). Determina la distancia entre el centro de la cara superior \( EFGH \) y el centro de la base inferior \( ABCD \).}
{\( 4\,\mathrm{cm} \)}{\( 6\,\mathrm{cm} \)}{\( 7.5\,\mathrm{cm} \)}{\( \sqrt{61}\,\mathrm{cm} \)}{}{}}
\End

\Data\ID{1103055804}
\Updated{2020-07-15 03:07:12}
\Level{A}
\SubArea{Lines and planes: distances and angles}
\IMG{1103055804.pdf}
\LANGen{
\Question{
The rectangular box \( ABCDEFGH \) shown in the picture has edges of lengths \( |AB|=6\,\mathrm{cm} \), \( |BC|=4\,\mathrm{cm} \) and \( |AE|=8\,\mathrm{cm} \). Find the distance between the lines \( AH \) and \( FC \).}
{\( 6\,\mathrm{cm} \)}{\( 2\sqrt{13}\,\mathrm{cm} \)}{\( 0 \,\mathrm{cm} \)}{\( 4\,\mathrm{cm} \)}{}{}}
\LANGpl{
\Question{
Prostopadłościan \( ABCDEFGH \) ma krawędzie o długości  \( |AB|=6\,\mathrm{cm} \), \( |BC|=4\,\mathrm{cm} \) i \( |AE|=8\,\mathrm{cm} \). Oblicz odległość między prostymi \( AH \) i \( FC \). Zobacz rysunek.}
{\( 6\,\mathrm{cm} \)}{\( 2\sqrt{13}\,\mathrm{cm} \)}{\( 0 \,\mathrm{cm} \)}{\( 4\,\mathrm{cm} \)}{}{}}
\LANGcs{
\Question{
V kvádru \( ABCDEFGH \) platí:  \( |AB|=6\,\mathrm{cm} \), \( |BC|=4\,\mathrm{cm} \) a \( |AE|=8\,\mathrm{cm} \). Určete vzdálenost přímek \( AH \) a \( FC \), viz obrázek.}
{\( 6\,\mathrm{cm} \)}{\( 2\sqrt{13}\,\mathrm{cm} \)}{\( 0 \,\mathrm{cm} \)}{\( 4\,\mathrm{cm} \)}{}{}}
\LANGsk{
\Question{
V kvádri \( ABCDEFGH \) platí:  \( |AB|=6\,\mathrm{cm} \), \( |BC|=4\,\mathrm{cm} \) a \( |AE|=8\,\mathrm{cm} \). Určte vzdialenosť priamok \( AH \) a \( FC \), viď obrázok.}
{\( 6\,\mathrm{cm} \)}{\( 2\sqrt{13}\,\mathrm{cm} \)}{\( 0 \,\mathrm{cm} \)}{\( 4\,\mathrm{cm} \)}{}{}}
\LANGes{
\Question{
El ortoedro \( ABCDEFGH \) tiene las longitudes de las aristas\( |AB|=6\,\mathrm{cm} \), \( |BC|=4\,\mathrm{cm} \) y \( |AE|=8\,\mathrm{cm} \). Determina la distancia entre las rectas \( AH \) y \( FC \).}
{\( 6\,\mathrm{cm} \)}{\( 2\sqrt{13}\,\mathrm{cm} \)}{\( 0 \,\mathrm{cm} \)}{\( 4\,\mathrm{cm} \)}{}{}}
\End

\Data\ID{1103055805}
\Updated{2020-07-15 03:07:12}
\Level{A}
\SubArea{Lines and planes: distances and angles}
\IMG{1103055805.pdf}
\LANGen{
\Question{
The rectangular box \( ABCDEFGH \) shown in the picture has edges of lengths \( |AB|=6\,\mathrm{cm} \), \( |BC|=4\,\mathrm{cm} \) and \( |AE|=8\,\mathrm{cm} \). Find the distance between the lines \( AB \) and \( HG \).}
{\( 4\sqrt5\,\mathrm{cm} \)}{\( 8\sqrt2\,\mathrm{cm} \)}{\( 10\,\mathrm{cm} \)}{\( 8\,\mathrm{cm} \)}{}{}}
\LANGpl{
\Question{
Prostopadłościan  \( ABCDEFGH \) ma krawędzie o długościach  \( |AB|=6\,\mathrm{cm} \), \( |BC|=4\,\mathrm{cm} \) i \( |AE|=8\,\mathrm{cm} \). Oblicz odległość między prostymi \( AB \) i \( HG \). Zobacz rysunek.}
{\( 4\sqrt5\,\mathrm{cm} \)}{\( 8\sqrt2\,\mathrm{cm} \)}{\( 10\,\mathrm{cm} \)}{\( 8\,\mathrm{cm} \)}{}{}}
\LANGcs{
\Question{
V kvádru \( ABCDEFGH \) platí: \( |AB|=6\,\mathrm{cm} \), \( |BC|=4\,\mathrm{cm} \), \( |AE|=8\,\mathrm{cm} \). Určete vzdálenost přímek \( AB \) a \( HG \), viz obrázek.}
{\( 4\sqrt5\,\mathrm{cm} \)}{\( 8\sqrt2\,\mathrm{cm} \)}{\( 10\,\mathrm{cm} \)}{\( 8\,\mathrm{cm} \)}{}{}}
\LANGsk{
\Question{
V kvádri \( ABCDEFGH \) platí: \( |AB|=6\,\mathrm{cm} \), \( |BC|=4\,\mathrm{cm} \), \( |AE|=8\,\mathrm{cm} \). Určte vzdialenosť priamok \( AB \) a \( HG \), viď obrázok.}
{\( 4\sqrt5\,\mathrm{cm} \)}{\( 8\sqrt2\,\mathrm{cm} \)}{\( 10\,\mathrm{cm} \)}{\( 8\,\mathrm{cm} \)}{}{}}
\LANGes{
\Question{
El ortoedro \( ABCDEFGH \) tiene las longitudes de las aristas\( |AB|=6\,\mathrm{cm} \), \( |BC|=4\,\mathrm{cm} \) y \( |AE|=8\,\mathrm{cm} \). Determina la distancia entre las rectas \( AB \) y \( HG \).}
{\( 4\sqrt5\,\mathrm{cm} \)}{\( 8\sqrt2\,\mathrm{cm} \)}{\( 10\,\mathrm{cm} \)}{\( 8\,\mathrm{cm} \)}{}{}}
\End

\Data\ID{1103055806}
\Updated{2020-07-15 03:07:12}
\Level{A}
\SubArea{Lines and planes: distances and angles}
\IMG{1103055806_0.pdf}
\LANGen{
\Question{
The rectangular box \( ABCDEFGH \) shown in the picture has edges of lengths \( |AB|=6\,\mathrm{cm} \), \( |BC|=4\,\mathrm{cm} \) and \( |AE|=8\,\mathrm{cm} \). Find the distance between the lines \(ES_{FG} \) and \( DS_{BC} \), where \( S_{FG} \) is the  center of  \(FG\) and \( S_{BC} \) is the  center of  \(BC\).}
{\( 8 \,\mathrm{cm} \)}{\( 10\,\mathrm{cm} \)}{\( 4\sqrt5\,\mathrm{cm} \)}{\( 8\sqrt2\,\mathrm{cm} \)}{}{}}
\LANGpl{
\Question{
Prostopadłościan  \( ABCDEFGH \) ma krawędzie o długościach  \( |AB|=6\,\mathrm{cm} \), \( |BC|=4\,\mathrm{cm} \) i \( |AE|=8\,\mathrm{cm} \). Oblicz odległość między prostymi \(ES_{FG} \) i \( DS_{BC} \), jeśli \( S_{FG} \) to środek  \(FG\) i \( S_{BC} \) to środek  \(BC\).  Zobacz rysunek.}
{\( 8 \,\mathrm{cm} \)}{\( 10\,\mathrm{cm} \)}{\( 4\sqrt5\,\mathrm{cm} \)}{\( 8\sqrt2\,\mathrm{cm} \)}{}{}}
\LANGcs{
\Question{
V kvádru \( ABCDEFGH \) platí:  \( |AB|=6\,\mathrm{cm} \), \( |BC|=4\,\mathrm{cm} \), \( |AE|=8\,\mathrm{cm} \). Určete vzdálenost přímek \(ES_{FG} \) a \( DS_{BC} \), kde bod \( S_{FG} \) je střed  hrany \(FG\) a bod \( S_{BC} \) je střed hrany \(BC\), viz obrázek.}
{\( 8 \,\mathrm{cm} \)}{\( 10\,\mathrm{cm} \)}{\( 4\sqrt5\,\mathrm{cm} \)}{\( 8\sqrt2\,\mathrm{cm} \)}{}{}}
\LANGsk{
\Question{
V kvádri \( ABCDEFGH \) platí:  \( |AB|=6\,\mathrm{cm} \), \( |BC|=4\,\mathrm{cm} \), \( |AE|=8\,\mathrm{cm} \). Určte vzdialenosť priamok \(ES_{FG} \) a \( DS_{BC} \), kde bod \( S_{FG} \) je stred  hrany \(FG\) a bod \( S_{BC} \) je stred hrany \(BC\), viď obrázok.}
{\( 8 \,\mathrm{cm} \)}{\( 10\,\mathrm{cm} \)}{\( 4\sqrt5\,\mathrm{cm} \)}{\( 8\sqrt2\,\mathrm{cm} \)}{}{}}
\LANGes{
\Question{
El ortoedro\( ABCDEFGH \) tiene las longitudes de las aristas \( |AB|=6\,\mathrm{cm} \), \( |BC|=4\,\mathrm{cm} \) y \( |AE|=8\,\mathrm{cm} \). Determina la distancia entre las rectas \(ES_{FG} \) y \( DS_{BC} \), donde \( S_{FG} \) es el centro de \(FG\) y \( S_{BC} \) es el centro de \(BC\).}
{\( 8 \,\mathrm{cm} \)}{\( 10\,\mathrm{cm} \)}{\( 4\sqrt5\,\mathrm{cm} \)}{\( 8\sqrt2\,\mathrm{cm} \)}{}{}}
\End

\Data\ID{1103056001}
\Updated{2020-07-15 03:07:21}
\Level{A}
\SubArea{Lines and planes: distances and angles}
\IMG{1103056001.pdf}
\LANGen{
\Question{
The cube \( ABCDEFGH \) shown in the picture has edges of length \( a \). Find the distance between the points \( E \) and  \( C \). Choose the correct expression for this distance:}
{\( a\sqrt3 \)}{\( \sqrt[3]3 \)}{\( a\sqrt2 \)}{\( \sqrt3 \)}{}{}}
\LANGpl{
\Question{
Sześcian \( ABCDEFGH \) przedstawiony na rysunku ma krawędzie o długości  \( a \). Oblicz odległość pomiędzy punktami \( E \) i  \( C \). Wybierz poprawne wyrażenie dla tej odległości:}
{\( a\sqrt3 \)}{\( \sqrt[3]3 \)}{\( a\sqrt2 \)}{\( \sqrt3 \)}{}{}}
\LANGcs{
\Question{
Je dána krychle \( ABCDEFGH \) s hranou délky \( a \). Vyberte vztah, který platí pro vzdálenost bodů \( E \) a  \( C \):}
{\( a\sqrt3 \)}{\( \sqrt[3]3 \)}{\( a\sqrt2 \)}{\( \sqrt3 \)}{}{}}
\LANGsk{
\Question{
Je daná kocka \( ABCDEFGH \) s hranou dĺžky \( a \). Vyberte vzťah, ktorý platí pre vzdialenosť bodov \( E \) a  \( C \):}
{\( a\sqrt3 \)}{\( \sqrt[3]3 \)}{\( a\sqrt2 \)}{\( \sqrt3 \)}{}{}}
\LANGes{
\Question{
El cubo \( ABCDEFGH \) tiene la longitud de la arista \( a \).  Elige la expresión correcta para la distancia entre los puntos \( E \) y  \( C \).}
{\( a\sqrt3 \)}{\( \sqrt[3]3 \)}{\( a\sqrt2 \)}{\( \sqrt3 \)}{}{}}
\End

\Data\ID{1103056002}
\Updated{2020-07-15 03:07:21}
\Level{A}
\SubArea{Lines and planes: distances and angles}
\IMG{1103056002.pdf}
\LANGen{
\Question{
The cube \( ABCDEFGH \) shown in the picture has edges of length \( a=6\,\mathrm{cm} \). Let \( S \) be the midpoint of the base \( ABCD \). Find the distance between the points \( H \) and \( S \).}
{\( 3\sqrt6\,\mathrm{cm} \)}{\( 6\sqrt5\,\mathrm{cm} \)}{\( 3\sqrt5\,\mathrm{cm} \)}{\( 6\sqrt3\,\mathrm{cm} \)}{}{}}
\LANGpl{
\Question{
Sześcian \( ABCDEFGH \) przedstawiony na rysunku ma krawędzie o długości \( a=6\,\mathrm{cm} \). Punkt \( S \) to środek podstawy \( ABCD \). Oblicz odległość pomiędzy \( H \) i \( S \).}
{\( 3\sqrt6\,\mathrm{cm} \)}{\( 6\sqrt5\,\mathrm{cm} \)}{\( 3\sqrt5\,\mathrm{cm} \)}{\( 6\sqrt3\,\mathrm{cm} \)}{}{}}
\LANGcs{
\Question{
Je dána krychle  \( ABCDEFGH \) s hranou délky \( a=6\,\mathrm{cm} \). Určete vzdálenost bodů  \( H \), \( S \), kde \( S \) je střed dolní podstavy \( ABCD \) (viz obrázek).}
{\( 3\sqrt6\,\mathrm{cm} \)}{\( 6\sqrt5\,\mathrm{cm} \)}{\( 3\sqrt5\,\mathrm{cm} \)}{\( 6\sqrt3\,\mathrm{cm} \)}{}{}}
\LANGsk{
\Question{
Je daná kocka  \( ABCDEFGH \) s hranou dĺžky \( a=6\,\mathrm{cm} \). Určte vzdialenosť bodov  \( H \), \( S \), kde \( S \) je stred dolnej podstavy \( ABCD \) (viď obrázok).}
{\( 3\sqrt6\,\mathrm{cm} \)}{\( 6\sqrt5\,\mathrm{cm} \)}{\( 3\sqrt5\,\mathrm{cm} \)}{\( 6\sqrt3\,\mathrm{cm} \)}{}{}}
\LANGes{
\Question{
El cubo \( ABCDEFGH \) tiene la longitud de la arista \( a=6\,\mathrm{cm} \). Determina la distancia entre los puntos \( H \) y \( S \) donde \( S \) es el centro de la base inferior \( ABCD \).}
{\( 3\sqrt6\,\mathrm{cm} \)}{\( 6\sqrt5\,\mathrm{cm} \)}{\( 3\sqrt5\,\mathrm{cm} \)}{\( 6\sqrt3\,\mathrm{cm} \)}{}{}}
\End

\Data\ID{1103056003}
\Updated{2020-07-15 03:07:21}
\Level{A}
\SubArea{Lines and planes: distances and angles}
\IMG{1103056003.pdf}
\LANGen{
\Question{
The  cube \( ABCDEFGH \) shown in the picture has edges of length \( a=6\,\mathrm{cm} \). Let \( S \) be the midpoint of the edge \( FG \). Find the distance between the points \( E \) and \( S \).}
{\( 3\sqrt5\,\mathrm{cm} \)}{\( 6\sqrt5\,\mathrm{cm} \)}{\( \sqrt5\,\mathrm{cm} \)}{\( 6\sqrt3\,\mathrm{cm} \)}{}{}}
\LANGpl{
\Question{
Sześcian \( ABCDEFGH \) przedstawiony na rysunku ma krawędzie o długości \( a=6\,\mathrm{cm} \). Punkt \( S \) to środek krawędzi \( FG \). Oblicz odległość pomiędzy punktami \( E \) i \( S \).}
{\( 3\sqrt5\,\mathrm{cm} \)}{\( 6\sqrt5\,\mathrm{cm} \)}{\( \sqrt5\,\mathrm{cm} \)}{\( 6\sqrt3\,\mathrm{cm} \)}{}{}}
\LANGcs{
\Question{
Je dána krychle \( ABCDEFGH \) s hranou délky \( a=6\,\mathrm{cm} \). Určete vzdálenost bodů \( E \), \( S \), kde \( S \) je střed hrany \( FG \) (viz obrázek).}
{\( 3\sqrt5\,\mathrm{cm} \)}{\( 6\sqrt5\,\mathrm{cm} \)}{\( \sqrt5\,\mathrm{cm} \)}{\( 6\sqrt3\,\mathrm{cm} \)}{}{}}
\LANGsk{
\Question{
Je daná kocka \( ABCDEFGH \) s hranou dĺžky \( a=6\,\mathrm{cm} \). Určte vzdialenosť bodov \( E \), \( S \), kde \( S \) je stred hrany \( FG \) (viď obrázok).}
{\( 3\sqrt5\,\mathrm{cm} \)}{\( 6\sqrt5\,\mathrm{cm} \)}{\( \sqrt5\,\mathrm{cm} \)}{\( 6\sqrt3\,\mathrm{cm} \)}{}{}}
\LANGes{
\Question{
El cubo \( ABCDEFGH \) tiene la longitud de la arista \( a=6\,\mathrm{cm} \). Determina la distancia entre los puntos \( E \) y \( S \) donde \( S \) es el centro de la arista \( FG \).}
{\( 3\sqrt5\,\mathrm{cm} \)}{\( 6\sqrt5\,\mathrm{cm} \)}{\( \sqrt5\,\mathrm{cm} \)}{\( 6\sqrt3\,\mathrm{cm} \)}{}{}}
\End

\Data\ID{1103056004}
\Updated{2020-07-15 03:07:21}
\Level{A}
\SubArea{Lines and planes: distances and angles}
\IMG{1103056004.pdf}
\LANGen{
\Question{
The cube \( ABCDEFGH \) shown in the picture has edges of length \( a=6\,\mathrm{cm} \). Let \( S_1 \) be the midpoint of the diagonal \( ED \) and let \( S_2 \) be the midpoint of the diagonal \( CH \). Find the distance between the points \( S_1 \) and \( S_2 \).}
{\( 3\sqrt2\,\mathrm{cm} \)}{\( 6\sqrt2\,\mathrm{cm} \)}{\( \sqrt2\,\mathrm{cm} \)}{\( 6\sqrt3\,\mathrm{cm} \)}{}{}}
\LANGpl{
\Question{
Sześcian \( ABCDEFGH \) przedstawiony na rysunku ma krawędzie o długości \( a=6\,\mathrm{cm} \). Punkt \( S_1 \) to środek przekątnej \( ED \), punkt \( S_2 \) to środek przekątnej \( CH \). Oblicz odległość pomiędzy punktami \( S_1 \) i \( S_2 \).}
{\( 3\sqrt2\,\mathrm{cm} \)}{\( 6\sqrt2\,\mathrm{cm} \)}{\( \sqrt2\,\mathrm{cm} \)}{\( 6\sqrt3\,\mathrm{cm} \)}{}{}}
\LANGcs{
\Question{
Je dána krychle \( ABCDEFGH \) s hranou délky \( a=6\,\mathrm{cm} \). Určete vzdálenost bodů \( S_1 \) a \( S_2 \), kde \( S_1 \) je střed úhlopříčky \( ED \) a \( S_2 \) je střed úhlopříčky \( CH \) (viz obrázek).}
{\( 3\sqrt2\,\mathrm{cm} \)}{\( 6\sqrt2\,\mathrm{cm} \)}{\( \sqrt2\,\mathrm{cm} \)}{\( 6\sqrt3\,\mathrm{cm} \)}{}{}}
\LANGsk{
\Question{
Je daná kocka \( ABCDEFGH \) s hranou dĺžky \( a=6\,\mathrm{cm} \). Určte vzdialenosť bodov \( S_1 \) a \( S_2 \), kde \( S_1 \) je stred uhlopriečky \( ED \) a \( S_2 \) je stred uhlopriečky \( CH \) (viď obrázok).}
{\( 3\sqrt2\,\mathrm{cm} \)}{\( 6\sqrt2\,\mathrm{cm} \)}{\( \sqrt2\,\mathrm{cm} \)}{\( 6\sqrt3\,\mathrm{cm} \)}{}{}}
\LANGes{
\Question{
El cubo  \( ABCDEFGH \) tiene la longitud de la arista \( a=6\,\mathrm{cm} \). Determina la distancia entre los puntos \( S_1 \) y \( S_2 \), donde \( S_1 \) es el centro de la diagonal \( ED \) y \( S_2 \) es el centro de la diagonal \( CH \).}
{\( 3\sqrt2\,\mathrm{cm} \)}{\( 6\sqrt2\,\mathrm{cm} \)}{\( \sqrt2\,\mathrm{cm} \)}{\( 6\sqrt3\,\mathrm{cm} \)}{}{}}
\End

\Data\ID{1103056005}
\Updated{2020-07-15 03:07:21}
\Level{A}
\SubArea{Lines and planes: distances and angles}
\IMG{1103056005.pdf}
\LANGen{
\Question{
The cube \( ABCDEFGH \) shown in the picture has edges of length \( a \). Let \( S \) be the midpoint of the diagonal \( CF \). Find the distance between the point \( S \) and the line \( AH \).}
{\( a \)}{\( a\sqrt2 \)}{\( a\frac{\sqrt2}2 \)}{\( a\sqrt3 \)}{}{}}
\LANGpl{
\Question{
Sześcian \( ABCDEFGH \) przedstawiony na rysunku ma krawędzie o długości \( a \). Punkt \( S \) to środek przekątnej \( CF \). Oblicz odległość pomiędzy punktem \( S \) a prostą \( AH \).}
{\( a \)}{\( a\sqrt2 \)}{\( a\frac{\sqrt2}2 \)}{\( a\sqrt3 \)}{}{}}
\LANGcs{
\Question{
Je dána krychle \( ABCDEFGH \) s hranou délky \( a \). Určete vzdálenost bodu \( S \) (střed úhlopříčky \( CF \)) od přímky \( AH \) (viz obrázek).}
{\( a \)}{\( a\sqrt2 \)}{\( a\frac{\sqrt2}2 \)}{\( a\sqrt3 \)}{}{}}
\LANGsk{
\Question{
Je daná kocka \( ABCDEFGH \) s hranou dĺžky \( a \). Určte vzdialenosť bodu \( S \) (stred uhlopriečky \( CF \)) od priamky \( AH \) (viď obrázok).}
{\( a \)}{\( a\sqrt2 \)}{\( a\frac{\sqrt2}2 \)}{\( a\sqrt3 \)}{}{}}
\LANGes{
\Question{
El cubo \( ABCDEFGH \)  tiene la longitud de la arista \( a \). Determina la distancia entre la recta \( AH \) y el punto \( S \) donde \( S \) es el centro de la diagonal \( CF \).}
{\( a \)}{\( a\sqrt2 \)}{\( a\frac{\sqrt2}2 \)}{\( a\sqrt3 \)}{}{}}
\End

\Data\ID{1103056006}
\Updated{2020-07-15 03:07:21}
\Level{A}
\SubArea{Lines and planes: distances and angles}
\IMG{1103056006.pdf}
\LANGen{
\Question{
The cube \( ABCDEFGH \) shown in the picture has edges of length \( a=6\,\mathrm{cm} \). Find the distance between the point \( B \) and the line \( EG \).}
{\( 3\sqrt6\,\mathrm{cm} \)}{\( 6\sqrt3\,\mathrm{cm} \)}{\( 3\sqrt5\,\mathrm{cm} \)}{\( \frac{3\sqrt6}2\,\mathrm{cm} \)}{}{}}
\LANGpl{
\Question{
Sześcian \( ABCDEFGH \) przedstawiony na rysunku ma krawędzie o długości \( a=6\,\mathrm{cm} \). Oblicz odległość pomiędzy punktem \( B \) a prostą \( EG \).}
{\( 3\sqrt6\,\mathrm{cm} \)}{\( 6\sqrt3\,\mathrm{cm} \)}{\( 3\sqrt5\,\mathrm{cm} \)}{\( \frac{3\sqrt6}2\,\mathrm{cm} \)}{}{}}
\LANGcs{
\Question{
Je dána krychle \( ABCDEFGH \) s hranou délky \( a=6\,\mathrm{cm} \). Určete vzdálenost bodu \( B \) od přímky \( EG \) (viz obrázek).}
{\( 3\sqrt6\,\mathrm{cm} \)}{\( 6\sqrt3\,\mathrm{cm} \)}{\( 3\sqrt5\,\mathrm{cm} \)}{\( \frac{3\sqrt6}2\,\mathrm{cm} \)}{}{}}
\LANGsk{
\Question{
Je daná kocka \( ABCDEFGH \) s dĺžkou hrany \( a=6\,\mathrm{cm} \). Určte vzdialenosť bodu \( B \) a priamky \( EG \).}
{\( 3\sqrt6\,\mathrm{cm} \)}{\( 6\sqrt3\,\mathrm{cm} \)}{\( 3\sqrt5\,\mathrm{cm} \)}{\( \frac{3\sqrt6}2\,\mathrm{cm} \)}{}{}}
\LANGes{
\Question{
El cubo  \( ABCDEFGH \) tiene la longitud de la arista \( a=6\,\mathrm{cm} \). Determina la distancia entre el punto \( B \) y la recta \( EG \).}
{\( 3\sqrt6\,\mathrm{cm} \)}{\( 6\sqrt3\,\mathrm{cm} \)}{\( 3\sqrt5\,\mathrm{cm} \)}{\( \frac{3\sqrt6}2\,\mathrm{cm} \)}{}{}}
\End

\Data\ID{1103056007}
\Updated{2020-07-15 03:07:21}
\Level{A}
\SubArea{Lines and planes: distances and angles}
\IMG{1103056007.pdf}
\LANGen{
\Question{
The cube \( ABCDEFGH \) shown in the picture has edges of length \( a \). Find the distance between the lines \( AB \) and \( HG \).}
{\( a\sqrt2 \)}{\( a\sqrt3 \)}{\( \frac{a\sqrt5}2 \)}{\( a \)}{}{}}
\LANGpl{
\Question{
Sześcian \( ABCDEFGH \) przedstawiony na rysunku ma krawędzie o długości \( a \). Oblicz odległość pomiędzy prostymi \( AB \) i \( HG \).}
{\( a\sqrt2 \)}{\( a\sqrt3 \)}{\( \frac{a\sqrt5}2 \)}{\( a \)}{}{}}
\LANGcs{
\Question{
Je dána krychle \( ABCDEFGH \) s hranou délky \( a \). Určete vzdálenost přímek \( AB \) a \( HG \) (viz obrázek).}
{\( a\sqrt2 \)}{\( a\sqrt3 \)}{\( \frac{a\sqrt5}2 \)}{\( a \)}{}{}}
\LANGsk{
\Question{
Je daná kocka \( ABCDEFGH \) s hranou dĺžky \( a \). Určte vzdialenosť priamok \( AB \) a \( HG \) (viď obrázok).}
{\( a\sqrt2 \)}{\( a\sqrt3 \)}{\( \frac{a\sqrt5}2 \)}{\( a \)}{}{}}
\LANGes{
\Question{
El cubo \( ABCDEFGH \) tiene la longitud de la arista \( a \). Determina la distancia entre las rectas\( AB \) y \( HG \).}
{\( a\sqrt2 \)}{\( a\sqrt3 \)}{\( \frac{a\sqrt5}2 \)}{\( a \)}{}{}}
\End

\Data\ID{1103061401}
\Updated{2020-07-15 03:07:21}
\Level{A}
\SubArea{Lines and planes: distances and angles}
\IMG{1103061401.pdf}
\LANGen{
\Question{
The rectangular box \( ABCDEFGH \) has sides of lengths \( |AB|=6\,\mathrm{cm} \), \( |BC|=4\,\mathrm{cm} \), \( |AE|=8\,\mathrm{cm} \). Find the angle between the plane \( EFB \) and the plane \( HGB \) (see the picture). Round the result to two decimal places.}
{\( 26.57^{\circ} \)}{\( 45^{\circ} \)}{\( 22.5^{\circ} \)}{\( 31.43^{\circ} \)}{}{}}
\LANGpl{
\Question{
Dany jest prostopadłościan \( ABCDEFGH \) o bokach \( |AB|=6\,\mathrm{cm} \), \( |BC|=4\,\mathrm{cm} \), \( |AE|=8\,\mathrm{cm} \). Oblicz kąt pomiędzy płaszczyzną \( EFB \) a  płaszczyzną \( HGB \) (spójrz na rysunek). Zaokrągli wynik do dwóch miejsc po przecinku.}
{\( 26{,}57^{\circ} \)}{\( 45^{\circ} \)}{\( 22{,}5^{\circ} \)}{\( 31{,}43^{\circ} \)}{}{}}
\LANGcs{
\Question{
V kvádru \( ABCDEFGH \) je dáno: \( |AB|=6\,\mathrm{cm} \), \( |BC|=4\,\mathrm{cm} \), \( |AE|=8\,\mathrm{cm} \). Určete odchylku rovin \( EFB \) a \( HGB \) (viz obrázek). Výsledek zaokrouhlete na dvě desetinná místa.}
{\( 26{,}57^{\circ} \)}{\( 45^{\circ} \)}{\( 22{,}5^{\circ} \)}{\( 31{,}43^{\circ} \)}{}{}}
\LANGsk{
\Question{
V kvádri \( ABCDEFGH \) je dané: \( |AB|=6\,\mathrm{cm} \), \( |BC|=4\,\mathrm{cm} \), \( |AE|=8\,\mathrm{cm} \). Určte odchýlku rovín \( EFB \) a \( HGB \) (viď obrázok). Výsledok zaokrúhlite na dve desatinné miesta.}
{\( 26{,}57^{\circ} \)}{\( 45^{\circ} \)}{\( 22{,}5^{\circ} \)}{\( 31{,}43^{\circ} \)}{}{}}
\LANGes{
\Question{
El ortoedro \( ABCDEFGH \) tiene las longitudes de las aristas \( |AB|=6\,\mathrm{cm} \), \( |BC|=4\,\mathrm{cm} \), \( |AE|=8\,\mathrm{cm} \). Determina el ángulo entre el plano \( EFB \) y el plano \( HGB \) (mira la imagen). Redondea a dos posiciones decimales.}
{\( 26.57^{\circ} \)}{\( 45^{\circ} \)}{\( 22.5^{\circ} \)}{\( 31.43^{\circ} \)}{}{}}
\End

\Data\ID{1103061402}
\Updated{2020-07-15 03:07:21}
\Level{A}
\SubArea{Lines and planes: distances and angles}
\IMG{1103061402.pdf}
\LANGen{
\Question{
The rectangular box \( ABCDEFGH \) has  sides of lengths \( |AB|=6\,\mathrm{cm} \), \( |BC|=4\,\mathrm{cm} \), \( |AE|=8\,\mathrm{cm} \). Find the angle between the lines \( AC \) and \( SC \), where \( S \) is the midpoint of the diagonal \( EG \) (see the picture). Round the result to two decimal places.}
{\( 65.74^{\circ} \)}{\( 70.28^{\circ} \)}{\( 75^{\circ} \)}{\( 45^{\circ} \)}{}{}}
\LANGpl{
\Question{
Dany jest prostopadłośćian \( ABCDEFGH \) o bokach \( |AB|=6\,\mathrm{cm} \), \( |BC|=4\,\mathrm{cm} \), \( |AE|=8\,\mathrm{cm} \). Oblicz kąt pomiędzy prostymi \( AC \) i \( SC \), gdzie \( S \) to środek przekątnej \( EG \) (spójrz na rysunek). Zaokrągli wynik do dwóch miejsc po przecinku.}
{\( 65{,}74^{\circ} \)}{\( 70{,}28^{\circ} \)}{\( 75^{\circ} \)}{\( 45^{\circ} \)}{}{}}
\LANGcs{
\Question{
V kvádru \( ABCDEFGH \) je dáno: \( |AB|=6\,\mathrm{cm} \), \( |BC|=4\,\mathrm{cm} \), \( |AE|=8\,\mathrm{cm} \). Určete odchylku přímek \( AC \) a \( SC \), kde \( S \) je střed úhlopříčky \( EG \) (viz obrázek). Výsledek zaokrouhlete na dvě desetinná místa.}
{\( 65{,}74^{\circ} \)}{\( 70{,}28^{\circ} \)}{\( 75^{\circ} \)}{\( 45^{\circ} \)}{}{}}
\LANGsk{
\Question{
V kvádri \( ABCDEFGH \) je dané: \( |AB|=6\,\mathrm{cm} \), \( |BC|=4\,\mathrm{cm} \), \( |AE|=8\,\mathrm{cm} \). Určte odchýlku priamok \( AC \) a \( SC \), kde \( S \) je stred uhlopriečky \( EG \) (viď obrázok). Výsledok zaokrúhlite na dve desatinné miesta.}
{\( 65{,}74^{\circ} \)}{\( 70{,}28^{\circ} \)}{\( 75^{\circ} \)}{\( 45^{\circ} \)}{}{}}
\LANGes{
\Question{
El ortoedro\( ABCDEFGH \) tiene las longitudes de las aristas \( |AB|=6\,\mathrm{cm} \), \( |BC|=4\,\mathrm{cm} \), \( |AE|=8\,\mathrm{cm} \). Determina el ángulo entre las rectas \( AC \) y \( SC \), donde \( S \) es el centro de la diagonal \( EG \) (mira la imagen). Redondea a dos posiciones decimales.}
{\( 65.74^{\circ} \)}{\( 70.28^{\circ} \)}{\( 75^{\circ} \)}{\( 45^{\circ} \)}{}{}}
\End

\Data\ID{1103061403}
\Updated{2020-07-15 03:07:21}
\Level{A}
\SubArea{Lines and planes: distances and angles}
\IMG{1103061403.pdf}
\LANGen{
\Question{
The rectangular box \( ABCDEFGH \) has  sides of lengths \( |AB|=6\,\mathrm{cm} \), \( |BC|=4\,\mathrm{cm} \), \( |AE|=8\,\mathrm{cm} \). Find the angle between the lines \( CG \) and \( SC \), where \( S \) is the midpoint of the diagonal \( EG \) (see the picture). Round the result to two decimal places.}
{\( 24.26^{\circ} \)}{\( 45^{\circ} \)}{\( 35.24^{\circ} \)}{\( 63.21^{\circ} \)}{}{}}
\LANGpl{
\Question{
Dany jest prostopadłościan  \( ABCDEFGH \) o bokach  \( |AB|=6\,\mathrm{cm} \), \( |BC|=4\,\mathrm{cm} \), \( |AE|=8\,\mathrm{cm} \). Oblicz kąt pomiędzy prostymi  \( CG \) i \( SC \), gdzie \( S \) to środek przekątnej  \( EG \) (spójrz na rysunek). Zaokrągli wynik do dwóch miejsc po przecinku.}
{\( 24{,}26^{\circ} \)}{\( 45^{\circ} \)}{\( 35{,}24^{\circ} \)}{\( 63{,}21^{\circ} \)}{}{}}
\LANGcs{
\Question{
V kvádru \( ABCDEFGH \) je dáno: \( |AB|=6\,\mathrm{cm} \), \( |BC|=4\,\mathrm{cm} \), \( |AE|=8\,\mathrm{cm} \). Určete odchylku přímek \( CG \) a \( SC \), kde \( S \) je střed úhlopříčky \( EG \) (viz obrázek). Výsledek zaokrouhlete na dvě desetinná místa.}
{\( 24{,}26^{\circ} \)}{\( 45^{\circ} \)}{\( 35{,}24^{\circ} \)}{\( 63{,}21^{\circ} \)}{}{}}
\LANGsk{
\Question{
V kvádri \( ABCDEFGH \) je dané: \( |AB|=6\,\mathrm{cm} \), \( |BC|=4\,\mathrm{cm} \), \( |AE|=8\,\mathrm{cm} \). Určte odchýlku priamok \( CG \) a \( SC \), kde \( S \) je stred uhlopriečky \( EG \) (viď obrázok). Výsledok zaokrúhlite na dve desatinné miesta.}
{\( 24{,}26^{\circ} \)}{\( 45^{\circ} \)}{\( 35{,}24^{\circ} \)}{\( 63{,}21^{\circ} \)}{}{}}
\LANGes{
\Question{
El ortoedro \( ABCDEFGH \) tiene las longitudes de las aristas \( |AB|=6\,\mathrm{cm} \), \( |BC|=4\,\mathrm{cm} \), \( |AE|=8\,\mathrm{cm} \). Determina el ángulo entre las rectas \( CG \) y \( SC \), donde \( S \) es el centro de la diagonal \( EG \) (mira la imagen). Redondea a dos posiciones decimales.}
{\( 24.26^{\circ} \)}{\( 45^{\circ} \)}{\( 35.24^{\circ} \)}{\( 63.21^{\circ} \)}{}{}}
\End

\Data\ID{1103061404}
\Updated{2020-07-15 03:07:21}
\Level{A}
\SubArea{Lines and planes: distances and angles}
\IMG{1103061404.pdf}
\LANGen{
\Question{
The rectangular box \( ABCDEFGH \) has sides of lengths \( |AB|=6\,\mathrm{cm} \), \( |BC|=4\,\mathrm{cm} \), \( |AE|=8\,\mathrm{cm} \). Find the angle between the lines \( HG \) and \( AH \) (see the picture). Round the result to two decimal places.}
{\( 90^{\circ} \)}{\( 85^{\circ} \)}{\( 56.15^{\circ} \)}{\( 33.85^{\circ} \)}{}{}}
\LANGpl{
\Question{
Dany jest prostopadłościan  \( ABCDEFGH \) o bokach \( |AB|=6\,\mathrm{cm} \), \( |BC|=4\,\mathrm{cm} \), \( |AE|=8\,\mathrm{cm} \). Oblicz kąt pomiędzy prostymi  \( HG \) i  \( AH \) (spójrz na rysunek). Zaokrągli wynik do dwóch miejsc po przecinku.}
{\( 90^{\circ} \)}{\( 85^{\circ} \)}{\( 56{,}15^{\circ} \)}{\( 33{,}85^{\circ} \)}{}{}}
\LANGcs{
\Question{
V kvádru  \( ABCDEFGH \) je dáno: \( |AB|=6\,\mathrm{cm} \), \( |BC|=4\,\mathrm{cm} \), \( |AE|=8\,\mathrm{cm} \). Určete odchylku přímek \( HG \) a \( AH \) (viz obrázek). Výsledek zaokrouhlete na dvě desetinná místa.}
{\( 90^{\circ} \)}{\( 85^{\circ} \)}{\( 56{,}15^{\circ} \)}{\( 33{,}85^{\circ} \)}{}{}}
\LANGsk{
\Question{
V kvádri  \( ABCDEFGH \) je dané: \( |AB|=6\,\mathrm{cm} \), \( |BC|=4\,\mathrm{cm} \), \( |AE|=8\,\mathrm{cm} \). Určte odchýlku priamok \( HG \) a \( AH \) (viď obrázok). Výsledok zaokrúhlite na dve desatinné miesta.}
{\( 90^{\circ} \)}{\( 85^{\circ} \)}{\( 56{,}15^{\circ} \)}{\( 33{,}85^{\circ} \)}{}{}}
\LANGes{
\Question{
El ortoedro \( ABCDEFGH \) tiene las longitudes de las aristas \( |AB|=6\,\mathrm{cm} \), \( |BC|=4\,\mathrm{cm} \), \( |AE|=8\,\mathrm{cm} \). Determina el ángulo entre las rectas \( HG \) y \( AH \) (mira la imagen). Redondea a dos posiciones decimales.}
{\( 90^{\circ} \)}{\( 85^{\circ} \)}{\( 56.15^{\circ} \)}{\( 33.85^{\circ} \)}{}{}}
\End

\Data\ID{1103061405}
\Updated{2020-07-15 03:07:21}
\Level{A}
\SubArea{Lines and planes: distances and angles}
\IMG{1103061405.pdf}
\LANGen{
\Question{
The rectangular box \( ABCDEFGH \) has sides of lengths \( |AB|=6\,\mathrm{cm} \), \( |BC|=4\,\mathrm{cm} \), \( |AE|=8\,\mathrm{cm} \). Find the angle between the lines \( AG \) and \( BH \) (see the picture). Round the result to two decimal places.}
{\( 67.71^{\circ} \)}{\( 45^{\circ} \)}{\( 37.08^{\circ} \)}{\( 49.1^{\circ} \)}{}{}}
\LANGpl{
\Question{
Dany jest prostopadłościan \( ABCDEFGH \) o bokach \( |AB|=6\,\mathrm{cm} \), \( |BC|=4\,\mathrm{cm} \), \( |AE|=8\,\mathrm{cm} \). Oblicz kąt pomiędzy prostymi  \( AG \) i \( BH \) (spójrz na rysunek). Zaokrągli wynik do dwóch miejsc po przecinku.}
{\( 67{,}71^{\circ} \)}{\( 45^{\circ} \)}{\( 37{,}08^{\circ} \)}{\( 49{,}1^{\circ} \)}{}{}}
\LANGcs{
\Question{
V kvádru \( ABCDEFGH \) je dáno: \( |AB|=6\,\mathrm{cm} \), \( |BC|=4\,\mathrm{cm} \), \( |AE|=8\,\mathrm{cm} \). Určete odchylku přímek \( AG \) a \( BH \) (viz obrázek). Výsledek zaokrouhlete na dvě desetinná místa.}
{\( 67{,}71^{\circ} \)}{\( 45^{\circ} \)}{\( 37{,}08^{\circ} \)}{\( 49{,}1^{\circ} \)}{}{}}
\LANGsk{
\Question{
V kvádri \( ABCDEFGH \) je dané: \( |AB|=6\,\mathrm{cm} \), \( |BC|=4\,\mathrm{cm} \), \( |AE|=8\,\mathrm{cm} \). Určte odchýlku priamok \( AG \) a \( BH \) (viď obrázok). Výsledok zaokrúhlite na dve desatinné miesta.}
{\( 67{,}71^{\circ} \)}{\( 45^{\circ} \)}{\( 37{,}08^{\circ} \)}{\( 49{,}1^{\circ} \)}{}{}}
\LANGes{
\Question{
El ortoedro \( ABCDEFGH \) tiene las longitudes de las aristas \( |AB|=6\,\mathrm{cm} \), \( |BC|=4\,\mathrm{cm} \), \( |AE|=8\,\mathrm{cm} \). Determina el ángulo entre las rectas \( AG \) y \( BH \) (mira la imagen). Redondea a dos posiciones decimales.}
{\( 67.71^{\circ} \)}{\( 45^{\circ} \)}{\( 37.08^{\circ} \)}{\( 49.1^{\circ} \)}{}{}}
\End

\Data\ID{1103061406}
\Updated{2020-07-15 03:07:21}
\Level{A}
\SubArea{Lines and planes: distances and angles}
\IMG{1103061406.pdf}
\LANGen{
\Question{
The rectangular box \( ABCDEFGH \) has sides of lengths \( |AB|=6\,\mathrm{cm} \), \( |BC|=4\,\mathrm{cm} \), \( |AE|=8\,\mathrm{cm} \). Find the angle between the planes \( ABC \) and \( EFC \) (see the picture). Round the result to two decimal places.}
{\( 63.43^{\circ} \)}{\( 26.57^{\circ} \)}{\( 45^{\circ} \)}{\( 56.2^{\circ} \)}{}{}}
\LANGpl{
\Question{
Dany jest prostopadłościan  \( ABCDEFGH \) o bokach \( |AB|=6\,\mathrm{cm} \), \( |BC|=4\,\mathrm{cm} \), \( |AE|=8\,\mathrm{cm} \). Oblicz kąt pomiędzy płaszczyzną  \( ABC \) a \( EFC \) (spójrz na rysunek). Zaokrągli wynik do dwóch miejsc po przecinku.}
{\( 63{,}43^{\circ} \)}{\( 26{,}57^{\circ} \)}{\( 45^{\circ} \)}{\( 56{,}2^{\circ} \)}{}{}}
\LANGcs{
\Question{
V kvádru \( ABCDEFGH \) je dáno: \( |AB|=6\,\mathrm{cm} \), \( |BC|=4\,\mathrm{cm} \), \( |AE|=8\,\mathrm{cm} \). Určete odchylku rovin \( ABC \) a \( EFC \) (viz obrázek). Výsledek zaokrouhlete na dvě desetinná místa.}
{\( 63{,}43^{\circ} \)}{\( 26{,}57^{\circ} \)}{\( 45^{\circ} \)}{\( 56{,}2^{\circ} \)}{}{}}
\LANGsk{
\Question{
V kvádri \( ABCDEFGH \) je dané: \( |AB|=6\,\mathrm{cm} \), \( |BC|=4\,\mathrm{cm} \), \( |AE|=8\,\mathrm{cm} \). Určte odchýlku rovín \( ABC \) a \( EFC \) (viď obrázok). Výsledok zaokrúhlite na dve desatinné miesta.}
{\( 63{,}43^{\circ} \)}{\( 26{,}57^{\circ} \)}{\( 45^{\circ} \)}{\( 56{,}2^{\circ} \)}{}{}}
\LANGes{
\Question{
El ortoedro \( ABCDEFGH \) tiene las longitudes de las aristas \( |AB|=6\,\mathrm{cm} \), \( |BC|=4\,\mathrm{cm} \), \( |AE|=8\,\mathrm{cm} \). Determina el ángulo entre los planos \( ABC \) y \( EFC \) (mira la imagen). Redondea a dos posiciones decimales.}
{\( 63.43^{\circ} \)}{\( 26.57^{\circ} \)}{\( 45^{\circ} \)}{\( 56.2^{\circ} \)}{}{}}
\End

\Data\ID{1103061407}
\Updated{2020-07-15 03:07:21}
\Level{A}
\SubArea{Lines and planes: distances and angles}
\IMG{1103061407.pdf}
\LANGen{
\Question{
The rectangular box \( ABCDEFGH \) has sides of lengths \( |AB|=6\,\mathrm{cm} \), \( |BC|=4\,\mathrm{cm} \), \( |AE|=8\,\mathrm{cm} \). Find the angle between the lines \( CG \) and \( EC \) (see the picture). Round the result to two decimal places.}
{\( 42.03^{\circ} \)}{\( 47.97^{\circ} \)}{\( 45^{\circ} \)}{\( 50.04^{\circ} \)}{}{}}
\LANGpl{
\Question{
Dany jest prostopadłościan  \( ABCDEFGH \) o bokach  \( |AB|=6\,\mathrm{cm} \), \( |BC|=4\,\mathrm{cm} \), \( |AE|=8\,\mathrm{cm} \). Oblicz kąt pomiędzy prostymi  \( CG \) i  \( EC \) (spójrz na rysunek). Zaokrągli wynik do dwóch miejsc po przecinku.}
{\( 42{,}03^{\circ} \)}{\( 47{,}97^{\circ} \)}{\( 45^{\circ} \)}{\( 50{,}04^{\circ} \)}{}{}}
\LANGcs{
\Question{
V kvádru \( ABCDEFGH \) je dáno: \( |AB|=6\,\mathrm{cm} \), \( |BC|=4\,\mathrm{cm} \), \( |AE|=8\,\mathrm{cm} \). Určete odchylku přímek \( CG \) a \( EC \) (viz obrázek). Výsledek zaokrouhlete na dvě desetinná místa.}
{\( 42{,}03^{\circ} \)}{\( 47{,}97^{\circ} \)}{\( 45^{\circ} \)}{\( 50{,}04^{\circ} \)}{}{}}
\LANGsk{
\Question{
V kvádri \( ABCDEFGH \) je dané: \( |AB|=6\,\mathrm{cm} \), \( |BC|=4\,\mathrm{cm} \), \( |AE|=8\,\mathrm{cm} \). Určte odchýlku priamok \( CG \) a \( EC \) (viď obrázok). Výsledok zaokrúhlite na dve desatinné miesta.}
{\( 42{,}03^{\circ} \)}{\( 47{,}97^{\circ} \)}{\( 45^{\circ} \)}{\( 50{,}04^{\circ} \)}{}{}}
\LANGes{
\Question{
El ortoedro \( ABCDEFGH \) tiene las longitudes de las aristas \( |AB|=6\,\mathrm{cm} \), \( |BC|=4\,\mathrm{cm} \), \( |AE|=8\,\mathrm{cm} \). Determina el ángulo entre las rectas \( CG \) y \( EC \) (mira la imagen). Redondea a dos posiciones decimales.}
{\( 42.03^{\circ} \)}{\( 47.97^{\circ} \)}{\( 45^{\circ} \)}{\( 50.04^{\circ} \)}{}{}}
\End

\Data\ID{1103061408}
\Updated{2020-07-15 03:07:21}
\Level{A}
\SubArea{Lines and planes: distances and angles}
\IMG{1103061408_0.pdf}
\LANGen{
\Question{
The rectangular box \( ABCDEFGH \) has sides of lengths \( |AB|=|BC|=6\,\mathrm{cm} \), \( |AE|=8\,\mathrm{cm} \). Find the angle between the planes \( ABC \) and \( AFH \) (see the picture). Round the result to two decimal places.}
{\( 62.06^{\circ} \)}{\( 53.13^{\circ} \)}{\( 45^{\circ} \)}{\( 60^{\circ} \)}{}{}}
\LANGpl{
\Question{
Dany jest prostopadłościan  \( ABCDEFGH \) o bokach  \( |AB|=|BC|=6\,\mathrm{cm} \), \( |AE|=8\,\mathrm{cm} \). Oblicz kąt pomiędzy płaszczyzną  \( ABC \) a \( AFH \) (spójrz na rysunek). Zaokrągli wynik do dwóch miejsc po przecinku.}
{\( 62{,}06^{\circ} \)}{\( 53{,}13^{\circ} \)}{\( 45^{\circ} \)}{\( 60^{\circ} \)}{}{}}
\LANGcs{
\Question{
V kvádru \( ABCDEFGH \) je dáno \( |AB|=|BC|=6\,\mathrm{cm} \), \( |AE|=8\,\mathrm{cm} \). Určete odchylku rovin \( ABC \) a \( AFH \) (viz obrázek). Výsledek zaokrouhlete na dvě desetinná místa.}
{\( 62{,}06^{\circ} \)}{\( 53{,}13^{\circ} \)}{\( 45^{\circ} \)}{\( 60^{\circ} \)}{}{}}
\LANGsk{
\Question{
V kvádri \( ABCDEFGH \) je dané \( |AB|=|BC|=6\,\mathrm{cm} \), \( |AE|=8\,\mathrm{cm} \). Určte odchýlku rovín \( ABC \) a \( AFH \) (viď obrázok). Výsledok zaokrúhlite na dve desatinné miesta.}
{\( 62{,}06^{\circ} \)}{\( 53{,}13^{\circ} \)}{\( 45^{\circ} \)}{\( 60^{\circ} \)}{}{}}
\LANGes{
\Question{
El ortoedro \( ABCDEFGH \) tiene las longitudes de las aristas \( |AB|=|BC|=6\,\mathrm{cm} \), \( |AE|=8\,\mathrm{cm} \). Determina el ángulo entre los planos \( ABC \) y \( AFH \) (mira la imagen). Redondea a dos posiciones decimales.}
{\( 62.06^{\circ} \)}{\( 53.13^{\circ} \)}{\( 45^{\circ} \)}{\( 60^{\circ} \)}{}{}}
\End

\Data\ID{9000045709}
\Updated{2020-07-15 03:07:34}
\Level{A}
\SubArea{Lines and planes: distances and angles}
\IMG{000457_09-figure0.pdf}
\LANGen{
\Question{
Let \(\omega \) be the
angle between the solid diagonal of a box and the base of this box. Find the expression which
allows to find \(\omega \).}
{\(\mathop{\mathrm{tg}}\nolimits \omega = \frac{\sqrt{2}} 
 {2} \)}{\(\cos \omega = \frac{\sqrt{2}} 
 {2} \)}{\(\sin \omega = \frac{\sqrt{2}} 
 {2} \)}{\(\mathop{\mathrm{cotg}}\nolimits \omega = \frac{\sqrt{2}} 
 {2} \)}{}{}}
\LANGpl{
\Question{
Kąt  \(\omega \) to kąt między przekątną sześcianu, a przekątną podstawy tego sześcianu. Wybierz wyrażenie, które pozwala obliczyć 
 \(\omega \).}
{\(\mathop{\mathrm{tg}}\nolimits \omega = \frac{\sqrt{2}} 
 {2} \)}{\(\cos \omega = \frac{\sqrt{2}} 
 {2} \)}{\(\sin \omega = \frac{\sqrt{2}} 
 {2} \)}{\(\mathop{\mathrm{cotg}}\nolimits \omega = \frac{\sqrt{2}} 
 {2} \)}{}{}}
\LANGcs{
\Question{
Je dána krychle s hranou délky \(a\).
Vyberte vztah, který platí pro odchylku
\(\omega \) 
tělesové úhlopříčky od roviny podstavy.}
{\(\mathop{\mathrm{tg}}\nolimits \omega = \frac{\sqrt{2}} 
 {2} \)}{\(\cos \omega = \frac{\sqrt{2}} 
 {2} \)}{\(\sin \omega = \frac{\sqrt{2}} 
 {2} \)}{\(\mathop{\mathrm{cotg}}\nolimits \omega = \frac{\sqrt{2}} 
 {2} \)}{}{}}
\LANGsk{
\Question{
Je daná kocka s hranou dĺžky \(a\).
Vyberte vzťah, ktorý platí pre odchýlku
\(\omega \) 
telesovej uhlopriečky od roviny podstavy.}
{\(\mathop{\mathrm{tg}}\nolimits \omega = \frac{\sqrt{2}} 
 {2} \)}{\(\cos \omega = \frac{\sqrt{2}} 
 {2} \)}{\(\sin \omega = \frac{\sqrt{2}} 
 {2} \)}{\(\mathop{\mathrm{cotg}}\nolimits \omega = \frac{\sqrt{2}} 
 {2} \)}{}{}}
\LANGes{
\Question{
Sea \(\omega \)  el ángulo entre la diagonal espacial y la base de un cubo. Determina la expresión correcta para \(\omega \).}
{\(\mathop{\mathrm{tg}}\nolimits \omega = \frac{\sqrt{2}} 
 {2} \)}{\(\cos \omega = \frac{\sqrt{2}} 
 {2} \)}{\(\sin \omega = \frac{\sqrt{2}} 
 {2} \)}{\(\mathop{\mathrm{cotg}}\nolimits \omega = \frac{\sqrt{2}} 
 {2} \)}{}{}}
\End

\Data\ID{9000046407}
\Updated{2021-06-29 03:06:07}
\Level{A}
\SubArea{Lines and planes: distances and angles}
\IMG{000464_07-figure0.pdf}
\LANGen{
\Question{
Find the angle between the side of a cube and the cube’s space diagonal. Round the result to two decimal places.}
{\(35.26^{\circ }\)}{\(45^{\circ }\)}{\(54.76^{\circ }\)}{}{}{}}
\LANGpl{
\Question{
Wyznacz kąt między przekątną sześcianu, a podstawą  sześcianu. Zaokrągli wynik do dwóch miejsc po przecinku.}
{\(35.26^{\circ }\)}{\(45^{\circ }\)}{\(54.76^{\circ }\)}{}{}{}}
\LANGcs{
\Question{
Určete odchylku tělesové úhlopříčky krychle od stěny krychle (výsledek je
zaokrouhlen na \(2\) 
desetinná místa).}
{\(35{,}26^{\circ }\)}{\(45^{\circ }\)}{\(54{,}76^{\circ }\)}{}{}{}}
\LANGsk{
\Question{
Určte odchýlku telesovej uhlopriečky kocky od steny kocky (výsledok je
zaokrúhlený na \(2\) 
desatinné miesta).}
{\(35{,}26^{\circ }\)}{\(45^{\circ }\)}{\(54{,}76^{\circ }\)}{}{}{}}
\LANGes{
\Question{
Determina el ángulo entre la diagonal espacial y la cara de un cubo. Redondea el resultado a dos posiciones decimales.}
{\(35.26^{\circ }\)}{\(45^{\circ }\)}{\(54.76^{\circ }\)}{}{}{}}
\End

\Data\ID{9000120302}
\Updated{2020-07-15 03:07:37}
\Level{A}
\SubArea{Lines and planes: distances and angles}
\IMG{001203_02-figure0.pdf}
\LANGen{
\Question{
The rectangular box has the sides \(a = 5\, \mathrm{cm}\),
\(b = 8\, \mathrm{cm}\) and
\(c = \sqrt{111}\, \mathrm{cm}\). Find the length of
the solid diagonal \(u\).}
{\(10\sqrt{2}\, \mathrm{cm}\)}{\(\sqrt{222}\, \mathrm{cm}\)}{\(20\, \mathrm{cm}\)}{\(2\sqrt{10}\, \mathrm{cm}\)}{\(5\sqrt{7}\, \mathrm{cm}\)}{}}
\LANGpl{
\Question{
Dany jest prostopadłościan  o bokach  \(a = 5\, \mathrm{cm}\),
\(b = 8\, \mathrm{cm}\) i
\(c = \sqrt{111}\, \mathrm{cm}\).  Oblicz długość przekątnej
 \(u\).}
{\(10\sqrt{2}\, \mathrm{cm}\)}{\(\sqrt{222}\, \mathrm{cm}\)}{\(20\, \mathrm{cm}\)}{\(2\sqrt{10}\, \mathrm{cm}\)}{\(5\sqrt{7}\, \mathrm{cm}\)}{}}
\LANGcs{
\Question{
Délky hran čtyřbokého hranolu jsou
\(a = 5\, \mathrm{cm}\),
\(b = 8\, \mathrm{cm}\),
\(c = \sqrt{111}\, \mathrm{cm}\).
Délka tělesové úhlopříčky je:}
{\(10\sqrt{2}\, \mathrm{cm}\)}{\(\sqrt{222}\, \mathrm{cm}\)}{\(20\, \mathrm{cm}\)}{\(2\sqrt{10}\, \mathrm{cm}\)}{\(5\sqrt{7}\, \mathrm{cm}\)}{}}
\LANGsk{
\Question{
Dĺžky hrán štvorbokého hranolu sú
\(a = 5\, \mathrm{cm}\),
\(b = 8\, \mathrm{cm}\),
\(c = \sqrt{111}\, \mathrm{cm}\).
Dĺžka telesovej uhlopriečky je:}
{\(10\sqrt{2}\, \mathrm{cm}\)}{\(\sqrt{222}\, \mathrm{cm}\)}{\(20\, \mathrm{cm}\)}{\(2\sqrt{10}\, \mathrm{cm}\)}{\(5\sqrt{7}\, \mathrm{cm}\)}{}}
\LANGes{
\Question{
El ortoedro tiene las aristas \(a = 5\, \mathrm{cm}\),
\(b = 8\, \mathrm{cm}\) y
\(c = \sqrt{111}\, \mathrm{cm}\). Determina la longitud de la diagonal espacial \(u\).}
{\(10\sqrt{2}\, \mathrm{cm}\)}{\(\sqrt{222}\, \mathrm{cm}\)}{\(20\, \mathrm{cm}\)}{\(2\sqrt{10}\, \mathrm{cm}\)}{\(5\sqrt{7}\, \mathrm{cm}\)}{}}
\End

\Data\ID{9000120303}
\Updated{2020-07-15 03:07:37}
\Level{A}
\SubArea{Lines and planes: distances and angles}
\IMG{001203_03-figure0.pdf}
\LANGen{
\Question{
Identify a valid relation involving the angle
\(\alpha \) 
defined as an angle between a solid diagonal and a face diagonal through the same
vertex in a cube.}
{\(\mathop{\mathrm{tg}}\nolimits \alpha = \frac{\sqrt{2}} 
 {2} \)}{\(\sin \alpha = \frac{\sqrt{3}} 
 {2} \)}{\(\cos \alpha = \frac{\sqrt{5}} 
 {3} \)}{\(\mathop{\mathrm{cotg}}\nolimits \alpha = \sqrt{3}\)}{\(\alpha = 45^{\circ }\)}{}}
\LANGpl{
\Question{
Wskaż właściwą zależność kąta 
\(\alpha \) 
między przekątną sześcianu, a przekątną ściany bocznej, obie przekątne przechodzą przez ten sam wierzchołek.}
{\(\mathop{\mathrm{tg}}\nolimits \alpha = \frac{\sqrt{2}} 
 {2} \)}{\(\sin \alpha = \frac{\sqrt{3}} 
 {2} \)}{\(\cos \alpha = \frac{\sqrt{5}} 
 {3} \)}{\(\mathop{\mathrm{cotg}}\nolimits \alpha = \sqrt{3}\)}{\(\alpha = 45^{\circ }\)}{}}
\LANGcs{
\Question{
Odchylka tělesové a stěnové úhlopříčky v krychli o hraně
\(a\) je
\(\alpha \).
Potom platí:}
{\(\mathop{\mathrm{tg}}\nolimits \alpha = \frac{\sqrt{2}} 
 {2} \)}{\(\sin \alpha = \frac{\sqrt{3}} 
 {2} \)}{\(\cos \alpha = \frac{\sqrt{5}} 
 {3} \)}{\(\mathop{\mathrm{cotg}}\nolimits \alpha = \sqrt{3}\)}{\(\alpha = 45^{\circ }\)}{}}
\LANGsk{
\Question{
Odchýlka telesovej a stenovej uhlopriečky v kocke s hranou
\(a\) je
\(\alpha \).
Potom platí:}
{\(\mathop{\mathrm{tg}}\nolimits \alpha = \frac{\sqrt{2}} 
 {2} \)}{\(\sin \alpha = \frac{\sqrt{3}} 
 {2} \)}{\(\cos \alpha = \frac{\sqrt{5}} 
 {3} \)}{\(\mathop{\mathrm{cotg}}\nolimits \alpha = \sqrt{3}\)}{\(\alpha = 45^{\circ }\)}{}}
\LANGes{
\Question{
El ángulo entre la diagonal espacial y la diagonal de cara del cubo con la arista 
\(a\) es
\(\alpha \).
Entonces vale:}
{\(\mathop{\mathrm{tg}}\nolimits \alpha = \frac{\sqrt{2}} 
 {2} \)}{\(\sin \alpha = \frac{\sqrt{3}} 
 {2} \)}{\(\cos \alpha = \frac{\sqrt{5}} 
 {3} \)}{\(\mathop{\mathrm{cotg}}\nolimits \alpha = \sqrt{3}\)}{\(\alpha = 45^{\circ }\)}{}}
\End

\Data\ID{9000120309}
\Updated{2021-06-29 03:06:37}
\Level{A}
\SubArea{Lines and planes: distances and angles}
\IMG{001203_09-figure0.pdf}
\LANGen{
\Question{
The sides of a rectangular box are \(a = 3\, \mathrm{cm}\),
\(b = 4\, \mathrm{cm}\) and
\(c = 12\, \mathrm{cm}\). The space diagonal
is \(u_{t}\) and the longest
face diagonal is \(u_{s}\).
Find the ratio \(u_{t} : u_{s}\).}
{\(13\sqrt{10} : 40\)}{\(13 : \sqrt{153}\)}{\(13 : 12\)}{\(4\sqrt{10} : 5\)}{\(4\sqrt{10} : 13\)}{}}
\LANGpl{
\Question{
Dany jest prostopadłościan o bokach \(a = 3\, \mathrm{cm}\),
\(b = 4\, \mathrm{cm}\) i
\(c = 12\, \mathrm{cm}\).  Stosunek długości przekątnej 
 \(u_{t}\) do najdłuższej przekątnej ściany bocznej
 \(u_{s}\) jest równy.}
{\(13\sqrt{10} : 40\)}{\(13 : \sqrt{153}\)}{\(13 : 12\)}{\(4\sqrt{10} : 5\)}{\(4\sqrt{10} : 13\)}{}}
\LANGcs{
\Question{
Délky hran kvádru jsou \(a = 3\, \mathrm{cm}\),
\(b = 4\, \mathrm{cm}\),
\(c = 12\, \mathrm{cm}\). Poměr délek tělesové
úhlopříčky \(u_{t}\) a nejdelší
stěnové úhlopříčky \(u_{s}\) 
je roven:}
{\(13\sqrt{10} : 40\)}{\(13 : \sqrt{153}\)}{\(13 : 12\)}{\(4\sqrt{10} : 5\)}{\(4\sqrt{10} : 13\)}{}}
\LANGsk{
\Question{
Dĺžky hrán kvádra sú \(a = 3\, \mathrm{cm}\),
\(b = 4\, \mathrm{cm}\),
\(c = 12\, \mathrm{cm}\). Pomer dĺžok telesovej
uhlopriečky \(u_{t}\) a najdlhšej
stenovej uhlopriečky \(u_{s}\) 
je rovná:}
{\(13\sqrt{10} : 40\)}{\(13 : \sqrt{153}\)}{\(13 : 12\)}{\(4\sqrt{10} : 5\)}{\(4\sqrt{10} : 13\)}{}}
\LANGes{
\Question{
Los lados de un ortoedro son  \(a = 3\, \mathrm{cm}\),
\(b = 4\, \mathrm{cm}\) y
\(c = 12\, \mathrm{cm}\). La diagonal espacial
es \(u_{t}\)  y la  diagonal de cara más larga es \(u_{s}\).
Determina la proporción \(u_{t} : u_{s}\).}
{\(13\sqrt{10} : 40\)}{\(13 : \sqrt{153}\)}{\(13 : 12\)}{\(4\sqrt{10} : 5\)}{\(4\sqrt{10} : 13\)}{}}
\End

\Data\ID{9000121001}
\Updated{2020-07-15 03:07:37}
\Level{A}
\SubArea{Lines and planes: distances and angles}
\IMG{001210_01-figure0.pdf}
\LANGen{
\Question{
In the cube \(ABCDEFGH\) find the
angle between the lines \(AB\) 
and \(AG\).}
{\(54.74^{\circ }\)}{\(60^{\circ }\)}{\(35.26^{\circ }\)}{\(39.23^{\circ }\)}{}{}}
\LANGpl{
\Question{
Dany jest sześcian \(ABCDEFGH\) wyznacz kąt między prostymi
 \(AB\) 
i  \(AG\).}
{\(54.74^{\circ }\)}{\(60^{\circ }\)}{\(35.26^{\circ }\)}{\(39.23^{\circ }\)}{}{}}
\LANGcs{
\Question{
Je dána krychle \(ABCDEFGH\). Vypočítejte
odchylku přímek \(AB\) 
a \(AG\).}
{\(54{,}74^{\circ }
\)}{\(60^{\circ }
\)}{\(35{,}26^{\circ }
\)}{\(39{,}23^{\circ }
\)}{}{}}
\LANGsk{
\Question{
Je daná kocka \(ABCDEFGH\). Vypočítajte
odchýlku priamok \(AB\) 
a \(AG\).}
{\(54{,}74^{\circ }
\)}{\(60^{\circ }
\)}{\(35{,}26^{\circ }
\)}{\(39{,}23^{\circ }
\)}{}{}}
\LANGes{
\Question{
En el cubo \(ABCDEFGH\) determina el ángulo entre las rectas \(AB\) 
y\(AG\).}
{\(54.74^{\circ }\)}{\(60^{\circ }\)}{\(35.26^{\circ }\)}{\(39.23^{\circ }\)}{}{}}
\End

\Data\ID{9000121002}
\Updated{2020-07-15 03:07:37}
\Level{A}
\SubArea{Lines and planes: distances and angles}
\IMG{001210_02-figure0.pdf}
\LANGen{
\Question{
In the cube \(ABCDEFGH\) find the
angle between the lines \(CD\) 
and \(BH\).}
{\(54.74^{\circ }\)}{\(60^{\circ }\)}{\(35.26^{\circ }\)}{\(39.23^{\circ }\)}{}{}}
\LANGpl{
\Question{
Dany jest sześcian \(ABCDEFGH\) wyznacz kąt między prostymi
 \(CD\) 
i \(BH\).}
{\(54.74^{\circ }\)}{\(60^{\circ }\)}{\(35.26^{\circ }\)}{\(39.23^{\circ }\)}{}{}}
\LANGcs{
\Question{
Je dána krychle \(ABCDEFGH\). Vypočítejte
odchylku přímek \(CD\) 
a \(BH\).}
{\(54{,}74^{\circ }
\)}{\(60^{\circ }
\)}{\(35{,}26^{\circ }
\)}{\(39{,}23^{\circ }
\)}{}{}}
\LANGsk{
\Question{
Je daná kocka \(ABCDEFGH\). Vypočítajte
odchýlku priamok \(CD\) 
a \(BH\).}
{\(54{,}74^{\circ }
\)}{\(60^{\circ }
\)}{\(35{,}26^{\circ }
\)}{\(39{,}23^{\circ }
\)}{}{}}
\LANGes{
\Question{
En el cubo \(ABCDEFGH\) determina el ángulo entre las rectas \(CD\) 
y \(BH\).}
{\(54.74^{\circ }\)}{\(60^{\circ }\)}{\(35.26^{\circ }\)}{\(39.23^{\circ }\)}{}{}}
\End

\Data\ID{9000121003}
\Updated{2020-07-15 03:07:37}
\Level{A}
\SubArea{Lines and planes: distances and angles}
\IMG{001210_03-figure0.pdf}
\LANGen{
\Question{
In the cube \(ABCDEFGH\) find the
angle between the lines \(BS_{DH}\) 
and \(AE\), where
the point \(S_{DH}\) is the
center of the edge \(DH\).}
{\(70.53^{\circ }\)}{\(29.47^{\circ }\)}{\(35.26^{\circ }\)}{\(54.74^{\circ }\)}{}{}}
\LANGpl{
\Question{
Dany jest sześcian \(ABCDEFGH\) wyznacz kąt między prostymi
 \(BS_{DH}\) 
i \(AE\), gdzie
punkt  \(S_{DH}\) jest środkiem
krawędzi  \(DH\).}
{\(70.53^{\circ }\)}{\(29.47^{\circ }\)}{\(35.26^{\circ }\)}{\(54.74^{\circ }\)}{}{}}
\LANGcs{
\Question{
Je dána krychle \(ABCDEFGH\). Vypočítejte
odchylku přímek \(BS_{DH}\) 
a \(AE\), kde bod
\(S_{DH}\) je střed
hrany \(DH\).}
{\(70{,}53^{\circ }
\)}{\(29{,}47^{\circ }
\)}{\(35{,}26^{\circ }
\)}{\(54{,}74^{\circ }
\)}{}{}}
\LANGsk{
\Question{
Je daná kocka \(ABCDEFGH\). Vypočítajte
odchýlku priamok \(BS_{DH}\) 
a \(AE\), kde bod
\(S_{DH}\) je stred
hrany \(DH\).}
{\(70{,}53^{\circ }
\)}{\(29{,}47^{\circ }
\)}{\(35{,}26^{\circ }
\)}{\(54{,}74^{\circ }
\)}{}{}}
\LANGes{
\Question{
En el cubo \(ABCDEFGH\)  determina el ángulo entre las rectas \(BS_{DH}\) 
y \(AE\), donde el punto \(S_{DH}\) es el centro de la arista \(DH\).}
{\(70.53^{\circ }\)}{\(29.47^{\circ }\)}{\(35.26^{\circ }\)}{\(54.74^{\circ }\)}{}{}}
\End

\Data\ID{9000121004}
\Updated{2020-07-15 03:07:37}
\Level{A}
\SubArea{Lines and planes: distances and angles}
\IMG{001210_04-figure0.pdf}
\LANGen{
\Question{
In the cube \(ABCDEFGH\) find the
angle between the lines \(S_{AE}S_{HC}\) 
and \(S_{HC}S_{BF}\),
where \(S_{AE}\),
\(S_{HC}\) and
\(S_{BF}\) are the centers
of the segments \(AE\),
\(HC\) and
\(BF\),
respectively.}
{\(53.13^{\circ }\)}{\(26.57^{\circ }\)}{\(60^{\circ }\)}{\(36.87^{\circ }\)}{}{}}
\LANGpl{
\Question{
Dany jest sześcian \(ABCDEFGH\) wyznacz kąt między prostymi
 \(S_{AE}S_{HC}\) 
i \(S_{HC}S_{BF}\),
gdzie \(S_{AE}\),
\(S_{HC}\) i
\(S_{BF}\) to środki odcinków
 \(AE\),
\(HC\) i
\(BF\).}
{\(53.13^{\circ }\)}{\(26.57^{\circ }\)}{\(60^{\circ }\)}{\(36.87^{\circ }\)}{}{}}
\LANGcs{
\Question{
Je dána krychle \(ABCDEFGH\). Vypočítejte
odchylku přímek \(S_{AE}S_{HC}\) 
a \(S_{HC}S_{BF}\), kde
body \(S_{AE}\),
\(S_{HC}\) a
\(S_{BF}\) jsou středy
úseček \(AE\),
\(HC\) a
\(BF\).}
{\(53{,}13^{\circ }
\)}{\(26{,}57^{\circ }
\)}{\(60^{\circ }
\)}{\(36{,}87^{\circ }
\)}{}{}}
\LANGsk{
\Question{
Je daná kocka \(ABCDEFGH\). Vypočítajte
odchýlku priamok \(S_{AE}S_{HC}\) 
a \(S_{HC}S_{BF}\), kde
body \(S_{AE}\),
\(S_{HC}\) a
\(S_{BF}\) sú stredy
úsečiek \(AE\),
\(HC\) a
\(BF\).}
{\(53{,}13^{\circ }
\)}{\(26{,}57^{\circ }
\)}{\(60^{\circ }
\)}{\(36{,}87^{\circ }
\)}{}{}}
\LANGes{
\Question{
En el cubo \(ABCDEFGH\)  determina el ángulo entre las rectas \(S_{AE}S_{HC}\) 
y \(S_{HC}S_{BF}\),
donde\(S_{AE}\),
\(S_{HC}\) y
\(S_{BF}\) son centros de los segmentos
 \(AE\),
\(HC\) y
\(BF\),
respectivamente.}
{\(53.13^{\circ }\)}{\(26.57^{\circ }\)}{\(60^{\circ }\)}{\(36.87^{\circ }\)}{}{}}
\End

\Data\ID{9000121005}
\Updated{2020-07-15 03:07:37}
\Level{A}
\SubArea{Lines and planes: distances and angles}
\IMG{001210_05-figure0.pdf}
\LANGen{
\Question{
In the cube \(ABCDEFGH\) find the
angle between the lines \(AS_{BC}\) 
and \(AD\), where
\(S_{BC}\) is the center
of the edge \(BC\).}
{\(63.43^{\circ }\)}{\(26.57^{\circ }\)}{\(53.13^{\circ }\)}{\(36.87^{\circ }\)}{}{}}
\LANGpl{
\Question{
Dany jest sześcian \(ABCDEFGH\) wyznacz kąt między prostymi
 \(AS_{BC}\) i  \(AD\), gdzie
\(S_{BC}\) to środek krawędzi
 \(BC\).}
{\(63.43^{\circ }\)}{\(26.57^{\circ }\)}{\(53.13^{\circ }\)}{\(36.87^{\circ }\)}{}{}}
\LANGcs{
\Question{
Je dána krychle \(ABCDEFGH\). Vypočítejte
odchylku přímek \(AS_{BC}\) 
a \(AD\), kde bod
\(S_{BC}\) je střed
hrany \(BC\).}
{\(63{,}43^{\circ }
\)}{\(26{,}57^{\circ }
\)}{\(53{,}13^{\circ }
\)}{\(36{,}87^{\circ }
\)}{}{}}
\LANGsk{
\Question{
Je daná kocka \(ABCDEFGH\). Vypočítajte
odchýlku priamok \(AS_{BC}\) 
a \(AD\), kde bod
\(S_{BC}\) je stred
hrany \(BC\).}
{\(63{,}43^{\circ }
\)}{\(26{,}57^{\circ }
\)}{\(53{,}13^{\circ }
\)}{\(36{,}87^{\circ }
\)}{}{}}
\LANGes{
\Question{
En el cubo \(ABCDEFGH\) determina el ángulo entre las rectas \(AS_{BC}\) 
y \(AD\), donde
\(S_{BC}\) es el centro de la arista \(BC\).}
{\(63.43^{\circ }\)}{\(26.57^{\circ }\)}{\(53.13^{\circ }\)}{\(36.87^{\circ }\)}{}{}}
\End

\Data\ID{9000121006}
\Updated{2020-07-15 03:07:37}
\Level{A}
\SubArea{Lines and planes: distances and angles}
\IMG{001210_06-figure0.pdf}
\LANGen{
\Question{
In the cube \(ABCDEFGH\) find the
angle between the lines \(BG\) 
and \(GD\).}
{\(60^{\circ }\)}{\(45^{\circ }\)}{\(36.87^{\circ }\)}{\(53.13^{\circ }\)}{}{}}
\LANGpl{
\Question{
Dany jest sześcian \(ABCDEFGH\) wyznacz kąt między prostymi
 \(BG\) 
i \(GD\).}
{\(60^{\circ }\)}{\(45^{\circ }\)}{\(36.87^{\circ }\)}{\(53.13^{\circ }\)}{}{}}
\LANGcs{
\Question{
Je dána krychle \(ABCDEFGH\). Vypočítejte
odchylku přímek \(BG\) 
a \(GD\).}
{\(60^{\circ }
\)}{\(45^{\circ }
\)}{\(36{,}87^{\circ }
\)}{\(53{,}13^{\circ }
\)}{}{}}
\LANGsk{
\Question{
Je daná kocka \(ABCDEFGH\). Vypočítajte
odchýlku priamok \(BG\) 
a \(GD\).}
{\(60^{\circ }
\)}{\(45^{\circ }
\)}{\(36{,}87^{\circ }
\)}{\(53{,}13^{\circ }
\)}{}{}}
\LANGes{
\Question{
En el cubo \(ABCDEFGH\) determina el ángulo entre las rectas\(BG\) 
y \(GD\).}
{\(60^{\circ }\)}{\(45^{\circ }\)}{\(36.87^{\circ }\)}{\(53.13^{\circ }\)}{}{}}
\End

\Data\ID{9000121007}
\Updated{2020-07-15 03:07:37}
\Level{A}
\SubArea{Lines and planes: distances and angles}
\IMG{001210_07-figure0.pdf}
\LANGen{
\Question{
In the cube \(ABCDEFGH\) find the
angle between the lines \(S_{BE}S_{AH}\) 
and \(HC\), where
\(S_{BE}\) and
\(S_{AH}\) are centers of
the segments \(BE\) 
and \(AH\),
respectively.}
{\(60^{\circ }\)}{\(26.57^{\circ }\)}{\(45^{\circ }\)}{\(53.13^{\circ }\)}{}{}}
\LANGpl{
\Question{
Dany jest sześcian \(ABCDEFGH\) wyznacz kąt między prostymi
 \(S_{BE}S_{AH}\) 
i \(HC\), gdzie
\(S_{BE}\) i
\(S_{AH}\) to środki odcinków \(BE\) 
i \(AH\).}
{\(60^{\circ }\)}{\(26.57^{\circ }\)}{\(45^{\circ }\)}{\(53.13^{\circ }\)}{}{}}
\LANGcs{
\Question{
Je dána krychle \(ABCDEFGH\). Vypočítejte
odchylku přímek \(S_{BE}S_{AH}\) 
a \(HC\), kde
body \(S_{BE}\) a
\(S_{AH}\) jsou středy
úseček \(BE\) 
a \(AH\).}
{\(60^{\circ }
\)}{\(26{,}57^{\circ }
\)}{\(45^{\circ }
\)}{\(53{,}13^{\circ }
\)}{}{}}
\LANGsk{
\Question{
Je daná kocka \(ABCDEFGH\). Vypočítajte
odchýlku priamok \(S_{BE}S_{AH}\) 
a \(HC\), kde
body \(S_{BE}\) a
\(S_{AH}\) sú stredy
úsečiek \(BE\) 
a \(AH\).}
{\(60^{\circ }
\)}{\(26{,}57^{\circ }
\)}{\(45^{\circ }
\)}{\(53{,}13^{\circ }
\)}{}{}}
\LANGes{
\Question{
En el cubo \(ABCDEFGH\) determina el ángulo entre las rectas \(S_{BE}S_{AH}\) 
y \(HC\), donde
\(S_{BE}\) y
\(S_{AH}\) son centros de los segmentos \(BE\) 
y \(AH\),
respectivamente.}
{\(60^{\circ }\)}{\(26.57^{\circ }\)}{\(45^{\circ }\)}{\(53.13^{\circ }\)}{}{}}
\End

\Data\ID{9000121008}
\Updated{2020-07-15 03:07:37}
\Level{A}
\SubArea{Lines and planes: distances and angles}
\IMG{001210_08-figure0.pdf}
\LANGen{
\Question{
In the cube \(ABCDEFGH\) find the
angle between the lines \(DB\) 
and \(AG\).}
{\(90^{\circ }\)}{\(45^{\circ }\)}{\(35.26^{\circ }\)}{\(53.13^{\circ }\)}{}{}}
\LANGpl{
\Question{
Dany jest sześcian \(ABCDEFGH\) wskaż kąt między prostymi
\(DB\) 
i \(AG\).}
{\(90^{\circ }\)}{\(45^{\circ }\)}{\(35.26^{\circ }\)}{\(53.13^{\circ }\)}{}{}}
\LANGcs{
\Question{
Je dána krychle \(ABCDEFGH\). Vypočítejte
odchylku přímek \(DB\) 
a \(AG\).}
{\(90^{\circ }
\)}{\(45^{\circ }
\)}{\(35{,}26^{\circ }
\)}{\(53{,}13^{\circ }
\)}{}{}}
\LANGsk{
\Question{
Je daná kocka \(ABCDEFGH\). Vypočítajte
odchýlku priamok \(DB\) 
a \(AG\).}
{\(90^{\circ }
\)}{\(45^{\circ }
\)}{\(35{,}26^{\circ }
\)}{\(53{,}13^{\circ }
\)}{}{}}
\LANGes{
\Question{
En el cubo \(ABCDEFGH\) determina el ángulo entre las rectas \(DB\) 
y \(AG\).}
{\(90^{\circ }\)}{\(45^{\circ }\)}{\(35.26^{\circ }\)}{\(53.13^{\circ }\)}{}{}}
\End

\Data\ID{9000121009}
\Updated{2020-07-15 03:07:37}
\Level{A}
\SubArea{Lines and planes: distances and angles}
\IMG{001210_09-figure0.pdf}
\LANGen{
\Question{
In the cube \(ABCDEFGH\) find the
angle between the lines \(S_{HD}S_{FC}\) 
and \(AB\), where
the points \(S_{HD}\) 
and \(S_{FC}\) are
centers of \(HD\) 
and \(FC\),
respectively.}
{\(26.57^{\circ }\)}{\(45^{\circ }\)}{\(54.74^{\circ }\)}{\(60^{\circ }\)}{}{}}
\LANGpl{
\Question{
Dany jest sześcian \(ABCDEFGH\) 
wskaż kąt między prostymi \(S_{HD}S_{FC}\) 
i \(AB\), gdzie punkty
 \(S_{HD}\) 
i \(S_{FC}\) to środki 
\(HD\) 
i \(FC\).}
{\(26.57^{\circ }\)}{\(45^{\circ }\)}{\(54.74^{\circ }\)}{\(60^{\circ }\)}{}{}}
\LANGcs{
\Question{
Je dána krychle \(ABCDEFGH\). Vypočítejte
odchylku přímek \(S_{HD}S_{FC}\) 
a \(AB\), kde
body \(S_{HD}\) a
\(S_{FC}\) jsou středy
úseček \(HD\) 
a \(FC\).}
{\(26{,}57^{\circ }
\)}{\(45^{\circ }
\)}{\(54{,}74^{\circ }
\)}{\(60^{\circ }
\)}{}{}}
\LANGsk{
\Question{
Je daná kocka \(ABCDEFGH\). Vypočítajte
odchýlku priamok \(S_{HD}S_{FC}\) 
a \(AB\), kde
body \(S_{HD}\) a
\(S_{FC}\) sú stredy
úsečiek \(HD\) 
a \(FC\).}
{\(26{,}57^{\circ }
\)}{\(45^{\circ }
\)}{\(54{,}74^{\circ }
\)}{\(60^{\circ }
\)}{}{}}
\LANGes{
\Question{
En el cubo \(ABCDEFGH\) determina el ángulo entre las rectas \(S_{HD}S_{FC}\) 
y \(AB\), donde
los puntos \(S_{HD}\) 
y \(S_{FC}\) son centros de los segmentos \(HD\) 
y \(FC\),
respectivamente.}
{\(26.57^{\circ }\)}{\(45^{\circ }\)}{\(54.74^{\circ }\)}{\(60^{\circ }\)}{}{}}
\End

\Data\ID{9000121010}
\Updated{2020-07-15 03:07:37}
\Level{A}
\SubArea{Lines and planes: distances and angles}
\IMG{001210_10-figure0.pdf}
\LANGen{
\Question{
In the cube \(ABCDEFGH\) find the
angle between the lines \(ES_{CG}\) 
and \(AG\), where
\(S_{CG}\) is the center
of the \(CG\).}
{\(54.74^{\circ }\)}{\(19.47^{\circ }\)}{\(35.26^{\circ }\)}{\(60^{\circ }\)}{}{}}
\LANGpl{
\Question{
Dany jest sześcian  \(ABCDEFGH\) wskaż kąt między prostymi
 \(ES_{CG}\) 
i \(AG\), gdzie
\(S_{CG}\) to środek
 \(CG\).}
{\(54.74^{\circ }\)}{\(19.47^{\circ }\)}{\(35.26^{\circ }\)}{\(60^{\circ }\)}{}{}}
\LANGcs{
\Question{
Je dána krychle \(ABCDEFGH\). Vypočítejte
odchylku přímek \(ES_{CG}\) 
a \(AG\), kde bod
\(S_{CG}\) je střed
hrany \(CG\).}
{\(54{,}74^{\circ }
\)}{\(19{,}47^{\circ }
\)}{\(35{,}26^{\circ }
\)}{\(60^{\circ }
\)}{}{}}
\LANGsk{
\Question{
Je daná kocka \(ABCDEFGH\). Vypočítajte
odchýlku priamok \(ES_{CG}\) 
a \(AG\), kde bod
\(S_{CG}\) je stred
hrany \(CG\).}
{\(54{,}74^{\circ }
\)}{\(19{,}47^{\circ }
\)}{\(35{,}26^{\circ }
\)}{\(60^{\circ }
\)}{}{}}
\LANGes{
\Question{
En el cubo \(ABCDEFGH\) determina el ángulo entre las rectas \(ES_{CG}\) 
y \(AG\), donde
\(S_{CG}\) es el centro de \(CG\).}
{\(54.74^{\circ }\)}{\(19.47^{\circ }\)}{\(35.26^{\circ }\)}{\(60^{\circ }\)}{}{}}
\End

\Data\ID{1103025301}
\Updated{2020-07-15 03:07:12}
\Level{B}
\SubArea{Lines and planes: distances and angles}
\IMG{1103025301.pdf}
\LANGen{
\Question{
The base \( ABCD \) of a square pyramid \( ABCDV \)  has an edge of \( 6\,\mathrm{cm} \). The height of the pyramid is \( 4\,\mathrm{cm} \). Find the distance between the point \( V \) and the line \( BC \).}
{\( 5\,\mathrm{cm} \)}{\( \sqrt{52}\,\mathrm{cm} \)}{\( 25\,\mathrm{cm} \)}{\( \sqrt{10}\,\mathrm{cm} \)}{}{}}
\LANGpl{
\Question{
Krawędź podstawy \( ABCD \)  ostrosłupa  czworokątnego \( ABCDV \)  jest równa \( 6\,\mathrm{cm} \). Wysokość ostrosłupa wynosi \( 4\,\mathrm{cm} \). Oblicz odległość od punktu \( V \) do prostej \( BC \).}
{\( 5\,\mathrm{cm} \)}{\( \sqrt{52}\,\mathrm{cm} \)}{\( 25\,\mathrm{cm} \)}{\( \sqrt{10}\,\mathrm{cm} \)}{}{}}
\LANGcs{
\Question{
V pravidelném čtyřbokém jehlanu \( ABCDV \)  s hlavním vrcholem \( V \) má podstavná hrana velikost \( 6\,\mathrm{cm} \) a výška jehlanu \( 4\,\mathrm{cm} \). Určete vzdálenost bodu \( V \)  od přímky \( BC \).}
{\( 5\,\mathrm{cm} \)}{\( \sqrt{52}\,\mathrm{cm} \)}{\( 25\,\mathrm{cm} \)}{\( \sqrt{10}\,\mathrm{cm} \)}{}{}}
\LANGsk{
\Question{
V pravidelnom štvorbokom ihlane \( ABCDV \)  s hlavným vrcholom \( V \) má hrana podstavy veľkosť \( 6\,\mathrm{cm} \) a výška ihlana \( 4\,\mathrm{cm} \). Určte vzdialenosť bodu \( V \)  od priamky \( BC \).}
{\( 5\,\mathrm{cm} \)}{\( \sqrt{52}\,\mathrm{cm} \)}{\( 25\,\mathrm{cm} \)}{\( \sqrt{10}\,\mathrm{cm} \)}{}{}}
\LANGes{
\Question{
En una pirámide cuadrangular regular \( ABCDV \)  con el vértice principal \( V \) mide la arista básica  \( 6\,\mathrm{cm} \) y la altura de la pirámide es \( 4\,\mathrm{cm} \). Determina la distancia entre el punto  \( V \)  y la recta \( BC \).}
{\( 5\,\mathrm{cm} \)}{\( \sqrt{52}\,\mathrm{cm} \)}{\( 25\,\mathrm{cm} \)}{\( \sqrt{10}\,\mathrm{cm} \)}{}{}}
\End

\Data\ID{1103025302}
\Updated{2020-07-15 03:07:12}
\Level{B}
\SubArea{Lines and planes: distances and angles}
\IMG{1103025302.pdf}
\LANGen{
\Question{
The base \( ABCD \) of a square pyramid  \( ABCDV \) has an edge of \( 6\,\mathrm{cm} \). The height of the pyramid is \( 4\,\mathrm{cm} \). Find the distance between the line \( S_{VB}S_{VC}\) and the line \( BC \). The points $S_{VB}$ and $S_{VC}$ are the midpoints of $VB$ and $VC$, respectively.}
{\( 2.5\,\mathrm{cm} \)}{\( 2\,\mathrm{cm} \)}{\( \frac{\sqrt{52}}2\,\mathrm{cm} \)}{\( \frac{25}2\,\mathrm{cm} \)}{}{}}
\LANGpl{
\Question{
Krawędź podstawy \( ABCD \) ostrosłupa  czworokątnego   \( ABCDV \) jest równa  \( 6\,\mathrm{cm} \). Wysokość ostrosłupa wynosi \( 4\,\mathrm{cm} \). Oblicz odległość pomiędzy prostą  \( S_{VB}S_{VC}\), a prostą \( BC \). Punkty $S_{VB}$ i $S_{VC}$ są środkiem krawędzi $VB$ i $VC$.}
{\( 2{,}5\,\mathrm{cm} \)}{\( 2\,\mathrm{cm} \)}{\( \frac{\sqrt{52}}2\,\mathrm{cm} \)}{\( \frac{25}2\,\mathrm{cm} \)}{}{}}
\LANGcs{
\Question{
V pravidelném čtyřbokém jehlanu \( ABCDV \)  s hlavním vrcholem \( V \) má podstavná hrana velikost \( 6\,\mathrm{cm} \) a výška jehlanu \( 4\,\mathrm{cm} \). Určete vzdálenost přímek \( S_{VB}S_{VC}\) a \( BC \). Bod $S_{VB}$ je středem hrany $VB$ a bod $S_{VC}$ je středem hrany $VC$.}
{\( 2{,}5\,\mathrm{cm} \)}{\( 2\,\mathrm{cm} \)}{\( \frac{\sqrt{52}}2\,\mathrm{cm} \)}{\( \frac{25}2\,\mathrm{cm} \)}{}{}}
\LANGsk{
\Question{
V pravidelnom štvorbokom ihlane \( ABCDV \)  s hlavným vrcholom \( V \) má hrana podstavy veľkosť \( 6\,\mathrm{cm} \) a výška ihlana \( 4\,\mathrm{cm} \). Určte vzdialenosť priamok \( S_{VB}S_{VC}\) a \( BC \). Bod $S_{VB}$ je stred hrany $VB$ a bod $S_{VC}$ je stred hrany $VC$.}
{\( 2{,}5\,\mathrm{cm} \)}{\( 2\,\mathrm{cm} \)}{\( \frac{\sqrt{52}}2\,\mathrm{cm} \)}{\( \frac{25}2\,\mathrm{cm} \)}{}{}}
\LANGes{
\Question{
En una pirámide cuadrangular regular \( ABCDV \)  con el vértice principal \( V \) mide la arista básica \( 6\,\mathrm{cm} \) y la altura de la pirámide es \( 4\,\mathrm{cm} \). Determina la distancia entre las rectas \( S_{VB}S_{VC}\) y \( BC \). El punto $S_{VB}$ es el centro de la arista $VB$ y el punto $S_{VC}$ es el centro de la arista $VC$.}
{\( 2.5\,\mathrm{cm} \)}{\( 2\,\mathrm{cm} \)}{\( \frac{\sqrt{52}}2\,\mathrm{cm} \)}{\( \frac{25}2\,\mathrm{cm} \)}{}{}}
\End

\Data\ID{1103025303}
\Updated{2020-07-15 03:07:12}
\Level{B}
\SubArea{Lines and planes: distances and angles}
\IMG{1103025303.pdf}
\LANGen{
\Question{
The base \( ABCD \) of a square pyramid \( ABCDV \) has an edge of \( 6\,\mathrm{cm} \). The height of the pyramid is \( 4\,\mathrm{cm} \). Find the distance between the line \( S_{VA}S_{VC} \) and the line \( AC \). The points $S_{VA}$ and $S_{VC}$ are the midpoints of $VA$ and $VC$, respectively.}
{\( 2\,\mathrm{cm} \)}{\( 2.5\,\mathrm{cm} \)}{\( \frac{\sqrt{52}}2\,\mathrm{cm} \)}{\( 4\,\mathrm{cm} \)}{}{}}
\LANGpl{
\Question{
Krawędź postawy  \( ABCD \) ostrosłupa  czworokątnego  \( ABCDV \) jest równa \( 6\,\mathrm{cm} \). Wysokość ostrosłupa wynosi \( 4\,\mathrm{cm} \). Oblicz kąt między prostą  \( S_{VA}S_{VC} \), a prostą  \( AC \). Punkty $S_{VA}$ i $S_{VC}$ są środkiem krawędzi $VA$ i $VC$.}
{\( 2\,\mathrm{cm} \)}{\( 2{,}5\,\mathrm{cm} \)}{\( \frac{\sqrt{52}}2\,\mathrm{cm} \)}{\( 4\,\mathrm{cm} \)}{}{}}
\LANGcs{
\Question{
V pravidelném čtyřbokém jehlanu \( ABCDV \)  s hlavním vrcholem \( V \) má podstavná hrana velikost \( 6\,\mathrm{cm} \) a výška jehlanu \( 4\,\mathrm{cm} \). Určete vzdálenost přímek \( S_{VA}S_{VC} \) a \( AC \). Bod $S_{VA}$ je středem hrany $VA$ a bod $S_{VC}$ je středem hrany $VC$.}
{\( 2\,\mathrm{cm} \)}{\( 2{,}5\,\mathrm{cm} \)}{\( \frac{\sqrt{52}}2\,\mathrm{cm} \)}{\( 4\,\mathrm{cm} \)}{}{}}
\LANGsk{
\Question{
V pravidelnom štvorbokom ihlane \( ABCDV \)  s hlavným vrcholom \( V \) má hrana podstavy veľkosť \( 6\,\mathrm{cm} \) a výška ihlanu \( 4\,\mathrm{cm} \). Určte vzdialenosť priamok \( S_{VA}S_{VC} \) a \( AC \). Bod $S_{VA}$ je stred hrany $VA$ a bod $S_{VC}$ je stred hrany $VC$.}
{\( 2\,\mathrm{cm} \)}{\( 2{,}5\,\mathrm{cm} \)}{\( \frac{\sqrt{52}}2\,\mathrm{cm} \)}{\( 4\,\mathrm{cm} \)}{}{}}
\LANGes{
\Question{
En una pirámide cuadrangular regular \( ABCDV \) con el vértice principal \( V \) mide la arista básica \( 6\,\mathrm{cm} \) y la altura de la pirámide es \( 4\,\mathrm{cm} \). Determina la distancia entre las rectas \( S_{VA}S_{VC} \) y \( AC \). El punto $S_{VA}$ es el centro de la arista  $VA$ y el punto $S_{VC}$ es el centro de la arista $VC$.}
{\( 2\,\mathrm{cm} \)}{\( 2.5\,\mathrm{cm} \)}{\( \frac{\sqrt{52}}2\,\mathrm{cm} \)}{\( 4\,\mathrm{cm} \)}{}{}}
\End

\Data\ID{1103025304}
\Updated{2020-07-15 03:07:12}
\Level{B}
\SubArea{Lines and planes: distances and angles}
\IMG{1103025304.pdf}
\LANGen{
\Question{
The base \( ABCD \) of a square pyramid \( ABCDV \) has an edge of \( 8\,\mathrm{cm} \). The height of the pyramid is \( 9\,\mathrm{cm} \). Find the distance between the line \( S_{VA}S_{VD} \) and the line \( BC \). The points $S_{VA}$ and $S_{VD}$ are the midpoints of $VA$ and $VD$, respectively.}
{\( 7.5\,\mathrm{cm} \)}{\( \frac23\sqrt{97}\,\mathrm{cm} \)}{\( \frac{\sqrt{97}}2\,\mathrm{cm} \)}{\( \sqrt{17}\,\mathrm{cm} \)}{}{}}
\LANGpl{
\Question{
Krawędź postawy \( ABCD \) ostrosłupa prawidłowego czworokątnego  \( ABCDV \) jest równa \( 8\,\mathrm{cm} \). Wysokość ostrosłupa wynosi \( 9\,\mathrm{cm} \). Oblicz odległość między prostą  \( S_{VA}S_{VD} \), a prostą \( BC \). Punkty $S_{VA}$ i $S{VD}$ są środkiem krawędzi $VA$ i $VD$.}
{\( 7{,}5\,\mathrm{cm} \)}{\( \frac23\sqrt{97}\,\mathrm{cm} \)}{\( \frac{\sqrt{97}}2\,\mathrm{cm} \)}{\( \sqrt{17}\,\mathrm{cm} \)}{}{}}
\LANGcs{
\Question{
V pravidelném čtyřbokém jehlanu \( ABCDV \)  s hlavním vrcholem \( V \) má podstavná hrana velikost \( 8\,\mathrm{cm} \) a výška jehlanu \( 9\,\mathrm{cm} \). Určete vzdálenost přímek \( S_{VA}S_{VD} \) a \( BC \). Bod $S_{VA}$ je středem hrany $VA$ a bod $S_{VD}$ je středem hrany $VD$.}
{\( 7{,}5\,\mathrm{cm} \)}{\( \frac23\sqrt{97}\,\mathrm{cm} \)}{\( \frac{\sqrt{97}}2\,\mathrm{cm} \)}{\( \sqrt{17}\,\mathrm{cm} \)}{}{}}
\LANGsk{
\Question{
V pravidelnom štvorbokom ihlane \( ABCDV \)  s hlavným vrcholom \( V \) má hrana podstavy veľkosť \( 8\,\mathrm{cm} \) a výška ihlanu \( 9\,\mathrm{cm} \). Určte vzdialenosť priamok \( S_{VA}S_{VD} \) a \( BC \). Bod $S_{VA}$ je stred hrany $VA$ a bod $S_{VD}$ je stred hrany $VD$.}
{\( 7{,}5\,\mathrm{cm} \)}{\( \frac23\sqrt{97}\,\mathrm{cm} \)}{\( \frac{\sqrt{97}}2\,\mathrm{cm} \)}{\( \sqrt{17}\,\mathrm{cm} \)}{}{}}
\LANGes{
\Question{
En una pirámide cuadrangular regula \( ABCDV \)  con el vértice principal \( V \) mide la arista básica \( 8\,\mathrm{cm} \) y la altura de la pirámide es \( 9\,\mathrm{cm} \). Determina la distancia entre las rectas \( S_{VA}S_{VD} \) y \( BC \). El punto $S_{VA}$ es el centro de la arista $VA$ y el punto $S_{VD}$ es el centro de la arista $VD$.}
{\( 7.5\,\mathrm{cm} \)}{\( \frac23\sqrt{97}\,\mathrm{cm} \)}{\( \frac{\sqrt{97}}2\,\mathrm{cm} \)}{\( \sqrt{17}\,\mathrm{cm} \)}{}{}}
\End

\Data\ID{1103025305}
\Updated{2020-07-15 03:07:12}
\Level{B}
\SubArea{Lines and planes: distances and angles}
\IMG{1103025305.pdf}
\LANGen{
\Question{
The base \( ABCD \) of a square pyramid \( ABCDV \) has an edge of \( 4\,\mathrm{cm} \). The height of the pyramid is \( 6\,\mathrm{cm} \). Find the distance between the point \( A \) and the point \( S_{VC} \),   where \( S_{VC} \) is the midpoint of the edge \( VC \).}
{\( 3\sqrt{3}\,\mathrm{cm} \)}{\( 4\sqrt{3}\,\mathrm{cm} \)}{\( \sqrt{10}\,\mathrm{cm} \)}{\( 3\sqrt{10}\,\mathrm{cm} \)}{}{}}
\LANGpl{
\Question{
Krawędź podstawy  \( ABCD \) ostrosłupa  czworokątnego  \( ABCDV \) jest równa \( 4\,\mathrm{cm} \). Wysokość ostrosłupa wynosi  \( 6\,\mathrm{cm} \). Oblicz odległość między punktami \( A \) i \( S_{VC} \),   gdzie \( S_{VC} \) jest środkiem krawędzi \( VC \).}
{\( 3\sqrt{3}\,\mathrm{cm} \)}{\( 4\sqrt{3}\,\mathrm{cm} \)}{\( \sqrt{10}\,\mathrm{cm} \)}{\( 3\sqrt{10}\,\mathrm{cm} \)}{}{}}
\LANGcs{
\Question{
V pravidelném čtyřbokém jehlanu \( ABCDV \)  s hlavním vrcholem \( V \) má podstavná hrana velikost \( 4\,\mathrm{cm} \) a výška jehlanu \( 6\,\mathrm{cm} \). Určete vzdálenost bodů \( A \) a \( S_{VC} \) (bod \( S_{VC} \) je střed boční hrany \( VC \)):}
{\( 3\sqrt{3}\,\mathrm{cm} \)}{\( 4\sqrt{3}\,\mathrm{cm} \)}{\( \sqrt{10}\,\mathrm{cm} \)}{\( 3\sqrt{10}\,\mathrm{cm} \)}{}{}}
\LANGsk{
\Question{
V pravidelnom štvorbokom ihlanu \( ABCDV \)  s hlavným vrcholom \( V \) má hrana podstavy veľkosť \( 4\,\mathrm{cm} \) a výška ihlanu \( 6\,\mathrm{cm} \). Určte vzdialenosť bodov \( A \) a \( S_{VC} \) (bod \( S_{VC} \) je stred boční hrany \( VC \)):}
{\( 3\sqrt{3}\,\mathrm{cm} \)}{\( 4\sqrt{3}\,\mathrm{cm} \)}{\( \sqrt{10}\,\mathrm{cm} \)}{\( 3\sqrt{10}\,\mathrm{cm} \)}{}{}}
\LANGes{
\Question{
En una pirámide cuadrangular regula  \( ABCDV \)  con el vértice principal \( V \) mide la arista básica  \( 4\,\mathrm{cm} \) y la altura de la pirámide es \( 6\,\mathrm{cm} \). Determina la distancia entre los  puntos \( A \) y \( S_{VC} \) (el punto \( S_{VC} \) es el centro de la arista \( VC \)):}
{\( 3\sqrt{3}\,\mathrm{cm} \)}{\( 4\sqrt{3}\,\mathrm{cm} \)}{\( \sqrt{10}\,\mathrm{cm} \)}{\( 3\sqrt{10}\,\mathrm{cm} \)}{}{}}
\End

\Data\ID{1103025306}
\Updated{2020-07-15 03:07:12}
\Level{B}
\SubArea{Lines and planes: distances and angles}
\IMG{1103025306.pdf}
\LANGen{
\Question{
The base \( ABCD \) of a square pyramid \( ABCDV \) has an edge of \( 6\,\mathrm{cm} \). The height of the pyramid is \( 3\sqrt2\,\mathrm{cm} \). Find the distance between the point \( A \) and the line \( BV \).}
{\( 3\sqrt3\,\mathrm{cm} \)}{\( 3\sqrt2\,\mathrm{cm} \)}{\( \frac32\sqrt3\,\mathrm{cm} \)}{\( 3\sqrt6\,\mathrm{cm} \)}{}{}}
\LANGpl{
\Question{
Krawędź podstawy  \( ABCD \) ostrosłupa  czworokątnego  \( ABCDV \) jest równa \( 6\,\mathrm{cm} \). Wysokość ostrosłupa wynosi \( 3\sqrt2\,\mathrm{cm} \). Oblicz odległość między punktem \( A \), a prostą  \( BV \).}
{\( 3\sqrt3\,\mathrm{cm} \)}{\( 3\sqrt2\,\mathrm{cm} \)}{\( \frac32\sqrt3\,\mathrm{cm} \)}{\( 3\sqrt6\,\mathrm{cm} \)}{}{}}
\LANGcs{
\Question{
V pravidelném čtyřbokém jehlanu \( ABCDV \)  s hlavním vrcholem \( V \) má podstavná hrana velikost \( 6\,\mathrm{cm} \) a výška jehlanu \( 3\sqrt2\,\mathrm{cm} \). Určete vzdálenost bodu \( A \) od přímky \( BV \):}
{\( 3\sqrt3\,\mathrm{cm} \)}{\( 3\sqrt2\,\mathrm{cm} \)}{\( \frac32\sqrt3\,\mathrm{cm} \)}{\( 3\sqrt6\,\mathrm{cm} \)}{}{}}
\LANGsk{
\Question{
V pravidelnom štvorbokom ihlane \( ABCDV \)  s hlavným vrcholom \( V \) má hrana podstavy veľkosť \( 6\,\mathrm{cm} \) a výška ihlanu \( 3\sqrt2\,\mathrm{cm} \). Určte vzdialenosť bodu \( A \) od priamky \( BV \):}
{\( 3\sqrt3\,\mathrm{cm} \)}{\( 3\sqrt2\,\mathrm{cm} \)}{\( \frac32\sqrt3\,\mathrm{cm} \)}{\( 3\sqrt6\,\mathrm{cm} \)}{}{}}
\LANGes{
\Question{
En una pirámide cuadrangular regula  \( ABCDV \)  con el vértice principal \( V \) mide la arista básica \( 6\,\mathrm{cm} \) y la altura de la pirámide es \( 3\sqrt2\,\mathrm{cm} \). Determina la distancia entre el punto \( A \) y la recta \( BV \):}
{\( 3\sqrt3\,\mathrm{cm} \)}{\( 3\sqrt2\,\mathrm{cm} \)}{\( \frac32\sqrt3\,\mathrm{cm} \)}{\( 3\sqrt6\,\mathrm{cm} \)}{}{}}
\End

\Data\ID{9000046408}
\Updated{2020-07-15 03:07:34}
\Level{B}
\SubArea{Lines and planes: distances and angles}
\LANGen{
\Question{
Consider a cone of base radius \(r\) 
and a special shape: the shape is such that the volume of the cone is related to the base radius
by the formula \(V =\pi r^{3}\).
Find the angle between the side of the cone and the base. Round your answer to two
decimal places.}
{\(71.57^{\circ }\)}{\(45^{\circ }\)}{\(63.43^{\circ }\)}{}{}{}}
\LANGpl{
\Question{
Objętość stożka o promieniu  \(r\) 
wynosi  \(V =\pi r^{3}\).
Oblicz kąt pomiędzy ścianą boczną, a jego podstawą. Zaokrągli  odpowiedź do dwóch miejsc po przecinku.}
{\(71.57^{\circ }\)}{\(45^{\circ }\)}{\(63.43^{\circ }\)}{}{}{}}
\LANGcs{
\Question{
Objem rotačního kužele s poloměrem podstavy
\(r\) je
\(V =\pi r^{3}\).
Určete odchylku jeho strany od roviny podstavy (výsledek je zaokrouhlen na
\(2\) 
desetinná místa).}
{\(71{,}57^{\circ }\)}{\(45^{\circ }\)}{\(63{,}43^{\circ }\)}{}{}{}}
\LANGsk{
\Question{
Objem rotačného kužeľa s polomerom podstavy
\(r\) je
\(V =\pi r^{3}\).
Určte odchýlku jeho strany od roviny podstavy (výsledok je zaokrúhlený na
\(2\) 
desatinné miesta).}
{\(71{,}57^{\circ }\)}{\(45^{\circ }\)}{\(63{,}43^{\circ }\)}{}{}{}}
\LANGes{
\Question{
El volumen del cono con el radio \(r \)
es
\(V =\pi r^{3}\).
Determina el ángulo entre la cara y la base del cono. Redondea el resultado a dos posiciones decimales.}
{\(71.57^{\circ }\)}{\(45^{\circ }\)}{\(63.43^{\circ }\)}{}{}{}}
\End

\Data\ID{9000046409}
\Updated{2021-06-29 03:06:53}
\Level{B}
\SubArea{Lines and planes: distances and angles}
\LANGen{
\Question{
The base of a pyramid is a square with the side of \(2\, \mathrm{cm}\). The height of the pyramid is \(4\, \mathrm{cm}\). Find the angle between the lateral side of the pyramid and the base. Round your result to two decimal places.}
{\(75.96^{\circ }\)}{\(70.52^{\circ }\)}{\(79.98^{\circ }\)}{}{}{}}
\LANGpl{
\Question{
Podstawą ostrosłupa jest kwadrat o boku
\(2\, \mathrm{cm}\). Wysokość ostrosłupa jest równa
 \(4\, \mathrm{cm}\).
Oblicz kąt między ścianą boczna ostrosłupa,  a jego podstawą. Zaokrągli odpowiedź do dwóch miejsc po przecinku.}
{\(75.96^{\circ }\)}{\(70.52^{\circ }\)}{\(79.98^{\circ }\)}{}{}{}}
\LANGcs{
\Question{
Pravidelný čtyřboký jehlan má podstavnou hranu o velikosti
\(2\, \mathrm{cm}\) a výšku
o velikosti \(4\, \mathrm{cm}\).
Určete odchylku jeho boční stěny od roviny podstavy (výsledek je zaokrouhlen
na \(2\) 
desetinná místa).}
{\(75{,}96^{\circ }\)}{\(70{,}52^{\circ }\)}{\(79{,}98^{\circ }\)}{}{}{}}
\LANGsk{
\Question{
Pravidelný štvorboký ihlan má hranu podstavy o veľkosti
\(2\, \mathrm{cm}\) a výšku
o veľkosti \(4\, \mathrm{cm}\).
Určte odchýlku jeho bočnej steny od roviny podstavy (výsledok je zaokrúhlený
na \(2\) 
desatinné miesta).}
{\(75{,}96^{\circ }\)}{\(70{,}52^{\circ }\)}{\(79{,}98^{\circ }\)}{}{}{}}
\LANGes{
\Question{
La base de una pirámide es un cuadrado y su lado mide
\(2 \, \mathrm {cm} \). La altura de
la pirámide es \(4 \, \mathrm {cm} \).
Determina el ángulo entre el lado de la pirámide y la base. Redondea el resultado a dos posiciones decimales.}
{\(75.96^{\circ }\)}{\(70.52^{\circ }\)}{\(79.98^{\circ }\)}{}{}{}}
\End

\Data\ID{9000128801}
\Updated{2020-07-15 03:07:37}
\Level{B}
\SubArea{Lines and planes: distances and angles}
\IMG{001288_01-figure0.pdf}
\LANGen{
\Question{
The base \(ABCD\) of a square
pyramid \(ABCDV \) has side
\(6\, \mathrm{cm}\). The height of the
pyramid is \(4\, \mathrm{cm}\). The point
\(M\) is the middle of the
side \(CV \). Find the distance
between the point \(M\) 
and the plane \(ABC\).}
{\(2\, \mathrm{cm}\)}{\(\frac{\sqrt{34}}
 {2} \, \mathrm{cm}\)}{\(\frac{5}
{2}\, \mathrm{cm}\)}{}{}{}}
\LANGpl{
\Question{
Bok podstawy \(ABCD\) ostrosłupa prawidłowego czworokątnego 
 \(ABCDV \) jest równy 
\(6\, \mathrm{cm}\). Wysokość ostrosłupa to 
 \(4\, \mathrm{cm}\). Punkt 
\(M\) to środek boku
 \(CV \). Oblicz odległość pomiędzy punktem
 \(M\) a płaszczyzną 
 \(ABC\).}
{\(2\, \mathrm{cm}\)}{\(\frac{\sqrt{34}}
 {2} \, \mathrm{cm}\)}{\(\frac{5}
{2}\, \mathrm{cm}\)}{}{}{}}
\LANGcs{
\Question{
Bod \(M\) je středem hrany
\(CV \) pravidelného čtyřbokého
jehlanu \(ABCDV \) s hlavním
vrcholem \(V \). Podstavná
hrana jehlanu má velikost \(6\, \mathrm{cm}\) 
a výška jehlanu je \(4\, \mathrm{cm}\).
Určete vzdálenost bodu \(M\) 
a roviny \(ABC\).}
{\(2\, \mathrm{cm}\)}{\(\frac{\sqrt{34}}
 {2} \, \mathrm{cm}\)}{\(\frac{5}
{2}\, \mathrm{cm}\)}{}{}{}}
\LANGsk{
\Question{
Bod \(M\) je stredom hrany
\(CV \) pravidelného štvorbokého
ihlanu \(ABCDV \) s hlavným
vrcholom \(V \). 
Hrana podstavy ihlanu má veľkosť \(6\, \mathrm{cm}\) 
a výška ihlanu je \(4\, \mathrm{cm}\).
Určte vzdialenosť bodu \(M\) 
a roviny \(ABC\).}
{\(2\, \mathrm{cm}\)}{\(\frac{\sqrt{34}}
 {2} \, \mathrm{cm}\)}{\(\frac{5}
{2}\, \mathrm{cm}\)}{}{}{}}
\LANGes{
\Question{
La base \(ABCD\) de una pirámide cuadrangular \(ABCDV \) tiene su arista de 
\(6\, \mathrm{cm}\). La altura de la pirámide es \(4\, \mathrm{cm}\). El punto
\(M\) es el centro de la arista \(CV \). Determina la distancia entre el punto \(M\) 
y el plano \(ABC\).}
{\(2\, \mathrm{cm}\)}{\(\frac{\sqrt{34}}
 {2} \, \mathrm{cm}\)}{\(\frac{5}
{2}\, \mathrm{cm}\)}{}{}{}}
\End

\Data\ID{9000128802}
\Updated{2020-07-15 03:07:37}
\Level{B}
\SubArea{Lines and planes: distances and angles}
\IMG{001288_02-figure0.pdf}
\LANGen{
\Question{
The base \(ABCD\) of a square
pyramid \(ABCDV \) has side
\(6\, \mathrm{cm}\). The height of the
pyramid is \(4\, \mathrm{cm}\). The point
\(M\) is the middle of the
side \(CV \). Find the distance
between the point \(M\) 
and the line \(BC\).}
{\(\frac{5}
{2}\, \mathrm{cm}\)}{\(\frac{\sqrt{34}}
 {2} \, \mathrm{cm}\)}{\(\frac{\sqrt{7}}
 {2} \, \mathrm{cm}\)}{}{}{}}
\LANGpl{
\Question{
Bok podstawy \(ABCD\) ostrosłupa prawidłowego czworokątnego 
 \(ABCDV \) jest równy
\(6\, \mathrm{cm}\). Wysokość ostrosłupa to 
 \(4\, \mathrm{cm}\). Punkt
\(M\) jest środkiem boku
 \(CV \). Oblicz odległość 
pomiędzy punktem  \(M\) 
a prostą  \(BC\).}
{\(\frac{5}
{2}\, \mathrm{cm}\)}{\(\frac{\sqrt{34}}
 {2} \, \mathrm{cm}\)}{\(\frac{\sqrt{7}}
 {2} \, \mathrm{cm}\)}{}{}{}}
\LANGcs{
\Question{
Bod \(M\) je středem hrany
\(CV \) pravidelného čtyřbokého
jehlanu \(ABCDV \) s hlavním
vrcholem \(V \). Podstavná
hrana jehlanu má velikost \(6\, \mathrm{cm}\) 
a výška jehlanu je \(4\, \mathrm{cm}\).
Určete vzdálenost bodu \(M\) 
a přímky \(BC\).}
{\(\frac{5}
{2}\, \mathrm{cm}\)}{\(\frac{\sqrt{34}}
 {2} \, \mathrm{cm}\)}{\(\frac{\sqrt{7}}
 {2} \, \mathrm{cm}\)}{}{}{}}
\LANGsk{
\Question{
Bod \(M\) je stredom hrany
\(CV \) pravidelného štvorbokého
ihlanu \(ABCDV \) s hlavným
vrcholom \(V \). 
Hrana podstavy ihlanu má veľkosť \(6\, \mathrm{cm}\) 
a výška ihlanu je \(4\, \mathrm{cm}\).
Určte vzdialenosť bodu \(M\) 
a priamky \(BC\).}
{\(\frac{5}
{2}\, \mathrm{cm}\)}{\(\frac{\sqrt{34}}
 {2} \, \mathrm{cm}\)}{\(\frac{\sqrt{7}}
 {2} \, \mathrm{cm}\)}{}{}{}}
\LANGes{
\Question{
La  base \(ABCD\) de una pirámide cuadrangular \(ABCDV \) tiene su arista de \(6\, \mathrm{cm}\). La altura de la pirámide es  \(4\, \mathrm{cm}\). El punto
\(M\) es el centro de la arista \(CV \). Determina la distancia entre el punto \(M\) 
y la recta \(BC\).}
{\(\frac{5}
{2}\, \mathrm{cm}\)}{\(\frac{\sqrt{34}}
 {2} \, \mathrm{cm}\)}{\(\frac{\sqrt{7}}
 {2} \, \mathrm{cm}\)}{}{}{}}
\End

\Data\ID{9000128803}
\Updated{2020-07-15 03:07:37}
\Level{B}
\SubArea{Lines and planes: distances and angles}
\IMG{001288_03-figure0.pdf}
\LANGen{
\Question{
The base \(ABCD\) of a square
pyramid \(ABCDV \) has side
\(6\, \mathrm{cm}\). The height of the
pyramid is \(4\, \mathrm{cm}\). The point
\(M\) is the middle of the
side \(CV \). Find the distance
between the point \(M\) 
and the line \(AD\).}
{\(\frac{\sqrt{97}}
 {2} \, \mathrm{cm}\)}{\(\frac{\sqrt{106}}
 {2} \, \mathrm{cm}\)}{\(\frac{\sqrt{65}}
 {2} \, \mathrm{cm}\)}{}{}{}}
\LANGpl{
\Question{
Bok podstawy  \(ABCD\) ostrosłupa prawidłowego czworokątnego 
 \(ABCDV \) wynosi
\(6\, \mathrm{cm}\). Wysokość ostrosłupa jest równa 
 \(4\, \mathrm{cm}\). Punkt
\(M\) to środek boku
 \(CV \). Oblicz odległość między punktem
 \(M\) 
a prostą \(AD\).}
{\(\frac{\sqrt{97}}
 {2} \, \mathrm{cm}\)}{\(\frac{\sqrt{106}}
 {2} \, \mathrm{cm}\)}{\(\frac{\sqrt{65}}
 {2} \, \mathrm{cm}\)}{}{}{}}
\LANGcs{
\Question{
Bod \(M\) je středem hrany
\(CV \) pravidelného čtyřbokého
jehlanu \(ABCDV \) s hlavním
vrcholem \(V \). Podstavná
hrana jehlanu má velikost \(6\, \mathrm{cm}\) 
a výška jehlanu je \(4\, \mathrm{cm}\).
Určete vzdálenost bodu \(M\) 
a přímky \(AD\).}
{\(\frac{\sqrt{97}}
 {2} \, \mathrm{cm}\)}{\(\frac{\sqrt{106}}
 {2} \, \mathrm{cm}\)}{\(\frac{\sqrt{65}}
 {2} \, \mathrm{cm}\)}{}{}{}}
\LANGsk{
\Question{
Bod \(M\) je stredom hrany
\(CV \) pravidelného štvorbokého
ihlanu \(ABCDV \) s hlavným
vrcholom \(V \).
Hrana podstavy ihlanu má veľkosť \(6\, \mathrm{cm}\) 
a výška ihlanu je \(4\, \mathrm{cm}\).
Určte vzdialenosť bodu \(M\) 
a priamky \(AD\).}
{\(\frac{\sqrt{97}}
 {2} \, \mathrm{cm}\)}{\(\frac{\sqrt{106}}
 {2} \, \mathrm{cm}\)}{\(\frac{\sqrt{65}}
 {2} \, \mathrm{cm}\)}{}{}{}}
\LANGes{
\Question{
La base \(ABCD\) de una pirámide cuadrangular \(ABCDV \)  tiene su arista de
\(6\, \mathrm{cm}\).  La altura de la pirámide es \(4\, \mathrm{cm}\). El punto
\(M\) es el centro de la arista \(CV \). Determina la distancia entre el punto \(M\) 
y la recta \(AD\).}
{\(\frac{\sqrt{97}}
 {2} \, \mathrm{cm}\)}{\(\frac{\sqrt{106}}
 {2} \, \mathrm{cm}\)}{\(\frac{\sqrt{65}}
 {2} \, \mathrm{cm}\)}{}{}{}}
\End

\Data\ID{9000128804}
\Updated{2020-07-15 03:07:37}
\Level{B}
\SubArea{Lines and planes: distances and angles}
\IMG{001288_04-figure0.pdf}
\LANGen{
\Question{
The base \(ABCD\) of a square
pyramid \(ABCDV \) has side
\(6\, \mathrm{cm}\). The height of the
pyramid is \(4\, \mathrm{cm}\). Find the
distance between the line \(AD\) 
and the plane \(BCV \).}
{\(\frac{24}
 {5} \, \mathrm{cm}\)}{\(\frac{15\sqrt{34}}
 {5} \, \mathrm{cm}\)}{\(6\, \mathrm{cm}\)}{}{}{}}
\LANGpl{
\Question{
Bok podstawy  \(ABCD\) ostrosłupa prawidłowego czworokątnego 
 \(ABCDV \) jest równy
\(6\, \mathrm{cm}\). Wysokość ostrosłupa wynosi
 \(4\, \mathrm{cm}\). Oblicz odległość pomiędzy
prostą  \(AD\), a
płaszczyzną  \(BCV \).}
{\(\frac{24}
 {5} \, \mathrm{cm}\)}{\(\frac{15\sqrt{34}}
 {5} \, \mathrm{cm}\)}{\(6\, \mathrm{cm}\)}{}{}{}}
\LANGcs{
\Question{
V pravidelném čtyřbokém jehlanu
\(ABCDV \) s hlavním vrcholem
\(V \) má podstavná hrana
velikost \(6\, \mathrm{cm}\) a výška
jehlanu je \(4\, \mathrm{cm}\). Určete
vzdálenost přímky \(AD\) 
a roviny \(BCV \).}
{\(\frac{24}
 {5} \, \mathrm{cm}\)}{\(\frac{15\sqrt{34}}
 {5} \, \mathrm{cm}\)}{\(6\, \mathrm{cm}\)}{}{}{}}
\LANGsk{
\Question{
V pravidelnom štvorbokom ihlane
\(ABCDV \) s hlavným vrcholom
\(V \) má hrana podstavy
veľkosť \(6\, \mathrm{cm}\) a výška
ihlanu je \(4\, \mathrm{cm}\). Určte
vzdialenosť priamky \(AD\) 
a roviny \(BCV \).}
{\(\frac{24}
 {5} \, \mathrm{cm}\)}{\(\frac{15\sqrt{34}}
 {5} \, \mathrm{cm}\)}{\(6\, \mathrm{cm}\)}{}{}{}}
\LANGes{
\Question{
La base \(ABCD\) de una pirámide cuadrangular  \(ABCDV \) tiene su arista de
\(6\, \mathrm{cm}\).  La altura de la pirámide es \(4\, \mathrm{cm}\). Determina la distancia entre la recta \(AD\) 
y el plano \(BCV \).}
{\(\frac{24}
 {5} \, \mathrm{cm}\)}{\(\frac{15\sqrt{34}}
 {5} \, \mathrm{cm}\)}{\(6\, \mathrm{cm}\)}{}{}{}}
\End

\Data\ID{9000128805}
\Updated{2020-07-15 03:07:37}
\Level{B}
\SubArea{Lines and planes: distances and angles}
\IMG{001288_05-figure0.pdf}
\LANGen{
\Question{
The base \(ABCD\) of a square
pyramid \(ABCDV \) has side
\(6\, \mathrm{cm}\). The height of the
pyramid is \(4\, \mathrm{cm}\). Find the
angle between the line \(BV \) 
and the plane \(ABC\).
Round to two decimal places.}
{\(43.31^{\circ }\)}{\(59.04^{\circ }\)}{\(45^{\circ }\)}{}{}{}}
\LANGpl{
\Question{
Bok podstawy  \(ABCD\) ostrosłupa prawidłowego czworokątnego 
 \(ABCDV \) jest równy
\(6\, \mathrm{cm}\). Wysokość ostrosłupa wynosi
 \(4\, \mathrm{cm}\). Oblicz kąt między
prostą  \(BV \) a 
płaszczyzną  \(ABC\).
Zaokrągli wynik do dwóch miejsc po przecinku.}
{\(43.31^{\circ }\)}{\(59.04^{\circ }\)}{\(45^{\circ }\)}{}{}{}}
\LANGcs{
\Question{
V pravidelném čtyřbokém jehlanu
\(ABCDV \) s hlavním vrcholem
\(V \) má podstavná hrana
velikost \(6\, \mathrm{cm}\) a výška
jehlanu je \(4\, \mathrm{cm}\). Určete
odchylku přímky \(BV \) 
a roviny \(ABC\).
Výsledek zaokrouhlete na dvě desetinná místa.}
{\(43{,}31^{\circ }\)}{\(59{,}04^{\circ }\)}{\(45^{\circ }\)}{}{}{}}
\LANGsk{
\Question{
V pravidelnom štvorbokom ihlane
\(ABCDV \) s hlavným vrcholom
\(V \) má hrana podstavy
veľkosť \(6\, \mathrm{cm}\) a výška
ihlanu je \(4\, \mathrm{cm}\). Určte
odchýlku priamky \(BV \) 
a roviny \(ABC\).
Výsledok zaokrúhlite na dve desatinné miesta.}
{\(43{,}31^{\circ }\)}{\(59{,}04^{\circ }\)}{\(45^{\circ }\)}{}{}{}}
\LANGes{
\Question{
La base \(ABCD\) de una pirámide cuadrangular  \(ABCDV \)  tiene su arista de
\(6\, \mathrm{cm}\).  La altura de la pirámide es \(4\, \mathrm{cm}\). Determina el ángulo entre la recta \(BV \) 
y el plano \(ABC\).
Redondea a dos posiciones decimales.}
{\(43.31^{\circ }\)}{\(59.04^{\circ }\)}{\(45^{\circ }\)}{}{}{}}
\End

\Data\ID{9000128806}
\Updated{2020-07-15 03:07:37}
\Level{B}
\SubArea{Lines and planes: distances and angles}
\IMG{001288_06-figure0.pdf}
\LANGen{
\Question{
The base \(ABCD\) of a square
pyramid \(ABCDV \) has side
\(6\, \mathrm{cm}\). The height of the
pyramid is \(4\, \mathrm{cm}\). The point
\(M\) is the middle of the
side \(CV \). Find the angle
between the line \(AM\) 
and the plane \(ABC\).
Round to two decimal places.}
{\(17.45^{\circ }\)}{\(34.50^{\circ }\)}{\(18.32^{\circ }\)}{}{}{}}
\LANGpl{
\Question{
Bok podstawy  \(ABCD\)  ostrosłupa prawidłowego czworokątnego 
 \(ABCDV \) jest równy
\(6\, \mathrm{cm}\). Wysokość ostrosłupa 
wynosi \(4\, \mathrm{cm}\). Punkt 
\(M\) to środek boku
 \(CV \). Wskaż kąt między 
prostą  \(AM\) 
a płaszczyzną  \(ABC\).
Zaokrągli wynik do dwóch miejsc po przecinku.}
{\(17.45^{\circ }\)}{\(34.50^{\circ }\)}{\(18.32^{\circ }\)}{}{}{}}
\LANGcs{
\Question{
Bod \(M\) je středem hrany
\(CV \) pravidelného čtyřbokého
jehlanu \(ABCDV \) s hlavním
vrcholem \(V \). Podstavná
hrana jehlanu má velikost \(6\, \mathrm{cm}\) 
a výška jehlanu je \(4\, \mathrm{cm}\).
Určete odchylku přímky \(AM\) 
a roviny \(ABC\).
Výsledek zaokrouhlete na dvě desetinná místa.}
{\(17{,}45^{\circ }\)}{\(34{,}50^{\circ }\)}{\(18{,}32^{\circ }\)}{}{}{}}
\LANGsk{
\Question{
Bod \(M\) je stredom hrany
\(CV \) pravidelného štvorbokého
ihlanu \(ABCDV \) s hlavným
vrcholom \(V \). Hrana podstavy ihlanu má veľkosť \(6\, \mathrm{cm}\) 
a výška ihlanu je \(4\, \mathrm{cm}\).
Určte odchýlku priamky \(AM\) 
a roviny \(ABC\).
Výsledok zaokrúhlite na dve desatinné miesta.}
{\(17{,}45^{\circ }\)}{\(34{,}50^{\circ }\)}{\(18{,}32^{\circ }\)}{}{}{}}
\LANGes{
\Question{
La base \(ABCD\) de una pirámide cuadrangular \(ABCDV \)  tiene su arista de
\(6\, \mathrm{cm}\).  La altura de la pirámide es \(4\, \mathrm{cm}\). El punto
\(M\) es el centro de la arista \(CV \). Determina el ángulo entre la recta \(AM\) 
y el plano \(ABC\).
Redondea a dos posiciones decimales.}
{\(17.45^{\circ }\)}{\(34.50^{\circ }\)}{\(18.32^{\circ }\)}{}{}{}}
\End

\Data\ID{9000128807}
\Updated{2020-07-15 03:07:37}
\Level{B}
\SubArea{Lines and planes: distances and angles}
\IMG{001288_07-figure0.pdf}
\LANGen{
\Question{
The base \(ABCD\) of a square
pyramid \(ABCDV \) has side
\(6\, \mathrm{cm}\). The height of the
pyramid is \(4\, \mathrm{cm}\). Find the
angle between the planes \(DCV \) 
and \(ABC\).
Round to two decimal places.}
{\(53.13^{\circ }\)}{\(59.04^{\circ }\)}{\(43.31^{\circ }\)}{}{}{}}
\LANGpl{
\Question{
Bok podstawy  \(ABCD\)  ostrosłupa prawidłowego czworokątnego 
 \(ABCDV \) jest równy
\(6\, \mathrm{cm}\). Wysokość ostrosłupa wynosi
 \(4\, \mathrm{cm}\). Oblicz kąt 
między płaszczyznami  \(DCV \) 
i  \(ABC\).
Zaokrągli wynik do dwóch miejsc po przecinku.}
{\(53.13^{\circ }\)}{\(59.04^{\circ }\)}{\(43.31^{\circ }\)}{}{}{}}
\LANGcs{
\Question{
V pravidelném čtyřbokém jehlanu
\(ABCDV \) s hlavním vrcholem
\(V \) má podstavná
hrana velikost \(6\, \mathrm{cm}\) a
výška jehlanu je \(4\, \mathrm{cm}\).
Určete odchylku rovin \(DCV \) 
a \(ABC\).
Výsledek zaokrouhlete na dvě desetinná místa.}
{\(53{,}13^{\circ }\)}{\(59{,}04^{\circ }\)}{\(43{,}31^{\circ }\)}{}{}{}}
\LANGsk{
\Question{
V pravidelnom štvorbokom ihlane
\(ABCDV \) s hlavným vrcholom
\(V \) má hrana podstavy veľkosť \(6\, \mathrm{cm}\) a
výška ihlanu je \(4\, \mathrm{cm}\).
Určte odchýlku rovín \(DCV \) 
a \(ABC\).
Výsledok zaokrúhlite na dve desatinné miesta.}
{\(53{,}13^{\circ }\)}{\(59{,}04^{\circ }\)}{\(43{,}31^{\circ }\)}{}{}{}}
\LANGes{
\Question{
La base \(ABCD\) de una pirámide cuadrangular  \(ABCDV \)  tiene su arista de
\(6\, \mathrm{cm}\).  La altura de la pirámide es \(4\, \mathrm{cm}\).  Determina el ángulo entre los planos \(DCV \) 
y \(ABC\).
Redondea a dos posiciones decimales.}
{\(53.13^{\circ }\)}{\(59.04^{\circ }\)}{\(43.31^{\circ }\)}{}{}{}}
\End

\Data\ID{9000128808}
\Updated{2020-07-15 03:07:37}
\Level{B}
\SubArea{Lines and planes: distances and angles}
\IMG{001288_08-figure0.pdf}
\LANGen{
\Question{
The base \(ABCD\) of a square
pyramid \(ABCDV \) has side
\(6\, \mathrm{cm}\). The height of the
pyramid is \(4\, \mathrm{cm}\). Find the
angle between the planes \(ADV \) 
and \(BCV \).
Round to two decimal places.}
{\(73.74^{\circ }\)}{\(36.87^{\circ }\)}{\(61.93^{\circ }\)}{}{}{}}
\LANGpl{
\Question{
Bok podstawy  \(ABCD\) ostrosłupa prawidłowego czworokątnego 
\(ABCDV \) jest  równy
\(6\, \mathrm{cm}\). Wysokość ostrosłupa wynosi
 \(4\, \mathrm{cm}\). Oblicz kąt między płaszczyznami 
 \(ADV \) 
i \(BCV \).
Zaokrągli wynik do dwóch miejsc po przecinku.}
{\(73.74^{\circ }\)}{\(36.87^{\circ }\)}{\(61.93^{\circ }\)}{}{}{}}
\LANGcs{
\Question{
V pravidelném čtyřbokém jehlanu
\(ABCDV \) s hlavním vrcholem
\(V \) má podstavná
hrana velikost \(6\, \mathrm{cm}\) a
výška jehlanu je \(4\, \mathrm{cm}\).
Určete odchylku rovin \(ADV \) 
a \(BCV \).
Výsledek zaokrouhlete na dvě desetinná místa.}
{\(73{,}74^{\circ }\)}{\(36{,}87^{\circ }\)}{\(61{,}93^{\circ }\)}{}{}{}}
\LANGsk{
\Question{
V pravidelnom štvorbokom ihlane
\(ABCDV \) s hlavným vrcholom
\(V \) má hrana podstavy veľkosť \(6\, \mathrm{cm}\) a
výška ihlanu je \(4\, \mathrm{cm}\).
Určte odchýlku rovín \(ADV \) 
a \(BCV \).
Výsledok zaokrúhlite na dve desatinné miesta.}
{\(73{,}74^{\circ }\)}{\(36{,}87^{\circ }\)}{\(61{,}93^{\circ }\)}{}{}{}}
\LANGes{
\Question{
La base \(ABCD\) de una pirámide cuadrangular  \(ABCDV \)  tiene su arista de
\(6\, \mathrm{cm}\).  La altura de la pirámide es \(4\, \mathrm{cm}\). Determina el ángulo entre los planos \(ADV \) 
y \(BCV \).
Redondea a dos posiciones decimales.}
{\(73.74^{\circ }\)}{\(36.87^{\circ }\)}{\(61.93^{\circ }\)}{}{}{}}
\End

\Data\ID{9000153601}
\Updated{2020-07-15 03:07:37}
\Level{B}
\SubArea{Lines and planes: distances and angles}
\IMG{001536_01-figure0.pdf}
\LANGen{
\Question{
Give a verbal description to the angle shown in the picture.}
{The angle between the triangular face and the base.}{The angle between the edge and the base.}{The angle between two triangular faces.}{The angle between the edge of a triangular face and the base edge.}{}{}}
\LANGpl{
\Question{
Wybierz prawidłowy opis kąta pokazanego na rysunku.}
{Kąt między ścianą boczną a podstawą.}{Kąt między krawędzią  a podstawą.}{Kąt między ścianami bocznymi.}{Kąt między krawędzią ściany bocznej a krawędzią podstawy.}{}{}}
\LANGcs{
\Question{
Úhel vyznačený na obrázku znázorňuje:}
{Odchylku boční stěny a podstavy.}{Odchylku boční hrany a podstavy.}{Odchylku dvou sousedních bočních stěn.}{Odchylku boční hrany a podstavné hrany.}{}{}}
\LANGsk{
\Question{
Uhol vyznačený na obrázku znázorňuje:}
{Odchýlku bočnej steny a podstavy.}{Odchýlku bočnej hrany a podstavy.}{Odchýlku dvoch susedných bočných stien.}{Odchýlku bočnej hrany a hrany podstavy.}{}{}}
\LANGes{
\Question{
El ángulo marcado en la figura muestra:}
{El ángulo entre la cara triangular y la base.}{El ángulo entre la arista y la base.}{El ángulo entre dos  caras triangulares.}{El ángulo entre la arista de una cara triangular y la arista básica.}{}{}}
\End

\Data\ID{9000153602}
\Updated{2020-07-15 03:07:37}
\Level{B}
\SubArea{Lines and planes: distances and angles}
\IMG{001536_02-figure0.pdf}
\LANGen{
\Question{
Give a verbal description to the angle shown in the picture.}
{The angle between the edge and the base.}{The angle between the triangular face and the square base.}{The angle between two opposite edges.}{The angle between the edge on a face and the base edge.}{}{}}
\LANGpl{
\Question{
Wybierz prawidłowy opis kąta pokazanego na rysunku.}
{Kąt między krawędzią ściany bocznej a  podstawą.}{Kąt miedzy ścianą boczną  a podstawą.}{Kąt pomiędzy przeciwległymi krawędziami.}{Kąt między krawędzią ściany a krawędzią podstawy.}{}{}}
\LANGcs{
\Question{
Úhel vyznačený na obrázku znázorňuje:}
{Odchylku boční hrany a podstavy.}{Odchylku boční stěny a podstavy.}{Odchylku dvou protilehlých hran.}{Odchylku podstavné hrany a boční hrany.}{}{}}
\LANGsk{
\Question{
Uhol vyznačený na obrázku znázorňuje:}
{Odchýlku bočnej hrany a podstavy.}{Odchýlku bočnej steny a podstavy.}{Odchýlku dvoch protiľahlých hrán.}{Odchýlku hrany podstavy a bočnej hrany.}{}{}}
\LANGes{
\Question{
El ángulo marcado en la figura muestra:}
{El ángulo entre la arista y la base.}{El ángulo entre la cara triangular y la base cuadrangular.}{El ángulo entre dos arista opuestas.}{El ángulo entre la arista de una cara y la arista básica.}{}{}}
\End

\Data\ID{9000153603}
\Updated{2020-07-15 03:07:37}
\Level{B}
\SubArea{Lines and planes: distances and angles}
\IMG{001536_03-figure0.pdf}
\LANGen{
\Question{
Give a verbal description to the angle shown in the picture.}
{The angle between two opposite triangular faces.}{The angle between a triangular face and the base.}{The angle between two edges in the same triangular face.}{The angle between two triangular faces having a common edge.}{}{}}
\LANGpl{
\Question{
Wybierz prawidłowy opis kąta pokazanego na rysunku.}
{Kąt miedzy dwoma przeciwległymi ścianami bocznymi.}{Kat między ścianą boczną a podstawą.}{Kąt między dwoma krawędziami ściany bocznej.}{Kąt między dwoma ścianami bocznymi o wspólnej krawędzi.}{}{}}
\LANGcs{
\Question{
Úhel vyznačený na obrázku znázorňuje:}
{Odchylku dvou protilehlých bočních stěn.}{Odchylku boční stěny a podstavy.}{Odchylku dvou sousedních bočních hran.}{Odchylku dvou sousedních bočních stěn.}{}{}}
\LANGsk{
\Question{
Uhol vyznačený na obrázku znázorňuje:}
{Odchýlku dvoch protiľahlých bočných stien.}{Odchýlku bočnej steny a podstavy.}{Odchýlku dvoch susedných bočných hrán.}{Odchýlku dvoch susedených bočných stien.}{}{}}
\LANGes{
\Question{
El ángulo marcado en la figura muestra:}
{El ángulo entre dos caras triangulares opuestas.}{El ángulo entre la cara triangular y la base.}{El ángulo entre dos aristas de la misma cara triangular.}{El ángulo entre dos caras triangulares que tienen una arista común.}{}{}}
\End

\Data\ID{9000153604}
\Updated{2020-07-15 03:07:37}
\Level{B}
\SubArea{Lines and planes: distances and angles}
\IMG{001536_04-figure0.pdf}
\LANGen{
\Question{
Give a verbal description to the angle shown in the picture.}
{The angle between two edges in a common triangular face.}{The angle between two opposite triangular faces.}{The angle between two opposite edges.}{The angle between two triangular faces having common edge.}{}{}}
\LANGpl{
\Question{
Wybierz prawidłowy opis kąta pokazanego na rysunku.}
{Kąt między krawędziami  ściany bocznej.}{Kat między przeciwległymi ścianami bocznymi.}{Kąt między przeciwległymi krawędziami.}{Kąt między ścianami bocznymi mającymi wspólną krawędź.}{}{}}
\LANGcs{
\Question{
Úhel vyznačený na obrázku znázorňuje:}
{Odchylku dvou sousedních bočních hran.}{Odchylku dvou protilehlých bočních stěn.}{Odchylku dvou protilehlých bočních hran.}{Odchylku dvou sousedních bočních stěn.}{}{}}
\LANGsk{
\Question{
Uhol vyznačený na obrázku znázorňuje:}
{Odchýlku dvoch susedených bočných hrán.}{Odchýlku dvoch protiľahlých bočných stien.}{Odchýlku dvoch protiľahlých bočných hrán.}{Odchýlku dvoch susedených bočných stien.}{}{}}
\LANGes{
\Question{
El ángulo marcado en la figura muestra:}
{El ángulo entre dos aristas que tienen la cara triangular común.}{El ángulo entre dos caras triangulares opuestas.}{El ángulo entre dos aristas opuestas.}{El ángulo entre dos caras triangulares que tienen una arista común.}{}{}}
\End

\Data\ID{9000153605}
\Updated{2020-07-15 03:07:37}
\Level{B}
\SubArea{Lines and planes: distances and angles}
\IMG{001536_05-figure0.pdf}
\LANGen{
\Question{
Give a verbal description to the angle shown in the picture.}
{The angle between opposite edges.}{The angle between a triangular face and an edge from the opposite triangular face.}{The angle between two opposite triangular faces.}{The angle between two triangular faces having a common edge.}{}{}}
\LANGpl{
\Question{
Wybierz prawidłowy opis kąta pokazanego na rysunku.}
{Kąt między przeciwległymi krawędziami.}{Kąt między ścianą boczną a krawędzią przeciwległej ściany bocznej.}{Kąt między przeciwległymi ścianami bocznymi.}{Kąt między  ścianami bocznymi o wspólnej krawędzi.}{}{}}
\LANGcs{
\Question{
Úhel vyznačený na obrázku znázorňuje:}
{Odchylku dvou protilehlých bočních hran.}{Odchylku boční stěny a boční hrany.}{Odchylku dvou protilehlých bočních stěn.}{Odchylku dvou sousedních bočních stěn.}{}{}}
\LANGsk{
\Question{
Uhol vyznačený na obrázku znázorňuje:}
{Odchýlku dvoch protiľahlých bočných hrán.}{Odchýlku bočnej steny a bočnej hrany.}{Odchýlku dvoch protiľahlých bočných stien.}{Odchýlku dvoch susedených bočných stien.}{}{}}
\LANGes{
\Question{
El ángulo marcado en la figura muestra:}
{El ángulo entre dos aristas opuestas.}{El ángulo entre la cara triangular y la arista de la cara triangular opuesta.}{El ángulo entre dos caras triangulares opuestas.}{El ángulo entre dos caras triangulares que tienen una arista común.}{}{}}
\End

\Data\ID{9000153606}
\Updated{2020-07-15 03:07:37}
\Level{B}
\SubArea{Lines and planes: distances and angles}
\IMG{001536_06-figure0.pdf}
\LANGen{
\Question{
Give a verbal description to the angle shown in the picture.}
{The angle between the edge on triangular face and the base edge from the same face.}{The angle between the triangular face and a base edge not in this face.}{The angle between two triangular faces having a common edge.}{The angle between a triangular face and square base.}{}{}}
\LANGpl{
\Question{
Wybierz prawidłowy opis kąta pokazanego na rysunku.}
{Kąt pomiędzy krawędzią ściany bocznej a krawędzią podstawy znajdujących się na tej samej  płaszczyźnie.}{Kąt pomiędzy krawędzią ściany bocznej a krawędzią podstawy znajdujących się na różnych płaszczyznach.}{Kąt między  ścianami bocznymi o wspólnej krawędzi.}{Kąt między ścianą boczną a podstawą.}{}{}}
\LANGcs{
\Question{
Úhel vyznačený na obrázku znázorňuje:}
{Odchylku boční hrany a podstavné hrany.}{Odchylku boční stěny a podstavné hrany.}{Odchylku dvou sousedních bočních stěn.}{Odchylku boční stěny a podstavy.}{}{}}
\LANGsk{
\Question{
Uhol vyznačený na obrázku znázorňuje:}
{Odchýlku bočnej hrany a podstavné hrany.}{Odchýlku bočnej steny a hrany podstavy.}{Odchýlku dvoch susedených bočných stien.}{Odchýlku bočnej steny a podstavy.}{}{}}
\LANGes{
\Question{
El ángulo marcado en la figura muestra:}
{El ángulo entre la arista de la cara  triangular  y la arista básica de la misma cara.}{El ángulo entre la cara triangular y la arista básica que no se encuentra en la misma cara triangular.}{El ángulo entre dos caras triangulares que tienen una arista común.}{El ángulo entre la cara triangular y la base cuadrangular..}{}{}}
\End

\Data\ID{9000153701}
\Updated{2020-07-15 03:07:37}
\Level{B}
\SubArea{Lines and planes: distances and angles}
\IMG{001537_01-figure0.pdf}
\LANGen{
\Question{
The picture shows a square pyramid. The side of a base square is
\(a = 4\; \mathrm{cm}\) and the height of
the pyramid is \(v = 6\; \mathrm{cm}\).
Find the angle \(\varphi \).}
{\(\mathop{\mathrm{tg}}\nolimits \varphi = \frac{6} 
{2}\mathrel{\implies }\varphi \mathop{\mathop{\doteq }}\nolimits 71^{\circ }34^{\prime}\)}{\(\mathop{\mathrm{tg}}\nolimits \varphi = \frac{6} 
{2\sqrt{2}}\mathrel{\implies }\varphi \mathop{\mathop{\doteq }}\nolimits 64^{\circ }46^{\prime}\)}{\(\mathop{\mathrm{tg}}\nolimits \frac{\varphi }
{2} = \frac{2} 
{6}\mathrel{\implies }\varphi \mathop{\mathop{\doteq }}\nolimits 36^{\circ }52^{\prime}\)}{\(\mathop{\mathrm{tg}}\nolimits \frac{\varphi }
{2} = \frac{2} 
{2\sqrt{10}}\mathrel{\implies }\varphi \mathop{\mathop{\doteq }}\nolimits 35^{\circ }6^{\prime}\)}{}{}}
\LANGpl{
\Question{
Rysunek przedstawia ostrosłup prawidłowy czworokątny o boku 
\(a = 4\; \mathrm{cm}\) i wysokości
 \(v = 6\; \mathrm{cm}\).
Wyznacz kąt \(\varphi \).}
{\(\mathop{\mathrm{tg}}\nolimits \varphi = \frac{6} 
{2}\mathrel{\implies }\varphi \mathop{\mathop{\doteq }}\nolimits 71^{\circ }34^{\prime}\)}{\(\mathop{\mathrm{tg}}\nolimits \varphi = \frac{6} 
{2\sqrt{2}}\mathrel{\implies }\varphi \mathop{\mathop{\doteq }}\nolimits 64^{\circ }46^{\prime}\)}{\(\mathop{\mathrm{tg}}\nolimits \frac{\varphi }
{2} = \frac{2} 
{6}\mathrel{\implies }\varphi \mathop{\mathop{\doteq }}\nolimits 36^{\circ }52^{\prime}\)}{\(\mathop{\mathrm{tg}}\nolimits \frac{\varphi }
{2} = \frac{2} 
{2\sqrt{10}}\mathrel{\implies }\varphi \mathop{\mathop{\doteq }}\nolimits 35^{\circ }6^{\prime}\)}{}{}}
\LANGcs{
\Question{
Na obrázku je pravidelný čtyřboký jehlan se čtvercovou podstavou o hraně
\(a = 4\; \mathrm{cm}\) s tělesovou
výškou \(v = 6\; \mathrm{cm}\).
Pro velikost vyznačené odchylky platí:}
{\(\mathop{\mathrm{tg}}\nolimits \varphi = \frac{6} 
{2}\mathrel{\implies }\varphi \mathop{\mathop{\doteq }}\nolimits 71^{\circ }34^{\prime}\)}{\(\mathop{\mathrm{tg}}\nolimits \varphi = \frac{6} 
{2\sqrt{2}}\mathrel{\implies }\varphi \mathop{\mathop{\doteq }}\nolimits 64^{\circ }46^{\prime}\)}{\(\mathop{\mathrm{tg}}\nolimits \frac{\varphi }
{2} = \frac{2} 
{6}\mathrel{\implies }\varphi \mathop{\mathop{\doteq }}\nolimits 36^{\circ }52^{\prime}\)}{\(\mathop{\mathrm{tg}}\nolimits \frac{\varphi }
{2} = \frac{2} 
{2\sqrt{10}}\mathrel{\implies }\varphi \mathop{\mathop{\doteq }}\nolimits 35^{\circ }6^{\prime}\)}{}{}}
\LANGsk{
\Question{
Na obrázku je pravidelný štvorboký ihlan so štvorcovou podstavou s hranou
\(a = 4\; \mathrm{cm}\) s telesovou
výškou \(v = 6\; \mathrm{cm}\).
Pre veľkosť vyznačenej odchýlky platí:}
{\(\mathop{\mathrm{tg}}\nolimits \varphi = \frac{6} 
{2}\mathrel{\implies }\varphi \mathop{\mathop{\doteq }}\nolimits 71^{\circ }34^{\prime}\)}{\(\mathop{\mathrm{tg}}\nolimits \varphi = \frac{6} 
{2\sqrt{2}}\mathrel{\implies }\varphi \mathop{\mathop{\doteq }}\nolimits 64^{\circ }46^{\prime}\)}{\(\mathop{\mathrm{tg}}\nolimits \frac{\varphi }
{2} = \frac{2} 
{6}\mathrel{\implies }\varphi \mathop{\mathop{\doteq }}\nolimits 36^{\circ }52^{\prime}\)}{\(\mathop{\mathrm{tg}}\nolimits \frac{\varphi }
{2} = \frac{2} 
{2\sqrt{10}}\mathrel{\implies }\varphi \mathop{\mathop{\doteq }}\nolimits 35^{\circ }6^{\prime}\)}{}{}}
\LANGes{
\Question{
La imagen muestra una pirámide cuadrangular. La arista de una base cuadrada es
\(a = 4\; \mathrm{cm}\) y la altura de
la pirámide es \(v = 6\; \mathrm{cm}\).
Determina el ángulo \(\varphi \).}
{\(\mathop{\mathrm{tg}}\nolimits \varphi = \frac{6} 
{2}\mathrel{\implies }\varphi \mathop{\mathop{\doteq }}\nolimits 71^{\circ }34^{\prime}\)}{\(\mathop{\mathrm{tg}}\nolimits \varphi = \frac{6} 
{2\sqrt{2}}\mathrel{\implies }\varphi \mathop{\mathop{\doteq }}\nolimits 64^{\circ }46^{\prime}\)}{\(\mathop{\mathrm{tg}}\nolimits \frac{\varphi }
{2} = \frac{2} 
{6}\mathrel{\implies }\varphi \mathop{\mathop{\doteq }}\nolimits 36^{\circ }52^{\prime}\)}{\(\mathop{\mathrm{tg}}\nolimits \frac{\varphi }
{2} = \frac{2} 
{2\sqrt{10}}\mathrel{\implies }\varphi \mathop{\mathop{\doteq }}\nolimits 35^{\circ }6^{\prime}\)}{}{}}
\End

\Data\ID{9000153702}
\Updated{2020-07-15 03:07:12}
\Level{B}
\SubArea{Lines and planes: distances and angles}
\IMG{001537_02-figure0.pdf}
\LANGen{
\Question{
The picture shows a square pyramid. The side of a base square is
\(a = 4\; \mathrm{cm}\) and the height of
the pyramid is \(v = 6\; \mathrm{cm}\).
Find the angle \(\varphi \).}
{\(\mathop{\mathrm{tg}}\nolimits \varphi = \frac{6} 
{2\sqrt{2}}\mathrel{\implies }\varphi \mathop{\mathop{\doteq }}\nolimits 64^{\circ }46^{\prime}\)}{\(\mathop{\mathrm{tg}}\nolimits \varphi = \frac{6} 
{2}\mathrel{\implies }\varphi \mathop{\mathop{\doteq }}\nolimits 71^{\circ }34^{\prime}\)}{\(\mathop{\mathrm{tg}}\nolimits \frac{\varphi }
{2} = \frac{2} 
{6}\mathrel{\implies }\varphi \mathop{\mathop{\doteq }}\nolimits 36^{\circ }52^{\prime}\)}{\(\mathop{\mathrm{tg}}\nolimits \varphi = \frac{2\sqrt{10}} 
 {2} \mathrel{\implies }\varphi \mathop{\mathop{\doteq }}\nolimits 72^{\circ }27^{\prime}\)}{}{}}
\LANGpl{
\Question{
Rysunek przedstawia ostrosłup prawidłowy czworokątny o boku podstawy  równym
\(a = 4\; \mathrm{cm}\) i wysokości 
 \(v = 6\; \mathrm{cm}\).
Wyznacz kąt \(\varphi \).}
{\(\mathop{\mathrm{tg}}\nolimits \varphi = \frac{6} 
{2\sqrt{2}}\mathrel{\implies }\varphi \mathop{\mathop{\doteq }}\nolimits 64^{\circ }46^{\prime}\)}{\(\mathop{\mathrm{tg}}\nolimits \varphi = \frac{6} 
{2}\mathrel{\implies }\varphi \mathop{\mathop{\doteq }}\nolimits 71^{\circ }34^{\prime}\)}{\(\mathop{\mathrm{tg}}\nolimits \frac{\varphi }
{2} = \frac{2} 
{6}\mathrel{\implies }\varphi \mathop{\mathop{\doteq }}\nolimits 36^{\circ }52^{\prime}\)}{\(\mathop{\mathrm{tg}}\nolimits \varphi = \frac{2\sqrt{10}} 
 {2} \mathrel{\implies }\varphi \mathop{\mathop{\doteq }}\nolimits 72^{\circ }27^{\prime}\)}{}{}}
\LANGcs{
\Question{
Na obrázku je pravidelný čtyřboký jehlan se čtvercovou podstavou o hraně
\(a = 4\; \mathrm{cm}\) s tělesovou
výškou \(v = 6\; \mathrm{cm}\).
Pro velikost vyznačené odchylky platí:}
{\(\mathop{\mathrm{tg}}\nolimits \varphi = \frac{6} 
{2\sqrt{2}}\mathrel{\implies }\varphi \mathop{\mathop{\doteq }}\nolimits 64^{\circ }46^{\prime}\)}{\(\mathop{\mathrm{tg}}\nolimits \varphi = \frac{6} 
{2}\mathrel{\implies }\varphi \mathop{\mathop{\doteq }}\nolimits 71^{\circ }34^{\prime}\)}{\(\mathop{\mathrm{tg}}\nolimits \frac{\varphi }
{2} = \frac{2} 
{6}\mathrel{\implies }\varphi \mathop{\mathop{\doteq }}\nolimits 36^{\circ }52^{\prime}\)}{\(\mathop{\mathrm{tg}}\nolimits \varphi = \frac{2\sqrt{10}} 
 {2} \mathrel{\implies }\varphi \mathop{\mathop{\doteq }}\nolimits 72^{\circ }27^{\prime}\)}{}{}}
\LANGsk{
\Question{
Na obrázku je pravidelný štvorboký ihlan so štvorcovou podstavou s hranou
\(a = 4\; \mathrm{cm}\) s telesovou
výškou \(v = 6\; \mathrm{cm}\).
Pre veľkosť vyznačenej odchýlky platí:}
{\(\mathop{\mathrm{tg}}\nolimits \varphi = \frac{6} 
{2\sqrt{2}}\mathrel{\implies }\varphi \mathop{\mathop{\doteq }}\nolimits 64^{\circ }46^{\prime}\)}{\(\mathop{\mathrm{tg}}\nolimits \varphi = \frac{6} 
{2}\mathrel{\implies }\varphi \mathop{\mathop{\doteq }}\nolimits 71^{\circ }34^{\prime}\)}{\(\mathop{\mathrm{tg}}\nolimits \frac{\varphi }
{2} = \frac{2} 
{6}\mathrel{\implies }\varphi \mathop{\mathop{\doteq }}\nolimits 36^{\circ }52^{\prime}\)}{\(\mathop{\mathrm{tg}}\nolimits \varphi = \frac{2\sqrt{10}} 
 {2} \mathrel{\implies }\varphi \mathop{\mathop{\doteq }}\nolimits 72^{\circ }27^{\prime}\)}{}{}}
\LANGes{
\Question{
La imagen muestra una pirámide cuadrangular. La arista de una base cuadrada es
\(a = 4\; \mathrm{cm}\) y la altura de
la pirámide es \(v = 6\; \mathrm{cm}\).
Determina el ángulo \(\varphi \).}
{\(\mathop{\mathrm{tg}}\nolimits \varphi = \frac{6} 
{2\sqrt{2}}\mathrel{\implies }\varphi \mathop{\mathop{\doteq }}\nolimits 64^{\circ }46^{\prime}\)}{\(\mathop{\mathrm{tg}}\nolimits \varphi = \frac{6} 
{2}\mathrel{\implies }\varphi \mathop{\mathop{\doteq }}\nolimits 71^{\circ }34^{\prime}\)}{\(\mathop{\mathrm{tg}}\nolimits \frac{\varphi }
{2} = \frac{2} 
{6}\mathrel{\implies }\varphi \mathop{\mathop{\doteq }}\nolimits 36^{\circ }52^{\prime}\)}{\(\mathop{\mathrm{tg}}\nolimits \varphi = \frac{2\sqrt{10}} 
 {2} \mathrel{\implies }\varphi \mathop{\mathop{\doteq }}\nolimits 72^{\circ }27^{\prime}\)}{}{}}
\End

\Data\ID{9000153703}
\Updated{2020-07-15 03:07:12}
\Level{B}
\SubArea{Lines and planes: distances and angles}
\IMG{001537_03-figure0.pdf}
\LANGen{
\Question{
The picture shows a square pyramid. The side of a base square is
\(a = 4\; \mathrm{cm}\) and the height of
the pyramid is \(v = 6\; \mathrm{cm}\).
Find the angle \(\varphi \).}
{\(\mathop{\mathrm{tg}}\nolimits \frac{\varphi }
{2} = \frac{2} 
{6}\mathrel{\implies }\varphi \mathop{\mathop{\doteq }}\nolimits 36^{\circ }52^{\prime}\)}{\(\mathop{\mathrm{tg}}\nolimits \varphi = \frac{6} 
{2\sqrt{2}}\mathrel{\implies }\varphi \mathop{\mathop{\doteq }}\nolimits 64^{\circ }46^{\prime}\)}{\(\mathop{\mathrm{tg}}\nolimits \frac{\varphi }
{2} = \frac{2} 
{2\sqrt{10}}\mathrel{\implies }\varphi \mathop{\mathop{\doteq }}\nolimits 35^{\circ }6^{\prime}\)}{\(\mathop{\mathrm{tg}}\nolimits \frac{\varphi }
{2} = \frac{2\sqrt{2}} 
 {6} \mathrel{\implies }\varphi \mathop{\mathop{\doteq }}\nolimits 50^{\circ }29^{\prime}\)}{}{}}
\LANGpl{
\Question{
Rysunek przedstawia ostrosłup prawidłowy czworokątny o boku podstawy  równym
\(a = 4\; \mathrm{cm}\) i wysokości
 \(v = 6\; \mathrm{cm}\).
Wyznacz kąt \(\varphi \).}
{\(\mathop{\mathrm{tg}}\nolimits \frac{\varphi }
{2} = \frac{2} 
{6}\mathrel{\implies }\varphi \mathop{\mathop{\doteq }}\nolimits 36^{\circ }52^{\prime}\)}{\(\mathop{\mathrm{tg}}\nolimits \varphi = \frac{6} 
{2\sqrt{2}}\mathrel{\implies }\varphi \mathop{\mathop{\doteq }}\nolimits 64^{\circ }46^{\prime}\)}{\(\mathop{\mathrm{tg}}\nolimits \frac{\varphi }
{2} = \frac{2} 
{2\sqrt{10}}\mathrel{\implies }\varphi \mathop{\mathop{\doteq }}\nolimits 35^{\circ }6^{\prime}\)}{\(\mathop{\mathrm{tg}}\nolimits \frac{\varphi }
{2} = \frac{2\sqrt{2}} 
 {6} \mathrel{\implies }\varphi \mathop{\mathop{\doteq }}\nolimits 50^{\circ }29^{\prime}\)}{}{}}
\LANGcs{
\Question{
Na obrázku je pravidelný čtyřboký jehlan se čtvercovou podstavou o hraně
\(a = 4\; \mathrm{cm}\) s tělesovou
výškou \(v = 6\; \mathrm{cm}\).
Pro velikost vyznačené odchylky platí:}
{\(\mathop{\mathrm{tg}}\nolimits \frac{\varphi }
{2} = \frac{2} 
{6}\mathrel{\implies }\varphi \mathop{\mathop{\doteq }}\nolimits 36^{\circ }52^{\prime}\)}{\(\mathop{\mathrm{tg}}\nolimits \varphi = \frac{6} 
{2\sqrt{2}}\mathrel{\implies }\varphi \mathop{\mathop{\doteq }}\nolimits 64^{\circ }46^{\prime}\)}{\(\mathop{\mathrm{tg}}\nolimits \frac{\varphi }
{2} = \frac{2} 
{2\sqrt{10}}\mathrel{\implies }\varphi \mathop{\mathop{\doteq }}\nolimits 35^{\circ }6^{\prime}\)}{\(\mathop{\mathrm{tg}}\nolimits \frac{\varphi }
{2} = \frac{2\sqrt{2}} 
 {6} \mathrel{\implies }\varphi \mathop{\mathop{\doteq }}\nolimits 50^{\circ }29^{\prime}\)}{}{}}
\LANGsk{
\Question{
Na obrázku je pravidelný štvorboký ihlan so štvorcovou podstavou s hranou
\(a = 4\; \mathrm{cm}\) s telesovou
výškou \(v = 6\; \mathrm{cm}\).
Pre veľkosť vyznačenej odchýlky platí:}
{\(\mathop{\mathrm{tg}}\nolimits \frac{\varphi }
{2} = \frac{2} 
{6}\mathrel{\implies }\varphi \mathop{\mathop{\doteq }}\nolimits 36^{\circ }52^{\prime}\)}{\(\mathop{\mathrm{tg}}\nolimits \varphi = \frac{6} 
{2\sqrt{2}}\mathrel{\implies }\varphi \mathop{\mathop{\doteq }}\nolimits 64^{\circ }46^{\prime}\)}{\(\mathop{\mathrm{tg}}\nolimits \frac{\varphi }
{2} = \frac{2} 
{2\sqrt{10}}\mathrel{\implies }\varphi \mathop{\mathop{\doteq }}\nolimits 35^{\circ }6^{\prime}\)}{\(\mathop{\mathrm{tg}}\nolimits \frac{\varphi }
{2} = \frac{2\sqrt{2}} 
 {6} \mathrel{\implies }\varphi \mathop{\mathop{\doteq }}\nolimits 50^{\circ }29^{\prime}\)}{}{}}
\LANGes{
\Question{
La imagen muestra una pirámide cuadrangular. La arista de una base cuadrada es
\(a = 4\; \mathrm{cm}\) y la altura de la pirámide es \(v = 6\; \mathrm{cm}\).
Determina el ángulo \(\varphi \).}
{\(\mathop{\mathrm{tg}}\nolimits \frac{\varphi }
{2} = \frac{2} 
{6}\mathrel{\implies }\varphi \mathop{\mathop{\doteq }}\nolimits 36^{\circ }52^{\prime}\)}{\(\mathop{\mathrm{tg}}\nolimits \varphi = \frac{6} 
{2\sqrt{2}}\mathrel{\implies }\varphi \mathop{\mathop{\doteq }}\nolimits 64^{\circ }46^{\prime}\)}{\(\mathop{\mathrm{tg}}\nolimits \frac{\varphi }
{2} = \frac{2} 
{2\sqrt{10}}\mathrel{\implies }\varphi \mathop{\mathop{\doteq }}\nolimits 35^{\circ }6^{\prime}\)}{\(\mathop{\mathrm{tg}}\nolimits \frac{\varphi }
{2} = \frac{2\sqrt{2}} 
 {6} \mathrel{\implies }\varphi \mathop{\mathop{\doteq }}\nolimits 50^{\circ }29^{\prime}\)}{}{}}
\End

\Data\ID{9000153704}
\Updated{2020-07-15 03:07:12}
\Level{B}
\SubArea{Lines and planes: distances and angles}
\IMG{001537_04-figure0.pdf}
\LANGen{
\Question{
The picture shows a square pyramid. The side of a base square is
\(a = 4\; \mathrm{cm}\) and the height of
the pyramid is \(v = 6\; \mathrm{cm}\).
Find the angle \(\varphi \).}
{\(\mathop{\mathrm{tg}}\nolimits \frac{\varphi }
{2} = \frac{2} 
{2\sqrt{10}}\mathrel{\implies }\varphi \mathop{\mathop{\doteq }}\nolimits 35^{\circ }6^{\prime}\)}{\(\mathop{\mathrm{tg}}\nolimits \varphi = \frac{6} 
{2\sqrt{2}}\mathrel{\implies }\varphi \mathop{\mathop{\doteq }}\nolimits 64^{\circ }46^{\prime}\)}{\(\mathop{\mathrm{tg}}\nolimits \frac{\varphi }
{2} = \frac{2} 
{6}\mathrel{\implies }\varphi \mathop{\mathop{\doteq }}\nolimits 36^{\circ }52^{\prime}\)}{\(\mathop{\mathrm{tg}}\nolimits \frac{\varphi }
{2} = \frac{2\sqrt{2}} 
 {6} \mathrel{\implies }\varphi \mathop{\mathop{\doteq }}\nolimits 50^{\circ }29^{\prime}\)}{}{}}
\LANGpl{
\Question{
Rysunek przedstawia ostrosłup prawidłowy czworokątny o boku podstawy  równym
\(a = 4\; \mathrm{cm}\) i wysokości
 \(v = 6\; \mathrm{cm}\).
Wyznacz kąt \(\varphi \).}
{\(\mathop{\mathrm{tg}}\nolimits \frac{\varphi }
{2} = \frac{2} 
{2\sqrt{10}}\mathrel{\implies }\varphi \mathop{\mathop{\doteq }}\nolimits 35^{\circ }6^{\prime}\)}{\(\mathop{\mathrm{tg}}\nolimits \varphi = \frac{6} 
{2\sqrt{2}}\mathrel{\implies }\varphi \mathop{\mathop{\doteq }}\nolimits 64^{\circ }46^{\prime}\)}{\(\mathop{\mathrm{tg}}\nolimits \frac{\varphi }
{2} = \frac{2} 
{6}\mathrel{\implies }\varphi \mathop{\mathop{\doteq }}\nolimits 36^{\circ }52^{\prime}\)}{\(\mathop{\mathrm{tg}}\nolimits \frac{\varphi }
{2} = \frac{2\sqrt{2}} 
 {6} \mathrel{\implies }\varphi \mathop{\mathop{\doteq }}\nolimits 50^{\circ }29^{\prime}\)}{}{}}
\LANGcs{
\Question{
Na obrázku je pravidelný čtyřboký jehlan se čtvercovou podstavou o hraně
\(a = 4\; \mathrm{cm}\) s tělesovou
výškou \(v = 6\; \mathrm{cm}\).
Pro velikost vyznačené odchylky platí:}
{\(\mathop{\mathrm{tg}}\nolimits \frac{\varphi }
{2} = \frac{2} 
{2\sqrt{10}}\mathrel{\implies }\varphi \mathop{\mathop{\doteq }}\nolimits 35^{\circ }6^{\prime}\)}{\(\mathop{\mathrm{tg}}\nolimits \varphi = \frac{6} 
{2\sqrt{2}}\mathrel{\implies }\varphi \mathop{\mathop{\doteq }}\nolimits 64^{\circ }46^{\prime}\)}{\(\mathop{\mathrm{tg}}\nolimits \frac{\varphi }
{2} = \frac{2} 
{6}\mathrel{\implies }\varphi \mathop{\mathop{\doteq }}\nolimits 36^{\circ }52^{\prime}\)}{\(\mathop{\mathrm{tg}}\nolimits \frac{\varphi }
{2} = \frac{2\sqrt{2}} 
 {6} \mathrel{\implies }\varphi \mathop{\mathop{\doteq }}\nolimits 50^{\circ }29^{\prime}\)}{}{}}
\LANGsk{
\Question{
Na obrázku je pravidelný štvorboký ihlan so štvorcovou podstavou s hranou
\(a = 4\; \mathrm{cm}\) s telesovou
výškou \(v = 6\; \mathrm{cm}\).
Pre veľkosť vyznačené odchýlky platí:}
{\(\mathop{\mathrm{tg}}\nolimits \frac{\varphi }
{2} = \frac{2} 
{2\sqrt{10}}\mathrel{\implies }\varphi \mathop{\mathop{\doteq }}\nolimits 35^{\circ }6^{\prime}\)}{\(\mathop{\mathrm{tg}}\nolimits \varphi = \frac{6} 
{2\sqrt{2}}\mathrel{\implies }\varphi \mathop{\mathop{\doteq }}\nolimits 64^{\circ }46^{\prime}\)}{\(\mathop{\mathrm{tg}}\nolimits \frac{\varphi }
{2} = \frac{2} 
{6}\mathrel{\implies }\varphi \mathop{\mathop{\doteq }}\nolimits 36^{\circ }52^{\prime}\)}{\(\mathop{\mathrm{tg}}\nolimits \frac{\varphi }
{2} = \frac{2\sqrt{2}} 
 {6} \mathrel{\implies }\varphi \mathop{\mathop{\doteq }}\nolimits 50^{\circ }29^{\prime}\)}{}{}}
\LANGes{
\Question{
La imagen muestra una pirámide cuadrangular. La arista de una base cuadrada es
\(a = 4\; \mathrm{cm}\) y la altura de la pirámide es \(v = 6\; \mathrm{cm}\).
Determina el ángulo  \(\varphi \).}
{\(\mathop{\mathrm{tg}}\nolimits \frac{\varphi }
{2} = \frac{2} 
{2\sqrt{10}}\mathrel{\implies }\varphi \mathop{\mathop{\doteq }}\nolimits 35^{\circ }6^{\prime}\)}{\(\mathop{\mathrm{tg}}\nolimits \varphi = \frac{6} 
{2\sqrt{2}}\mathrel{\implies }\varphi \mathop{\mathop{\doteq }}\nolimits 64^{\circ }46^{\prime}\)}{\(\mathop{\mathrm{tg}}\nolimits \frac{\varphi }
{2} = \frac{2} 
{6}\mathrel{\implies }\varphi \mathop{\mathop{\doteq }}\nolimits 36^{\circ }52^{\prime}\)}{\(\mathop{\mathrm{tg}}\nolimits \frac{\varphi }
{2} = \frac{2\sqrt{2}} 
 {6} \mathrel{\implies }\varphi \mathop{\mathop{\doteq }}\nolimits 50^{\circ }29^{\prime}\)}{}{}}
\End

\Data\ID{9000153705}
\Updated{2020-07-15 03:07:12}
\Level{B}
\SubArea{Lines and planes: distances and angles}
\IMG{001537_05-figure0.pdf}
\LANGen{
\Question{
The picture shows a square pyramid. The side of a base square is
\(a = 4\; \mathrm{cm}\) and the height of
the pyramid is \(v = 6\; \mathrm{cm}\).
Find the angle \(\varphi \).}
{\(\mathop{\mathrm{tg}}\nolimits \frac{\varphi }
{2} = \frac{2\sqrt{2}} 
 {6} \mathrel{\implies }\varphi \mathop{\mathop{\doteq }}\nolimits 50^{\circ }29^{\prime}\)}{\(\mathop{\mathrm{tg}}\nolimits \varphi = \frac{6} 
{2\sqrt{2}}\mathrel{\implies }\varphi \mathop{\mathop{\doteq }}\nolimits 64^{\circ }46^{\prime}\)}{\(\mathop{\mathrm{tg}}\nolimits \frac{\varphi }
{2} = \frac{2} 
{6}\mathrel{\implies }\varphi \mathop{\mathop{\doteq }}\nolimits 36^{\circ }52^{\prime}\)}{\(\mathop{\mathrm{tg}}\nolimits \frac{\varphi }
{2} = \frac{2} 
{2\sqrt{10}}\mathrel{\implies }\varphi \mathop{\mathop{\doteq }}\nolimits 35^{\circ }6^{\prime}\)}{}{}}
\LANGpl{
\Question{
Rysunek przedstawia ostrosłup prawidłowy czworokątny o boku podstawy  równym
\(a = 4\; \mathrm{cm}\)   i wysokości
 \(v = 6\; \mathrm{cm}\).
Wyznacz kąt  \(\varphi \).}
{\(\mathop{\mathrm{tg}}\nolimits \frac{\varphi }
{2} = \frac{2\sqrt{2}} 
 {6} \mathrel{\implies }\varphi \mathop{\mathop{\doteq }}\nolimits 50^{\circ }29^{\prime}\)}{\(\mathop{\mathrm{tg}}\nolimits \varphi = \frac{6} 
{2\sqrt{2}}\mathrel{\implies }\varphi \mathop{\mathop{\doteq }}\nolimits 64^{\circ }46^{\prime}\)}{\(\mathop{\mathrm{tg}}\nolimits \frac{\varphi }
{2} = \frac{2} 
{6}\mathrel{\implies }\varphi \mathop{\mathop{\doteq }}\nolimits 36^{\circ }52^{\prime}\)}{\(\mathop{\mathrm{tg}}\nolimits \frac{\varphi }
{2} = \frac{2} 
{2\sqrt{10}}\mathrel{\implies }\varphi \mathop{\mathop{\doteq }}\nolimits 35^{\circ }6^{\prime}\)}{}{}}
\LANGcs{
\Question{
Na obrázku je pravidelný čtyřboký jehlan se čtvercovou podstavou o hraně
\(a = 4\; \mathrm{cm}\) s tělesovou
výškou \(v = 6\; \mathrm{cm}\).
Pro velikost vyznačené odchylky platí:}
{\(\mathop{\mathrm{tg}}\nolimits \frac{\varphi }
{2} = \frac{2\sqrt{2}} 
 {6} \mathrel{\implies }\varphi \mathop{\mathop{\doteq }}\nolimits 50^{\circ }29^{\prime}\)}{\(\mathop{\mathrm{tg}}\nolimits \varphi = \frac{6} 
{2\sqrt{2}}\mathrel{\implies }\varphi \mathop{\mathop{\doteq }}\nolimits 64^{\circ }46^{\prime}\)}{\(\mathop{\mathrm{tg}}\nolimits \frac{\varphi }
{2} = \frac{2} 
{6}\mathrel{\implies }\varphi \mathop{\mathop{\doteq }}\nolimits 36^{\circ }52^{\prime}\)}{\(\mathop{\mathrm{tg}}\nolimits \frac{\varphi }
{2} = \frac{2} 
{2\sqrt{10}}\mathrel{\implies }\varphi \mathop{\mathop{\doteq }}\nolimits 35^{\circ }6^{\prime}\)}{}{}}
\LANGsk{
\Question{
Na obrázku je pravidelný štvorboký ihlan so štvorcovou podstavou s hranou
\(a = 4\; \mathrm{cm}\) s telesovou
výškou \(v = 6\; \mathrm{cm}\).
Pre veľkosť vyznačenej odchýlky platí:}
{\(\mathop{\mathrm{tg}}\nolimits \frac{\varphi }
{2} = \frac{2\sqrt{2}} 
 {6} \mathrel{\implies }\varphi \mathop{\mathop{\doteq }}\nolimits 50^{\circ }29^{\prime}\)}{\(\mathop{\mathrm{tg}}\nolimits \varphi = \frac{6} 
{2\sqrt{2}}\mathrel{\implies }\varphi \mathop{\mathop{\doteq }}\nolimits 64^{\circ }46^{\prime}\)}{\(\mathop{\mathrm{tg}}\nolimits \frac{\varphi }
{2} = \frac{2} 
{6}\mathrel{\implies }\varphi \mathop{\mathop{\doteq }}\nolimits 36^{\circ }52^{\prime}\)}{\(\mathop{\mathrm{tg}}\nolimits \frac{\varphi }
{2} = \frac{2} 
{2\sqrt{10}}\mathrel{\implies }\varphi \mathop{\mathop{\doteq }}\nolimits 35^{\circ }6^{\prime}\)}{}{}}
\LANGes{
\Question{
La imagen muestra una pirámide cuadrangular. La arista de una base cuadrada es
\(a = 4\; \mathrm{cm}\) y la altura de la pirámide es \(v = 6\; \mathrm{cm}\).
Determina el ángulo  \(\varphi \).}
{\(\mathop{\mathrm{tg}}\nolimits \frac{\varphi }
{2} = \frac{2\sqrt{2}} 
 {6} \mathrel{\implies }\varphi \mathop{\mathop{\doteq }}\nolimits 50^{\circ }29^{\prime}\)}{\(\mathop{\mathrm{tg}}\nolimits \varphi = \frac{6} 
{2\sqrt{2}}\mathrel{\implies }\varphi \mathop{\mathop{\doteq }}\nolimits 64^{\circ }46^{\prime}\)}{\(\mathop{\mathrm{tg}}\nolimits \frac{\varphi }
{2} = \frac{2} 
{6}\mathrel{\implies }\varphi \mathop{\mathop{\doteq }}\nolimits 36^{\circ }52^{\prime}\)}{\(\mathop{\mathrm{tg}}\nolimits \frac{\varphi }
{2} = \frac{2} 
{2\sqrt{10}}\mathrel{\implies }\varphi \mathop{\mathop{\doteq }}\nolimits 35^{\circ }6^{\prime}\)}{}{}}
\End

\Data\ID{9000153706}
\Updated{2020-07-15 03:07:12}
\Level{B}
\SubArea{Lines and planes: distances and angles}
\IMG{001537_06-figure0.pdf}
\LANGen{
\Question{
The picture shows a square pyramid. The side of a base square is
\(a = 4\; \mathrm{cm}\) and the height of
the pyramid is \(v = 6\; \mathrm{cm}\).
Find the angle \(\varphi \).}
{\(\mathop{\mathrm{tg}}\nolimits \varphi = \frac{2\sqrt{10}} 
 {2} \mathrel{\implies }\varphi \mathop{\mathop{\doteq }}\nolimits 72^{\circ }27^{\prime}\)}{\(\mathop{\mathrm{tg}}\nolimits \varphi = \frac{6} 
{2\sqrt{2}}\mathrel{\implies }\varphi \mathop{\mathop{\doteq }}\nolimits 64^{\circ }46^{\prime}\)}{\(\mathop{\mathrm{tg}}\nolimits \varphi = \frac{6} 
{2}\mathrel{\implies }\varphi \mathop{\mathop{\doteq }}\nolimits 71^{\circ }34^{\prime}\)}{\(\mathop{\mathrm{tg}}\nolimits \frac{\varphi }
{2} = \frac{2\sqrt{2}} 
 {6} \mathrel{\implies }\varphi \mathop{\mathop{\doteq }}\nolimits 50^{\circ }29^{\prime}\)}{}{}}
\LANGpl{
\Question{
Rysunek przedstawia ostrosłup prawidłowy czworokątny o boku podstawy  równym
\(a = 4\; \mathrm{cm}\) i wysokości
 \(v = 6\; \mathrm{cm}\).
Wyznacz kąt \(\varphi \).}
{\(\mathop{\mathrm{tg}}\nolimits \varphi = \frac{2\sqrt{10}} 
 {2} \mathrel{\implies }\varphi \mathop{\mathop{\doteq }}\nolimits 72^{\circ }27^{\prime}\)}{\(\mathop{\mathrm{tg}}\nolimits \varphi = \frac{6} 
{2\sqrt{2}}\mathrel{\implies }\varphi \mathop{\mathop{\doteq }}\nolimits 64^{\circ }46^{\prime}\)}{\(\mathop{\mathrm{tg}}\nolimits \varphi = \frac{6} 
{2}\mathrel{\implies }\varphi \mathop{\mathop{\doteq }}\nolimits 71^{\circ }34^{\prime}\)}{\(\mathop{\mathrm{tg}}\nolimits \frac{\varphi }
{2} = \frac{2\sqrt{2}} 
 {6} \mathrel{\implies }\varphi \mathop{\mathop{\doteq }}\nolimits 50^{\circ }29^{\prime}\)}{}{}}
\LANGcs{
\Question{
Na obrázku je pravidelný čtyřboký jehlan se čtvercovou podstavou o hraně
\(a = 4\; \mathrm{cm}\) s tělesovou
výškou \(v = 6\; \mathrm{cm}\).
Pro velikost vyznačené odchylky platí:}
{\(\mathop{\mathrm{tg}}\nolimits \varphi = \frac{2\sqrt{10}} 
 {2} \mathrel{\implies }\varphi \mathop{\mathop{\doteq }}\nolimits 72^{\circ }27^{\prime}\)}{\(\mathop{\mathrm{tg}}\nolimits \varphi = \frac{6} 
{2\sqrt{2}}\mathrel{\implies }\varphi \mathop{\mathop{\doteq }}\nolimits 64^{\circ }46^{\prime}\)}{\(\mathop{\mathrm{tg}}\nolimits \varphi = \frac{6} 
{2}\mathrel{\implies }\varphi \mathop{\mathop{\doteq }}\nolimits 71^{\circ }34^{\prime}\)}{\(\mathop{\mathrm{tg}}\nolimits \frac{\varphi }
{2} = \frac{2\sqrt{2}} 
 {6} \mathrel{\implies }\varphi \mathop{\mathop{\doteq }}\nolimits 50^{\circ }29^{\prime}\)}{}{}}
\LANGsk{
\Question{
Na obrázku je pravidelný štvorboký ihlan so štvorcovou podstavou s hranou
\(a = 4\; \mathrm{cm}\) s telesovou
výškou \(v = 6\; \mathrm{cm}\).
Pre veľkosť vyznačenej odchýlky platí:}
{\(\mathop{\mathrm{tg}}\nolimits \varphi = \frac{2\sqrt{10}} 
 {2} \mathrel{\implies }\varphi \mathop{\mathop{\doteq }}\nolimits 72^{\circ }27^{\prime}\)}{\(\mathop{\mathrm{tg}}\nolimits \varphi = \frac{6} 
{2\sqrt{2}}\mathrel{\implies }\varphi \mathop{\mathop{\doteq }}\nolimits 64^{\circ }46^{\prime}\)}{\(\mathop{\mathrm{tg}}\nolimits \varphi = \frac{6} 
{2}\mathrel{\implies }\varphi \mathop{\mathop{\doteq }}\nolimits 71^{\circ }34^{\prime}\)}{\(\mathop{\mathrm{tg}}\nolimits \frac{\varphi }
{2} = \frac{2\sqrt{2}} 
 {6} \mathrel{\implies }\varphi \mathop{\mathop{\doteq }}\nolimits 50^{\circ }29^{\prime}\)}{}{}}
\LANGes{
\Question{
La imagen muestra una pirámide cuadrangular. La arista de una base cuadrada es 
\(a = 4\; \mathrm{cm}\) y la altura de la pirámide es \(v = 6\; \mathrm{cm}\).
Determina el ángulo \(\varphi \).}
{\(\mathop{\mathrm{tg}}\nolimits \varphi = \frac{2\sqrt{10}} 
 {2} \mathrel{\implies }\varphi \mathop{\mathop{\doteq }}\nolimits 72^{\circ }27^{\prime}\)}{\(\mathop{\mathrm{tg}}\nolimits \varphi = \frac{6} 
{2\sqrt{2}}\mathrel{\implies }\varphi \mathop{\mathop{\doteq }}\nolimits 64^{\circ }46^{\prime}\)}{\(\mathop{\mathrm{tg}}\nolimits \varphi = \frac{6} 
{2}\mathrel{\implies }\varphi \mathop{\mathop{\doteq }}\nolimits 71^{\circ }34^{\prime}\)}{\(\mathop{\mathrm{tg}}\nolimits \frac{\varphi }
{2} = \frac{2\sqrt{2}} 
 {6} \mathrel{\implies }\varphi \mathop{\mathop{\doteq }}\nolimits 50^{\circ }29^{\prime}\)}{}{}}
\End

\Data\ID{1103025401}
\Updated{2020-07-15 03:07:12}
\Level{C}
\SubArea{Lines and planes: distances and angles}
\IMG{1103025401.pdf}
\LANGen{
\Question{
A regular hexagonal prism \( ABCDEFA'B'C'D'E'F' \) has the side \( a \) of the length \( 3\,\mathrm{cm} \) and the height \( v \) of the length \( 8\,\mathrm{cm} \). Find the angle between the lines \( AD' \) and \( BD' \). Round the result to two decimal places.}
{\( 17.46^{\circ} \)}{\( 72.54^{\circ} \)}{\( 16.70^{\circ} \)}{\( 20.57^{\circ} \)}{}{}}
\LANGpl{
\Question{
Dany jest graniastosłup prawidłowy sześciokątny  \( ABCDEFA'B'C'D'E'F' \) o boku \( a \) o długości  \( 3\,\mathrm{cm} \), wysokości  \( v \) równej \( 8\,\mathrm{cm} \). Wyznacz kąt między prostymi \( AD' \) i \( BD' \). Zaokrągli wynik do dwóch miejsc po przecinku.}
{\( 17{,}46^{\circ} \)}{\( 72{,}54^{\circ} \)}{\( 16{,}70^{\circ} \)}{\( 20{,}57^{\circ} \)}{}{}}
\LANGcs{
\Question{
Je dán pravidelný šestiboký hranol \( ABCDEFA'B'C'D'E'F' \), délka podstavné hrany \( a = 3\,\mathrm{cm} \), výška \( v = 8\,\mathrm{cm} \). Určete odchylku přímek \( AD' \) a \( BD' \). Výsledek zaokrouhlete na dvě desetinná místa.}
{\( 17{,}46^{\circ} \)}{\( 72{,}54^{\circ} \)}{\( 16{,}70^{\circ} \)}{\( 20{,}57^{\circ} \)}{}{}}
\LANGsk{
\Question{
Je daný pravidelný šesťboký hranol \( ABCDEFA'B'C'D'E'F' \), dĺžka hrany podstavy \( a = 3\,\mathrm{cm} \), výška \( v = 8\,\mathrm{cm} \). Určte odchýlku priamok \( AD' \) a \( BD' \). Výsledok zaokrúhlite na dve desatinná miesta.}
{\( 17{,}46^{\circ} \)}{\( 72{,}54^{\circ} \)}{\( 16{,}70^{\circ} \)}{\( 20{,}57^{\circ} \)}{}{}}
\LANGes{
\Question{
Un prisma hexagonal regular \( ABCDEFA'B'C'D'E'F' \) tiene la arista básica \( a \) de \( 3\,\mathrm{cm} \) y la altura \( v \) de \( 8\,\mathrm{cm} \). Determina el ángulo entre las rectas \( AD' \) y \( BD' \). Redondea el resultado a dos posiciones decimales.}
{\( 17.46^{\circ} \)}{\( 72.54^{\circ} \)}{\( 16.70^{\circ} \)}{\( 20.57^{\circ} \)}{}{}}
\End

\Data\ID{1103025402}
\Updated{2020-07-15 03:07:12}
\Level{C}
\SubArea{Lines and planes: distances and angles}
\IMG{1103025402.pdf}
\LANGen{
\Question{
A regular hexagonal prism \( ABCDEFA'B'C'D'E'F' \) has the side \( a \) of the length \( 3\,\mathrm{cm} \) and the height \( v \) of the length \( 8\,\mathrm{cm} \). Find the angle between the plane \( ABD' \) and the base plane \( ABC \). Round the result to two decimal places.}
{\( 57^{\circ} \)}{\( 53.13^{\circ} \)}{\( 33^{\circ} \)}{\( 72.01^{\circ} \)}{}{}}
\LANGpl{
\Question{
Dany jest graniastosłup prawidłowy sześciokątny  \( ABCDEFA'B'C'D'E'F' \) o boku \( a \) równym \( 3\,\mathrm{cm} \) i wysokości  \( v \) równej \( 8\,\mathrm{cm} \). Wyznacz kąt między płaszczyzną  \( ABD' \), a płaszczyzną  \( ABC \). Zaokrągli wynik do dwóch miejsc po przecinku.}
{\( 57^{\circ} \)}{\( 53{,}13^{\circ} \)}{\( 33^{\circ} \)}{\( 72{,}01^{\circ} \)}{}{}}
\LANGcs{
\Question{
Je dán pravidelný šestiboký hranol \( ABCDEFA'B'C'D'E'F' \), délka podstavné hrany \( a = 3\,\mathrm{cm} \), výška \( v = 8\,\mathrm{cm} \). Určete odchylku roviny \( ABD' \) a roviny dolní podstavy \( ABC \). Výsledek zaokrouhlete na dvě desetinná místa.}
{\( 57^{\circ} \)}{\( 53{,}13^{\circ} \)}{\( 33^{\circ} \)}{\( 72{,}01^{\circ} \)}{}{}}
\LANGsk{
\Question{
Je daný pravidelný šesťboký hranol \( ABCDEFA'B'C'D'E'F' \), dĺžka hrany podstavy \( a = 3\,\mathrm{cm} \), výška \( v = 8\,\mathrm{cm} \). Určte odchýlku roviny \( ABD' \) a roviny dolnej podstavy \( ABC \). Výsledok zaokrúhlite na dve desatinné miesta.}
{\( 57^{\circ} \)}{\( 53{,}13^{\circ} \)}{\( 33^{\circ} \)}{\( 72{,}01^{\circ} \)}{}{}}
\LANGes{
\Question{
Un prisma hexagonal regular \( ABCDEFA'B'C'D'E'F' \) tiene la arista básica \( a \) de \( 3\,\mathrm{cm} \) y la altura \( v \) de \( 8\,\mathrm{cm} \). Determina el ángulo entre el plano \( ABD' \) y el plano de la base \( ABC \). Redondea el resultado a dos posiciones decimales.}
{\( 57^{\circ} \)}{\( 53.13^{\circ} \)}{\( 33^{\circ} \)}{\( 72.01^{\circ} \)}{}{}}
\End

\Data\ID{1103025403}
\Updated{2020-07-15 03:07:12}
\Level{C}
\SubArea{Lines and planes: distances and angles}
\IMG{1103025403.pdf}
\LANGen{
\Question{
A regular hexagonal prism \( ABCDEFA'B'C'D'E'F' \) has the side \( a \) of the length \( 3\,\mathrm{cm} \) and the height \( v \) of the length \( 8\,\mathrm{cm} \). Find the distance between the lines \( FA' \) and \( CD' \).}
{\( 3\sqrt3 \)}{\( 6 \)}{\( 6\sqrt3 \)}{\( \frac32\sqrt3 \)}{}{}}
\LANGpl{
\Question{
Dany jest graniastosłup prawidłowy sześciokątny  \( ABCDEFA'B'C'D'E'F' \) o boku \( a \) równym \( 3\,\mathrm{cm} \) i wysokości \( v \) równej  \( 8\,\mathrm{cm} \). Oblicz odległość miedzy prostymi \( FA' \) i \( CD' \).}
{\( 3\sqrt3 \)}{\( 6 \)}{\( 6\sqrt3 \)}{\( \frac32\sqrt3 \)}{}{}}
\LANGcs{
\Question{
Je dán pravidelný šestiboký hranol \( ABCDEFA'B'C'D'E'F' \), délka podstavné hrany \( a = 3\,\mathrm{cm} \), výška \( v = 8\,\mathrm{cm} \). Určete vzdálenost přímek \( FA' \) a \( CD' \).}
{\( 3\sqrt3 \)}{\( 6 \)}{\( 6\sqrt3 \)}{\( \frac32\sqrt3 \)}{}{}}
\LANGsk{
\Question{
Je daný pravidelný šesťboký hranol \( ABCDEFA'B'C'D'E'F' \), dĺžka hrany podstavy \( a = 3\,\mathrm{cm} \), výška \( v = 8\,\mathrm{cm} \). Určte vzdialenosť priamok \( FA' \) a \( CD' \).}
{\( 3\sqrt3 \)}{\( 6 \)}{\( 6\sqrt3 \)}{\( \frac32\sqrt3 \)}{}{}}
\LANGes{
\Question{
Un prisma hexagonal regular \( ABCDEFA'B'C'D'E'F' \) tiene la arista básica \( a \) de \( 3\,\mathrm{cm} \) y la altura  \( v \) de \( 8\,\mathrm{cm} \). Determina la distancia entre las rectas \( FA' \) y \( CD' \).}
{\( 3\sqrt3 \)}{\( 6 \)}{\( 6\sqrt3 \)}{\( \frac32\sqrt3 \)}{}{}}
\End

\Data\ID{1103025404}
\Updated{2020-07-15 03:07:12}
\Level{C}
\SubArea{Lines and planes: distances and angles}
\IMG{1103025404.pdf}
\LANGen{
\Question{
A regular hexagonal prism \( ABCDEFA'B'C'D'E'F' \) has the side \( a \) of the length \( 3\,\mathrm{cm} \) and the height \( v \) of the length \( 8\,\mathrm{cm} \). Find the distance between the planes \( AEE' \) and \( BDD' \).}
{\( 3 \)}{\( \sqrt3 \)}{\( 2\sqrt3 \)}{\( \frac{\sqrt3}2 \)}{}{}}
\LANGpl{
\Question{
Dany jest graniastosłup prawidłowy sześciokątny  \( ABCDEFA'B'C'D'E'F' \) o boku \( a \) równym \( 3\,\mathrm{cm} \) i wysokość \( v \) równej \( 8\,\mathrm{cm} \). Oblicz odległość miedzy płaszczyznami \( AEE' \) i  \( BDD' \).}
{\( 3 \)}{\( \sqrt3 \)}{\( 2\sqrt3 \)}{\( \frac{\sqrt3}2 \)}{}{}}
\LANGcs{
\Question{
Je dán pravidelný šestiboký hranol \( ABCDEFA'B'C'D'E'F' \), délka podstavné hrany \( a = 3\,\mathrm{cm} \), výška \( v = 8\,\mathrm{cm} \). Určete vzdálenost rovin \( AEE' \) a \( BDD' \).}
{\( 3 \)}{\( \sqrt3 \)}{\( 2\sqrt3 \)}{\( \frac{\sqrt3}2 \)}{}{}}
\LANGsk{
\Question{
Je daný pravidelný šesťboký hranol \( ABCDEFA'B'C'D'E'F' \), dĺžka hrany podstavy \( a = 3\,\mathrm{cm} \), výška \( v = 8\,\mathrm{cm} \). Určte vzdialenosť rovín \( AEE' \) a \( BDD' \).}
{\( 3 \)}{\( \sqrt3 \)}{\( 2\sqrt3 \)}{\( \frac{\sqrt3}2 \)}{}{}}
\LANGes{
\Question{
Un prisma hexagonal regular \( ABCDEFA'B'C'D'E'F' \) tiene la arista básica \( a \) de \( 3\,\mathrm{cm} \) y la altura \( v \) de \( 8\,\mathrm{cm} \). Determina la distancia entre los planos \( AEE' \) y \( BDD' \).}
{\( 3 \)}{\( \sqrt3 \)}{\( 2\sqrt3 \)}{\( \frac{\sqrt3}2 \)}{}{}}
\End

\Data\ID{1103101201}
\Updated{2020-07-15 03:07:21}
\Level{C}
\SubArea{Lines and planes: distances and angles}
\IMG{1103101201.pdf}
\LANGen{
\Question{
Let \( ABCDEFV \) be a regular hexagonal pyramid with a base edge length of \( 4\,\mathrm{cm} \) and a height of \( 8\,\mathrm{cm} \). Find the distance between the point \( V \) and the line \( BC \) (see the picture).}
{\( 2\sqrt{19}\,\mathrm{cm} \)}{\( 8\,\mathrm{cm} \)}{\( \left(8+2\sqrt{3}\right)\,\mathrm{cm} \)}{\( 4\sqrt{5}\,\mathrm{cm} \)}{}{}}
\LANGpl{
\Question{
Dany jest ostrosłup sześciokątny foremny  \( ABCDEFV \), krawędź podstawy jest równa  \( 4\,\mathrm{cm} \), a wysokość  \( 8\,\mathrm{cm} \). Oblicz odległość punktu  \( V \) od prostej  \( BC \) (spójrz na rysunek).}
{\( 2\sqrt{19}\,\mathrm{cm} \)}{\( 8\,\mathrm{cm} \)}{\( \left(8+2\sqrt{3}\right)\,\mathrm{cm} \)}{\( 4\sqrt{5}\,\mathrm{cm} \)}{}{}}
\LANGcs{
\Question{
Je dán pravidelný šestiboký jehlan  \( ABCDEFV \). Podstavná hrana má velikost  \( a = 4\,\mathrm{cm} \) a výška \( v = 8\,\mathrm{cm} \). Určete vzdálenost vrcholu  \( V \) od přímky  \( BC \) (viz obrázek).}
{\( 2\sqrt{19}\,\mathrm{cm} \)}{\( 8\,\mathrm{cm} \)}{\( \left(8+2\sqrt{3}\right)\,\mathrm{cm} \)}{\( 4\sqrt{5}\,\mathrm{cm} \)}{}{}}
\LANGsk{
\Question{
Je daný pravidelný šesťboký ihlan  \( ABCDEFV \). Podstavná hrana má veľkosť  \( a = 4\,\mathrm{cm} \) a výška \( v = 8\,\mathrm{cm} \). Určte vzdialenosť vrcholu  \( V \) od priamky  \( BC \) (viď obrázok).}
{\( 2\sqrt{19}\,\mathrm{cm} \)}{\( 8\,\mathrm{cm} \)}{\( \left(8+2\sqrt{3}\right)\,\mathrm{cm} \)}{\( 4\sqrt{5}\,\mathrm{cm} \)}{}{}}
\LANGes{
\Question{
Sea  \( ABCDEFV \) una pirámide hexagonal regular con una longitud de la arista básica de \( 4\,\mathrm{cm} \) y una altura de \( 8\,\mathrm{cm} \). Determina la distancia entre el vértice \(V \) y la recta \(BC \) (mira la imagen).}
{\( 2\sqrt{19}\,\mathrm{cm} \)}{\( 8\,\mathrm{cm} \)}{\( \left(8+2\sqrt{3}\right)\,\mathrm{cm} \)}{\( 4\sqrt{5}\,\mathrm{cm} \)}{}{}}
\End

\Data\ID{1103101202}
\Updated{2020-07-15 03:07:21}
\Level{C}
\SubArea{Lines and planes: distances and angles}
\IMG{1103101202.pdf}
\LANGen{
\Question{
Let \( ABCDEFV \) be a regular hexagonal pyramid with a base edge length of \( 4\,\mathrm{cm} \) and a height of \( 8\,\mathrm{cm} \). Find the distance between the line \( AB \) and the line \( ED \) (see the picture).}
{\( 4\sqrt3\,\mathrm{cm} \)}{\( 8\,\mathrm{cm} \)}{\( 8\sqrt3\,\mathrm{cm} \)}{\( 2\sqrt3\,\mathrm{cm} \)}{}{}}
\LANGpl{
\Question{
Dany jest ostrosłup sześciokątny foremny  \( ABCDEFV \), krawędź podstawy jest równa  \( 4\,\mathrm{cm} \), a wysokość  \( 8\,\mathrm{cm} \). Oblicz odległość pomiędzy prostymi  \( AB \) i  \( ED \) (spójrz na rysunek).}
{\( 4\sqrt3\,\mathrm{cm} \)}{\( 8\,\mathrm{cm} \)}{\( 8\sqrt3\,\mathrm{cm} \)}{\( 2\sqrt3\,\mathrm{cm} \)}{}{}}
\LANGcs{
\Question{
Je dán pravidelný šestiboký jehlan  \( ABCDEFV \). Podstavná hrana má velikost  \( a = 4\,\mathrm{cm} \) a výška  \( v = 8\,\mathrm{cm} \). Určete vzdálenost přímek  \( AB \) a \( ED \) (viz obrázek).}
{\( 4\sqrt3\,\mathrm{cm} \)}{\( 8\,\mathrm{cm} \)}{\( 8\sqrt3\,\mathrm{cm} \)}{\( 2\sqrt3\,\mathrm{cm} \)}{}{}}
\LANGsk{
\Question{
Je daný pravidelný šesťboký ihlan  \( ABCDEFV \). Podstavná hrana má veľkosť  \( a = 4\,\mathrm{cm} \) a výška  \( v = 8\,\mathrm{cm} \). Určte vzdialenosť priamok  \( AB \) a \( ED \) (viď obrázok).}
{\( 4\sqrt3\,\mathrm{cm} \)}{\( 8\,\mathrm{cm} \)}{\( 8\sqrt3\,\mathrm{cm} \)}{\( 2\sqrt3\,\mathrm{cm} \)}{}{}}
\LANGes{
\Question{
Sea \( ABCDEFV \) una pirámide hexagonal regular con una longitud de la arista básica de \( 4\,\mathrm{cm} \) y una altura de  \( 8\,\mathrm{cm} \). Determina la distancia entre las rectas \( AB \) y \( ED \) (mira la imagen).}
{\( 4\sqrt3\,\mathrm{cm} \)}{\( 8\,\mathrm{cm} \)}{\( 8\sqrt3\,\mathrm{cm} \)}{\( 2\sqrt3\,\mathrm{cm} \)}{}{}}
\End

\Data\ID{1103101203}
\Updated{2020-07-15 03:07:21}
\Level{C}
\SubArea{Lines and planes: distances and angles}
\IMG{1103101203.pdf}
\LANGen{
\Question{
Let \( ABCDEFV \) be a regular hexagonal pyramid with a base edge length of \( 4\,\mathrm{cm} \) and a height of \( 4\sqrt3\,\mathrm{cm} \). Find the angle between the lines \( FV \) and \( CV \) (see the picture).}
{\( 60^{\circ} \)}{\( 45^{\circ} \)}{\( 72^{\circ} \)}{\( 30^{\circ} \)}{}{}}
\LANGpl{
\Question{
Dany jest ostrosłup sześciokątny foremny  \( ABCDEFV \), krawędź podstawy jest równa  \( 4\,\mathrm{cm} \) a wysokość \( 4\sqrt3\,\mathrm{cm} \). Oblicz miarę kąta pomiędzy prostymi  \( FV \) i \( CV \) (spójrz na rysunek).}
{\( 60^{\circ} \)}{\( 45^{\circ} \)}{\( 72^{\circ} \)}{\( 30^{\circ} \)}{}{}}
\LANGcs{
\Question{
Je dán pravidelný šestiboký jehlan  \( ABCDEFV \).  Podstavná hrana má velikost \( a = 4\,\mathrm{cm} \) a výška  \(v = 4\sqrt3\,\mathrm{cm} \). Určete odchylku přímek  \( FV \) a \( CV \) (viz obrázek).}
{\( 60^{\circ} \)}{\( 45^{\circ} \)}{\( 72^{\circ} \)}{\( 30^{\circ} \)}{}{}}
\LANGsk{
\Question{
Je daný pravidelný šesťboký ihlan  \( ABCDEFV \).  Podstavná hrana má veľkosť \( a = 4\,\mathrm{cm} \) a výška  \(v = 4\sqrt3\,\mathrm{cm} \). Určte odchýlku priamok  \( FV \) a \( CV \) (viď obrázok).}
{\( 60^{\circ} \)}{\( 45^{\circ} \)}{\( 72^{\circ} \)}{\( 30^{\circ} \)}{}{}}
\LANGes{
\Question{
Sea  \( ABCDEFV \) una pirámide hexagonal regular con una longitud de la arista básica de \( 4\,\mathrm{cm} \) y una altura de  \( 4\sqrt3\,\mathrm{cm} \). Determina el ángulo de las rectas \( FV \) y \( CV \) (mira la imagen).}
{\( 60^{\circ} \)}{\( 45^{\circ} \)}{\( 72^{\circ} \)}{\( 30^{\circ} \)}{}{}}
\End

\Data\ID{1103101204}
\Updated{2020-07-15 03:07:21}
\Level{C}
\SubArea{Lines and planes: distances and angles}
\IMG{1103101204.pdf}
\LANGen{
\Question{
Let \( ABCDEFV \) be a regular hexagonal pyramid with a base edge length of \( 4\,\mathrm{cm} \) and a height of \( 6\,\mathrm{cm} \). Let \( \varphi \) be the angle between the plane \( AFV \) and the base plane \( ABC \) (see the picture). Choose the correct expression for \( \varphi \):}
{\( \mathrm{tg}\,\varphi=\sqrt3 \)}{\( \sin\varphi=\sqrt3 \)}{\( \mathrm{tg}\,\varphi=\frac{\sqrt3}3 \)}{\( \mathrm{tg}\,\varphi=\frac32 \)}{}{}}
\LANGpl{
\Question{
Dany jest ostrosłup sześciokątny foremny  \( ABCDEFV \), krawędź podstawy jest równa  \( 4\,\mathrm{cm} \), a wysokość  \( 6\,\mathrm{cm} \). Niech  \( \varphi \) będzie kątem pomiędzy płaszczyzną  \( AFV \) a płaszczyzną podstawy \( ABC \) (spójrz na rysunek). Wybierz poprawne wyrażenie dla  \( \varphi \):}
{\( \mathrm{tg}\,\varphi=\sqrt3 \)}{\( \sin\varphi=\sqrt3 \)}{\( \mathrm{tg}\,\varphi=\frac{\sqrt3}3 \)}{\( \mathrm{tg}\,\varphi=\frac32 \)}{}{}}
\LANGcs{
\Question{
Je dán pravidelný šestiboký jehlan  \( ABCDEFV \). Podstavná hrana má velikost  \( a = 4\,\mathrm{cm} \) a výška  \( v = 6\,\mathrm{cm} \). Vyberte vztah, který platí pro odchylku  \( \varphi \) boční stěny  \( AFV \) a roviny podstavy \( ABC \) (viz obrázek):}
{\( \mathrm{tg}\,\varphi=\sqrt3 \)}{\( \sin\varphi=\sqrt3 \)}{\( \mathrm{tg}\,\varphi=\frac{\sqrt3}3 \)}{\( \mathrm{tg}\,\varphi=\frac32 \)}{}{}}
\LANGsk{
\Question{
Je daný pravidelný šesťboký ihlan  \( ABCDEFV \). Podstavná hrana má veľkosť  \( a = 4\,\mathrm{cm} \) a výška  \( v = 6\,\mathrm{cm} \). Vyberte vzťah, ktorý platí pre odchýlku  \( \varphi \) bočnej steny  \( AFV \) a roviny podstavy \( ABC \) (viď obrázok):}
{\( \mathrm{tg}\,\varphi=\sqrt3 \)}{\( \sin\varphi=\sqrt3 \)}{\( \mathrm{tg}\,\varphi=\frac{\sqrt3}3 \)}{\( \mathrm{tg}\,\varphi=\frac32 \)}{}{}}
\LANGes{
\Question{
Sea \( ABCDEFV \) una pirámide hexagonal regular con una longitud de la arista básica de \( 4\,\mathrm{cm} \)  y una altura de \( 6\,\mathrm{cm} \). Sea \( \varphi \) el ángulo entre el plano \( AFV \) y el plano de la base \( ABC \) (mira la imagen). Elige la expresión correcta para \( \varphi \):}
{\( \mathrm{tg}\,\varphi=\sqrt3 \)}{\( \sin\varphi=\sqrt3 \)}{\( \mathrm{tg}\,\varphi=\frac{\sqrt3}3 \)}{\( \mathrm{tg}\,\varphi=\frac32 \)}{}{}}
\End

\Data\ID{1103101205}
\Updated{2020-07-15 03:07:21}
\Level{C}
\SubArea{Lines and planes: distances and angles}
\IMG{1103101205.pdf}
\LANGen{
\Question{
Let \( ABCDEFV \) be a regular hexagonal pyramid with an edge \( AB \), where \( |AB| = 4\,\mathrm{cm} \), and an edge \( AV \), where \( |AV| = 8\,\mathrm{cm} \). Find the angle between the line \( AV \) and the base plane \( ABC \) (see the picture).}
{\( 60^{\circ} \)}{\( 45^{\circ} \)}{\( 72^{\circ} \)}{\( 30^{\circ} \)}{}{}}
\LANGpl{
\Question{
Dany jest ostrosłup sześciokątny foremny  \( ABCDEFV \) o krawędzi  \( AB \), gdzie \( |AB| = 4\,\mathrm{cm} \), i krawędzi  \( AV \), gdzie \( |AV| = 8\,\mathrm{cm} \). Oblicz miarę kąta pomiędzy prostą  \( AV \) a płaszczyzną podstawy  \( ABC \) (spójrz na rysunek).}
{\( 60^{\circ} \)}{\( 45^{\circ} \)}{\( 72^{\circ} \)}{\( 30^{\circ} \)}{}{}}
\LANGcs{
\Question{
Je dán pravidelný šestiboký jehlan  \( ABCDEFV \).  Podstavná hrana má velikost  \( |AB| = 4\,\mathrm{cm} \) a boční hrana má velikost \( |AV| = 8\,\mathrm{cm} \). Určete odchylku přímky  \( AV \) od roviny podstavy \( ABC \) (viz obrázek).}
{\( 60^{\circ} \)}{\( 45^{\circ} \)}{\( 72^{\circ} \)}{\( 30^{\circ} \)}{}{}}
\LANGsk{
\Question{
Je daný pravidelný šesťboký ihlan  \( ABCDEFV \).  Podstavná hrana má veľkosť  \( |AB| = 4\,\mathrm{cm} \) a bočná hrana má veľkosť \( |AV| = 8\,\mathrm{cm} \). Určte odchýlku priamky  \( AV \) od roviny podstavy \( ABC \) (viď obrázok).}
{\( 60^{\circ} \)}{\( 45^{\circ} \)}{\( 72^{\circ} \)}{\( 30^{\circ} \)}{}{}}
\LANGes{
\Question{
Sea \( ABCDEFV \) una pirámide hexagonal regular con la arista básica \( AB \), donde \( |AB| = 4\,\mathrm{cm} \), y la arista \( AV \), donde \( |AV| = 8\,\mathrm{cm} \). Determina el ángulo entre la recta \( AV \) y el plano de la base \( ABC \) (mira la imagen).}
{\( 60^{\circ} \)}{\( 45^{\circ} \)}{\( 72^{\circ} \)}{\( 30^{\circ} \)}{}{}}
\End

\Data\ID{1103101206}
\Updated{2020-07-15 03:07:21}
\Level{C}
\SubArea{Lines and planes: distances and angles}
\IMG{1103101206.pdf}
\LANGen{
\Question{
Let \( ABCDEFV \) be a regular hexagonal pyramid with a base edge length of \( 4\,\mathrm{cm} \) and a height of \( 10\,\mathrm{cm} \). Find the distance between the line \( AC \) and the line \( FD \) (see the picture).}
{\( 4\,\mathrm{cm} \)}{\( 2\sqrt3\,\mathrm{cm} \)}{\( 2\sqrt2\,\mathrm{cm} \)}{\( 4\sqrt3\,\mathrm{cm} \)}{}{}}
\LANGpl{
\Question{
Dany jest ostrosłup sześciokątny foremny  \( ABCDEFV \), krawędź podstawy jet równa \( 4\,\mathrm{cm} \) a wysokość \( 10\,\mathrm{cm} \). Oblicz odległość pomiędzy prostymi  \( AC \) i  \( FD \) (spójrz na rysunek).}
{\( 4\,\mathrm{cm} \)}{\( 2\sqrt3\,\mathrm{cm} \)}{\( 2\sqrt2\,\mathrm{cm} \)}{\( 4\sqrt3\,\mathrm{cm} \)}{}{}}
\LANGcs{
\Question{
Je dán pravidelný šestiboký jehlan  \( ABCDEFV \). Podstavná hrana má velikost  \( a = 4\,\mathrm{cm} \) a výška \( v = 10\,\mathrm{cm} \). Určete vzdálenost přímek  \( AC \) a \( FD \) (viz obrázek).}
{\( 4\,\mathrm{cm} \)}{\( 2\sqrt3\,\mathrm{cm} \)}{\( 2\sqrt2\,\mathrm{cm} \)}{\( 4\sqrt3\,\mathrm{cm} \)}{}{}}
\LANGsk{
\Question{
Je daný pravidelný šesťboký ihlan  \( ABCDEFV \). Podstavná hrana má veľkosť  \( a = 4\,\mathrm{cm} \) a výška \( v = 10\,\mathrm{cm} \). Určte vzdialenosť priamok  \( AC \) a \( FD \) (viď obrázok).}
{\( 4\,\mathrm{cm} \)}{\( 2\sqrt3\,\mathrm{cm} \)}{\( 2\sqrt2\,\mathrm{cm} \)}{\( 4\sqrt3\,\mathrm{cm} \)}{}{}}
\LANGes{
\Question{
Sea \( ABCDEFV \) una pirámide hexagonal regular con una longitud de la arista básica de \( 4\,\mathrm{cm} \) y una altura de  \( 10\,\mathrm{cm} \). Determina la distancia entre la recta \( AC \) y la recta \( FD \) (mira la imagen).}
{\( 4\,\mathrm{cm} \)}{\( 2\sqrt3\,\mathrm{cm} \)}{\( 2\sqrt2\,\mathrm{cm} \)}{\( 4\sqrt3\,\mathrm{cm} \)}{}{}}
\End

\Data\ID{1103101207}
\Updated{2020-07-15 03:07:21}
\Level{C}
\SubArea{Lines and planes: distances and angles}
\IMG{1103101207.pdf}
\LANGen{
\Question{
Let \( ABCD \) be a regular tetrahedron with an edge length of \( 5\,\mathrm{cm} \). Let \( \varphi \) be the angle between the edge \( AD \) and the base plane \( ABC \) (see the picture). Choose the correct expression for \( \varphi \):}
{\( \cos\varphi=\frac{\sqrt3}3 \)}{\( \mathrm{tg}\,\varphi=\sqrt3 \)}{\( \mathrm{tg}\,\varphi=\frac{\sqrt3}3 \)}{\( \cos\varphi=\sqrt3 \)}{}{}}
\LANGpl{
\Question{
Dany jest czworościan foremny  \( ABCD \), długość krawędzi jest równa  \( 5\,\mathrm{cm} \). Niech  \( \varphi \) będzie kątem pomiędzy krawędzią  \( AD \) a płaszczyzną podstawy \( ABC \) (spójrz na rysunek). Wybierz poprawne wyrażenie dla \( \varphi \):}
{\( \cos\varphi=\frac{\sqrt3}3 \)}{\( \mathrm{tg}\,\varphi=\sqrt3 \)}{\( \mathrm{tg}\,\varphi=\frac{\sqrt3}3 \)}{\( \cos\varphi=\sqrt3 \)}{}{}}
\LANGcs{
\Question{
Pravidelný čtyřstěn  \( ABCD \) má velikost hrany \( 5\,\mathrm{cm} \). Odchylku boční hrany \( AD \) od roviny podstavy \( ABC \) označme \( \varphi \).  Vyberte správný vztah pro odchylku \( \varphi \).}
{\( \cos\varphi=\frac{\sqrt3}3 \)}{\( \mathrm{tg}\,\varphi=\sqrt3 \)}{\( \mathrm{tg}\,\varphi=\frac{\sqrt3}3 \)}{\( \cos\varphi=\sqrt3 \)}{}{}}
\LANGsk{
\Question{
Pravidelný štvorsten  \( ABCD \) má veľkosť hrany \( 5\,\mathrm{cm} \). Odchýlku bočnej hrany \( AD \) od roviny podstavy \( ABC \) označme \( \varphi \).  Vyberte správny vzťah pre odchýlku \( \varphi \).}
{\( \cos\varphi=\frac{\sqrt3}3 \)}{\( \mathrm{tg}\,\varphi=\sqrt3 \)}{\( \mathrm{tg}\,\varphi=\frac{\sqrt3}3 \)}{\( \cos\varphi=\sqrt3 \)}{}{}}
\LANGes{
\Question{
Sea \( ABCD \)  un tetraedro regular con una longitud de la arista de \( 5\,\mathrm{cm} \). Sea \( \varphi \) el ángulo entre la arista \( AD \) y el plano de la base \( ABC \) (mira la imagen). Elige la expresión correcta para \( \varphi \):}
{\( \cos\varphi=\frac{\sqrt3}3 \)}{\( \mathrm{tg}\,\varphi=\sqrt3 \)}{\( \mathrm{tg}\,\varphi=\frac{\sqrt3}3 \)}{\( \cos\varphi=\sqrt3 \)}{}{}}
\End

\Data\ID{1103101208}
\Updated{2020-07-15 03:07:21}
\Level{C}
\SubArea{Lines and planes: distances and angles}
\IMG{1103101208.pdf}
\LANGen{
\Question{
Let \( ABCD \) be a regular tetrahedron with an edge length of \( 5\,\mathrm{cm} \). Let \( \varphi \) be the angle between the face \( ABD \) and the base plane \( ABC \) (see the picture). Choose the correct expression for \( \varphi \):}
{\( \cos\varphi=\frac13 \)}{\( \mathrm{tg}\,\varphi=6\sqrt3 \)}{\( \mathrm{tg}\,\varphi=3\sqrt2 \)}{\( \cos\varphi=\frac23 \)}{}{}}
\LANGpl{
\Question{
Dany jest czworościan foremny \( ABCD \) o krawędzi równej  \( 5\,\mathrm{cm} \). Niech \( \varphi \) będzie kątem pomiędzy ścianą  \( ABD \) a płaszczyzną podstawy  \( ABC \) (spójrz na rysunek). Wybierz poprawne wyrażenie dla  \( \varphi \):}
{\( \cos\varphi=\frac13 \)}{\( \mathrm{tg}\,\varphi=6\sqrt3 \)}{\( \mathrm{tg}\,\varphi=3\sqrt2 \)}{\( \cos\varphi=\frac23 \)}{}{}}
\LANGcs{
\Question{
Pravidelný čtyřstěn  \( ABCD \) má velikost hrany \( 5\,\mathrm{cm} \). Označme  \( \varphi \) odchylku boční stěny \( ABD \) a roviny podstavy \( ABC \) (viz obrázek). Vyberte správný vztah pro odchylku rovin \( \varphi \):}
{\( \cos\varphi=\frac13 \)}{\( \mathrm{tg}\,\varphi=6\sqrt3 \)}{\( \mathrm{tg}\,\varphi=3\sqrt2 \)}{\( \cos\varphi=\frac23 \)}{}{}}
\LANGsk{
\Question{
Pravidelný štvorsten  \( ABCD \) má veľkosť hrany \( 5\,\mathrm{cm} \). Označme  \( \varphi \) odchýlku bočnej steny \( ABD \) a roviny podstavy \( ABC \) (viď obrázok). Vyberte správny vzťah pre odchýlku rovín \( \varphi \):}
{\( \cos\varphi=\frac13 \)}{\( \mathrm{tg}\,\varphi=6\sqrt3 \)}{\( \mathrm{tg}\,\varphi=3\sqrt2 \)}{\( \cos\varphi=\frac23 \)}{}{}}
\LANGes{
\Question{
Sea \( ABCD \)  un tetraedro regular con una longitud de la arista de \( 5\,\mathrm{cm} \). Sea \( \varphi \) el ángulo entre el plano de cara  \( ABD \) y el plano de la base \( ABC \) (mira la imagen). Elige la expresión correcta para \( \varphi \):}
{\( \cos\varphi=\frac13 \)}{\( \mathrm{tg}\,\varphi=6\sqrt3 \)}{\( \mathrm{tg}\,\varphi=3\sqrt2 \)}{\( \cos\varphi=\frac23 \)}{}{}}
\End

\Data\ID{1103101209}
\Updated{2020-07-15 03:07:21}
\Level{C}
\SubArea{Lines and planes: distances and angles}
\IMG{1103101209.pdf}
\LANGen{
\Question{
Let \( ABCD \) be a regular tetrahedron with an edge length of \( 3\,\mathrm{cm} \). Find the height \( |DT| \) of the tetrahedron (see the picture).}
{\( \sqrt6\,\mathrm{cm} \)}{\( \frac{3\sqrt3}2\,\mathrm{cm} \)}{\( \frac{\sqrt3}2\,\mathrm{cm} \)}{\( 3\sqrt2\,\mathrm{cm} \)}{}{}}
\LANGpl{
\Question{
Dany jest czworościan foremny  \( ABCD \) o krawędzi równej  \( 3\,\mathrm{cm} \). Oblicz wysokość \( |DT| \) czworościanu  (spójrz na rysunek).}
{\( \sqrt6\,\mathrm{cm} \)}{\( \frac{3\sqrt3}2\,\mathrm{cm} \)}{\( \frac{\sqrt3}2\,\mathrm{cm} \)}{\( 3\sqrt2\,\mathrm{cm} \)}{}{}}
\LANGcs{
\Question{
Pravidelný čtyřstěn  \( ABCD \) má velikost hrany \( 3\,\mathrm{cm} \). Určete velikost výšky \( |DT| \) (viz obrázek).}
{\( \sqrt6\,\mathrm{cm} \)}{\( \frac{3\sqrt3}2\,\mathrm{cm} \)}{\( \frac{\sqrt3}2\,\mathrm{cm} \)}{\( 3\sqrt2\,\mathrm{cm} \)}{}{}}
\LANGsk{
\Question{
Pravidelný štvorsten  \( ABCD \) má veľkosť hrany \( 3\,\mathrm{cm} \). Určte veľkosť výšky \( |DT| \) (viď obrázok).}
{\( \sqrt6\,\mathrm{cm} \)}{\( \frac{3\sqrt3}2\,\mathrm{cm} \)}{\( \frac{\sqrt3}2\,\mathrm{cm} \)}{\( 3\sqrt2\,\mathrm{cm} \)}{}{}}
\LANGes{
\Question{
Sea \( ABCD \)  un tetraedro regular con una longitud de la arista de \( 3\,\mathrm{cm} \). Determina la altura \( |DT| \) del tetraedro (mira la imagen).}
{\( \sqrt6\,\mathrm{cm} \)}{\( \frac{3\sqrt3}2\,\mathrm{cm} \)}{\( \frac{\sqrt3}2\,\mathrm{cm} \)}{\( 3\sqrt2\,\mathrm{cm} \)}{}{}}
\End

\Data\ID{1103107001}
\Updated{2020-07-15 03:07:21}
\Level{C}
\SubArea{Lines and planes: distances and angles}
\IMG{1103107001.pdf}
\LANGen{
\Question{
Let \( ABCDEFA'B'C'D'E'F' \) be a regular hexagonal prism with the base edge length of \( 4\,\mathrm{cm} \) and the height of \( 8\,\mathrm{cm} \). Find the angle between the lines \( FC' \) and \( CF' \) (see the picture).}
{\( 90^{\circ} \)}{\( 45^{\circ} \)}{\( 60^{\circ} \)}{\( 72^{\circ} \)}{}{}}
\LANGpl{
\Question{
Dany jest graniastosłup prawidłowy sześciokątny  \( ABCDEFA'B'C'D'E'F' \), krawędź podstawy wynosi  \( 4\,\mathrm{cm} \) a wysokość  \( 8\,\mathrm{cm} \). Oblicz miarę kąta pomiędzy prostymi  \( FC' \) i \( CF' \) (spójrz na rysunek).}
{\( 90^{\circ} \)}{\( 45^{\circ} \)}{\( 60^{\circ} \)}{\( 72^{\circ} \)}{}{}}
\LANGcs{
\Question{
Pravidelný šestiboký hranol \( ABCDEFA'B'C'D'E'F' \) má délku podstavné hrany  \( 4\,\mathrm{cm} \) a výšku \( 8\,\mathrm{cm} \). Určete odchylku přímek  \( FC' \) a  \( CF' \) (viz obrázek).}
{\( 90^{\circ} \)}{\( 45^{\circ} \)}{\( 60^{\circ} \)}{\( 72^{\circ} \)}{}{}}
\LANGsk{
\Question{
Pravidelný šesťboký hranol \( ABCDEFA'B'C'D'E'F' \) má dĺžku podstavnej hrany  \( 4\,\mathrm{cm} \) a výšku \( 8\,\mathrm{cm} \). Určte odchýlku priamok  \( FC' \) a  \( CF' \) (viď obrázok).}
{\( 90^{\circ} \)}{\( 45^{\circ} \)}{\( 60^{\circ} \)}{\( 72^{\circ} \)}{}{}}
\LANGes{
\Question{
Sea \( ABCDEFA'B'C'D'E'F' \) un prisma hexagonal regular con la longitud de la arista básica de\( 4\,\mathrm{cm} \) y la altura de \( 8\,\mathrm{cm} \). Determina el ángulo entre las rectas \( FC' \) y \( CF' \) (mira la imagen).}
{\( 90^{\circ} \)}{\( 45^{\circ} \)}{\( 60^{\circ} \)}{\( 72^{\circ} \)}{}{}}
\End

\Data\ID{1103107002}
\Updated{2020-07-15 03:07:21}
\Level{C}
\SubArea{Lines and planes: distances and angles}
\IMG{1103107002.pdf}
\LANGen{
\Question{
Let \( ABCDEFA'B'C'D'E'F' \) be a regular hexagonal prism with the base edge length of \( 4\,\mathrm{cm} \) and the height of \( 8\,\mathrm{cm} \). Find the angle between the lines \( A'D' \) and \( C'D' \) (see the picture).}
{\( 60^{\circ} \)}{\( 45^{\circ} \)}{\( 72^{\circ} \)}{\( 54.7^{\circ} \)}{}{}}
\LANGpl{
\Question{
Dany jest graniastosłup prawidłowy sześciokątny \( ABCDEFA'B'C'D'E'F' \), krawędź podstawy  \( 4\,\mathrm{cm} \), a wysokość \( 8\,\mathrm{cm} \). Oblicz miarę kąta pomiędzy prostymi  \( A'D' \) i  \( C'D' \) (spójrz na rysunek).}
{\( 60^{\circ} \)}{\( 45^{\circ} \)}{\( 72^{\circ} \)}{\( 54{,}7^{\circ} \)}{}{}}
\LANGcs{
\Question{
Pravidelný šestiboký hranol  \( ABCDEFA'B'C'D'E'F' \) má délku podstavné hrany  \( 4\,\mathrm{cm} \) a výšku \( 8\,\mathrm{cm} \). Určete odchylku přímek  \( A'D' \) a \( C'D' \) (viz obrázek).}
{\( 60^{\circ} \)}{\( 45^{\circ} \)}{\( 72^{\circ} \)}{\( 54{,}7^{\circ} \)}{}{}}
\LANGsk{
\Question{
Pravidelný šesťboký hranol  \( ABCDEFA'B'C'D'E'F' \) má dĺžku podstavnej hrany  \( 4\,\mathrm{cm} \) a výšku \( 8\,\mathrm{cm} \). Určte odchýlku priamok  \( A'D' \) a \( C'D' \) (viď obrázok).}
{\( 60^{\circ} \)}{\( 45^{\circ} \)}{\( 72^{\circ} \)}{\( 54{,}7^{\circ} \)}{}{}}
\LANGes{
\Question{
Sea \( ABCDEFA'B'C'D'E'F' \) un prisma hexagonal regular con la longitud de la arista básica de \( 4\,\mathrm{cm} \) y la altura de \( 8\,\mathrm{cm} \). Determina el ángulo entre las rectas  \( A'D' \) y \( C'D' \) (mira la imagen).}
{\( 60^{\circ} \)}{\( 45^{\circ} \)}{\( 72^{\circ} \)}{\( 54.7^{\circ} \)}{}{}}
\End

\Data\ID{1103107003}
\Updated{2020-07-15 03:07:21}
\Level{C}
\SubArea{Lines and planes: distances and angles}
\IMG{1103107003.pdf}
\LANGen{
\Question{
Let \( ABCDEFA'B'C'D'E'F' \) be a regular hexagonal prism with the base edge length of \( 4\,\mathrm{cm} \) and the height of \( 8\,\mathrm{cm} \). Find the angle between the lines \( A'D' \) and \( CC' \) (see the picture).}
{\( 90^{\circ} \)}{\( 60^{\circ} \)}{\( 45^{\circ} \)}{\( 72^{\circ} \)}{}{}}
\LANGpl{
\Question{
Dany jest graniastosłup prawidłowy sześciokątny \( ABCDEFA'B'C'D'E'F' \), krawędź podstawy jest równa  \( 4\,\mathrm{cm} \), a wysokość \( 8\,\mathrm{cm} \). Oblicz miarę kąta pomiędzy prostymi \( A'D' \) i \( CC' \) (spójrz na rysunek).}
{\( 90^{\circ} \)}{\( 60^{\circ} \)}{\( 45^{\circ} \)}{\( 72^{\circ} \)}{}{}}
\LANGcs{
\Question{
Pravidelný šestiboký hranol  \( ABCDEFA'B'C'D'E'F' \) má délku podstavné hrany  \( 4\,\mathrm{cm} \) a výšku \( 8\,\mathrm{cm} \). Určete odchylku přímek  \( A'D' \) a  \( CC' \) (viz obrázek).}
{\( 90^{\circ} \)}{\( 60^{\circ} \)}{\( 45^{\circ} \)}{\( 72^{\circ} \)}{}{}}
\LANGsk{
\Question{
Pravidelný šesťboký hranol  \( ABCDEFA'B'C'D'E'F' \) má dĺžku podstavnej hrany  \( 4\,\mathrm{cm} \) a výšku \( 8\,\mathrm{cm} \). Určte odchýlku priamok  \( A'D' \) a  \( CC' \) (viď obrázok).}
{\( 90^{\circ} \)}{\( 60^{\circ} \)}{\( 45^{\circ} \)}{\( 72^{\circ} \)}{}{}}
\LANGes{
\Question{
Sea \( ABCDEFA'B'C'D'E'F' \)un prisma hexagonal regular con la longitud de la arista básica de \( 4\,\mathrm{cm} \) y la altura de \( 8\,\mathrm{cm} \). Determina el ángulo entre las rectas  \( A'D' \) y \( CC' \) (mira la imagen).}
{\( 90^{\circ} \)}{\( 60^{\circ} \)}{\( 45^{\circ} \)}{\( 72^{\circ} \)}{}{}}
\End

\Data\ID{1103107004}
\Updated{2020-07-15 03:07:21}
\Level{C}
\SubArea{Lines and planes: distances and angles}
\IMG{1103107004.pdf}
\LANGen{
\Question{
Let \( ABCDEFA'B'C'D'E'F' \) be a regular hexagonal prism with the base edge length of \( 4\,\mathrm{cm} \) and the height of \( 8\,\mathrm{cm} \). Find the angle between the plane \( ABB' \) and the plane \( AEE' \) (see the picture).}
{\( 90^{\circ} \)}{\( 60^{\circ} \)}{\( 45^{\circ} \)}{\( 72^{\circ} \)}{}{}}
\LANGpl{
\Question{
Dany jest graniastosłup prawidłowy sześciokątny \( ABCDEFA'B'C'D'E'F' \), krawędź podstawy jest równa \( 4\,\mathrm{cm} \), a wysokość  \( 8\,\mathrm{cm} \). Oblicz miarę kąta pomiędzy płaszczyzną \( ABB' \) a \( AEE' \) (spójrz na rysunek).}
{\( 90^{\circ} \)}{\( 60^{\circ} \)}{\( 45^{\circ} \)}{\( 72^{\circ} \)}{}{}}
\LANGcs{
\Question{
Pravidelný šestiboký hranol  \( ABCDEFA'B'C'D'E'F' \) má délku podstavné hrany  \( 4\,\mathrm{cm} \) a výšku \( 8\,\mathrm{cm} \). Určete odchylku rovin  \( ABB' \) a  \( AEE' \) (viz obrázek).}
{\( 90^{\circ} \)}{\( 60^{\circ} \)}{\( 45^{\circ} \)}{\( 72^{\circ} \)}{}{}}
\LANGsk{
\Question{
Pravidelný šesťboký hranol  \( ABCDEFA'B'C'D'E'F' \) má dĺžku podstavnej hrany  \( 4\,\mathrm{cm} \) a výšku \( 8\,\mathrm{cm} \). Určte odchýlku rovín  \( ABB' \) a  \( AEE' \) (viď obrázok).}
{\( 90^{\circ} \)}{\( 60^{\circ} \)}{\( 45^{\circ} \)}{\( 72^{\circ} \)}{}{}}
\LANGes{
\Question{
Sea \( ABCDEFA'B'C'D'E'F' \) un prisma hexagonal regular con la longitud de la arista básica de \( 4\,\mathrm{cm} \) y la altura de  \( 8\,\mathrm{cm} \). Determina el ángulo entre el plano \( ABB' \) y el plano \( AEE' \) (mira la imagen).}
{\( 90^{\circ} \)}{\( 60^{\circ} \)}{\( 45^{\circ} \)}{\( 72^{\circ} \)}{}{}}
\End

\Data\ID{1103107005}
\Updated{2020-07-15 03:07:21}
\Level{C}
\SubArea{Lines and planes: distances and angles}
\IMG{1103107005_0.pdf}
\LANGen{
\Question{
Let \( ABCDEFA'B'C'D'E'F' \) be a regular hexagonal prism with the base edge length of \( 4\,\mathrm{cm} \) and the height of \( 8\,\mathrm{cm} \). Find the angle between the plane \( ADD' \) and the plane \( AEE' \) (see the picture).}
{\( 30^{\circ} \)}{\( 60^{\circ} \)}{\( 45^{\circ} \)}{\( 22.5^{\circ} \)}{}{}}
\LANGpl{
\Question{
Dany jest graniastosłup prawidłowy sześciokątny \( ABCDEFA'B'C'D'E'F' \), krawędź podstawy jest równa  \( 4\,\mathrm{cm} \), a wysokość  \( 8\,\mathrm{cm} \). Oblicz miarę kąta pomiędzy płaszczyzną \( ADD' \) i  \( AEE' \) (spójrz na rysunek).}
{\( 30^{\circ} \)}{\( 60^{\circ} \)}{\( 45^{\circ} \)}{\( 22{,}5^{\circ} \)}{}{}}
\LANGcs{
\Question{
Pravidelný šestiboký hranol  \( ABCDEFA'B'C'D'E'F' \) má délku podstavné hrany  \( 4\,\mathrm{cm} \) a výšku  \( 8\,\mathrm{cm} \). Určete odchylku rovin   \( ADD' \) a \( AEE' \) (viz obrázek).}
{\( 30^{\circ} \)}{\( 60^{\circ} \)}{\( 45^{\circ} \)}{\( 22{,}5^{\circ} \)}{}{}}
\LANGsk{
\Question{
Pravidelný šesťboký hranol  \( ABCDEFA'B'C'D'E'F' \) má dĺžku podstavnej hrany  \( 4\,\mathrm{cm} \) a výšku  \( 8\,\mathrm{cm} \). Určte odchýlku rovín   \( ADD' \) a \( AEE' \) (viď obrázok).}
{\( 30^{\circ} \)}{\( 60^{\circ} \)}{\( 45^{\circ} \)}{\( 22{,}5^{\circ} \)}{}{}}
\LANGes{
\Question{
Sea   \( ABCDEFA'B'C'D'E'F' \) un prisma hexagonal regular con la longitud de la arista básica de \( 4\,\mathrm{cm} \) y la altura de \( 8\,\mathrm{cm} \). Determina el ángulo entre el plano \( ADD' \) y el plano \( AEE' \) (mira la imagen).}
{\( 30^{\circ} \)}{\( 60^{\circ} \)}{\( 45^{\circ} \)}{\( 22.5^{\circ} \)}{}{}}
\End

\Data\ID{1103107006}
\Updated{2020-07-15 03:07:21}
\Level{C}
\SubArea{Lines and planes: distances and angles}
\IMG{1103107006.pdf}
\LANGen{
\Question{
Let \( ABCDEFA'B'C'D'E'F' \) be a regular hexagonal prism with the base edge length of \( 4\,\mathrm{cm} \) and the height of \( 4\sqrt6\,\mathrm{cm} \). Find the distance between the lines \( FA \) and \( C'D' \) (see the picture).}
{\( 12\,\mathrm{cm} \)}{\( 8\sqrt3\,\mathrm{cm} \)}{\( 4\sqrt6\,\mathrm{cm} \)}{\( 10\,\mathrm{cm} \)}{}{}}
\LANGpl{
\Question{
Dany jest graniastosłup prawidłowy sześciokątny \( ABCDEFA'B'C'D'E'F' \), krawędź podstawy jest równa  \( 4\,\mathrm{cm} \), a wysokość \( 4\sqrt6\,\mathrm{cm} \). Oblicz odległość pomiędzy prostymi  \( FA \) i \( C'D' \) (spójrz na rysunek).}
{\( 12\,\mathrm{cm} \)}{\( 8\sqrt3\,\mathrm{cm} \)}{\( 4\sqrt6\,\mathrm{cm} \)}{\( 10\,\mathrm{cm} \)}{}{}}
\LANGcs{
\Question{
Pravidelný šestiboký hranol  \( ABCDEFA'B'C'D'E'F' \) má délku podstavné hrany  \( 4\,\mathrm{cm} \) a výšku  \( 4\sqrt6\,\mathrm{cm} \). Určete vzdálenost přímek \( FA \) a \( C'D' \) (viz obrázek).}
{\( 12\,\mathrm{cm} \)}{\( 8\sqrt3\,\mathrm{cm} \)}{\( 4\sqrt6\,\mathrm{cm} \)}{\( 10\,\mathrm{cm} \)}{}{}}
\LANGsk{
\Question{
Pravidelný šesťboký hranol  \( ABCDEFA'B'C'D'E'F' \) má dĺžku podstavnej hrany  \( 4\,\mathrm{cm} \) a výšku  \( 4\sqrt6\,\mathrm{cm} \). Určte vzdialenosť priamok \( FA \) a \( C'D' \) (viď obrázok).}
{\( 12\,\mathrm{cm} \)}{\( 8\sqrt3\,\mathrm{cm} \)}{\( 4\sqrt6\,\mathrm{cm} \)}{\( 10\,\mathrm{cm} \)}{}{}}
\LANGes{
\Question{
Sea \( ABCDEFA'B'C'D'E'F' \) un prisma hexagonal regular con la longitud de la arista básica de \( 4\,\mathrm{cm} \) y la altura de \( 4\sqrt6\,\mathrm{cm} \). Determina la distancia entre las rectas \( FA \) y \( C'D' \) (mira la imagen).}
{\( 12\,\mathrm{cm} \)}{\( 8\sqrt3\,\mathrm{cm} \)}{\( 4\sqrt6\,\mathrm{cm} \)}{\( 10\,\mathrm{cm} \)}{}{}}
\End

\Data\ID{1103107007}
\Updated{2020-07-15 03:07:21}
\Level{C}
\SubArea{Lines and planes: distances and angles}
\IMG{1103107007.pdf}
\LANGen{
\Question{
Let \( ABCDEFA'B'C'D'E'F' \) be a regular hexagonal prism with the base edge length of \( 4\,\mathrm{cm} \) and the height of \( 6\,\mathrm{cm} \). Find the distance between the lines \( FA \) and \( E'B' \) (see the picture).}
{\( 4\sqrt3\,\mathrm{cm} \)}{\( 2\sqrt{13}\,\mathrm{cm} \)}{\( 2\sqrt{10}\,\mathrm{cm} \)}{\( 2\sqrt3\,\mathrm{cm} \)}{}{}}
\LANGpl{
\Question{
Dany jest graniastosłup prawidłowy sześciokątny \( ABCDEFA'B'C'D'E'F' \), krawędź podstawy jest równa  \( 4\,\mathrm{cm} \), a wysokość \( 6\,\mathrm{cm} \). Oblicz odległość pomiędzy prostymi \( FA \) i \( E'B' \) (spójrz na rysunek).}
{\( 4\sqrt3\,\mathrm{cm} \)}{\( 2\sqrt{13}\,\mathrm{cm} \)}{\( 2\sqrt{10}\,\mathrm{cm} \)}{\( 2\sqrt3\,\mathrm{cm} \)}{}{}}
\LANGcs{
\Question{
Pravidelný šestiboký hranol  \( ABCDEFA'B'C'D'E'F' \) má délku podstavné hrany  \( 4\,\mathrm{cm} \) a výšku \( 6\,\mathrm{cm} \). Určete vzdálenost přímek \( FA \) a \( E'B' \) (viz obrázek).}
{\( 4\sqrt3\,\mathrm{cm} \)}{\( 2\sqrt{13}\,\mathrm{cm} \)}{\( 2\sqrt{10}\,\mathrm{cm} \)}{\( 2\sqrt3\,\mathrm{cm} \)}{}{}}
\LANGsk{
\Question{
Pravidelný šesťboký hranol  \( ABCDEFA'B'C'D'E'F' \) má dĺžku podstavnej hrany  \( 4\,\mathrm{cm} \) a výšku \( 6\,\mathrm{cm} \). Určte vzdialenosť priamok \( FA \) a \( E'B' \) (viď obrázok).}
{\( 4\sqrt3\,\mathrm{cm} \)}{\( 2\sqrt{13}\,\mathrm{cm} \)}{\( 2\sqrt{10}\,\mathrm{cm} \)}{\( 2\sqrt3\,\mathrm{cm} \)}{}{}}
\LANGes{
\Question{
Sea\( ABCDEFA'B'C'D'E'F' \)un prisma hexagonal regular con la longitud de la arista básica de \( 4\,\mathrm{cm} \) y la altura de \( 6\,\mathrm{cm} \). Determina la distancia entre las rectas \( FA \) y \( E'B' \) (mira la imagen).}
{\( 4\sqrt3\,\mathrm{cm} \)}{\( 2\sqrt{13}\,\mathrm{cm} \)}{\( 2\sqrt{10}\,\mathrm{cm} \)}{\( 2\sqrt3\,\mathrm{cm} \)}{}{}}
\End

\Data\ID{1103107008}
\Updated{2020-07-15 03:07:21}
\Level{C}
\SubArea{Lines and planes: distances and angles}
\IMG{1103107008.pdf}
\LANGen{
\Question{
A regular hexagonal prism \( ABCDEFA'B'C'D'E'F' \) has the  base edge length of \( 4\,\mathrm{cm} \) and the height of \( 6\,\mathrm{cm} \). Let \( S \) be the middle of the top face \( A'B'C'D'E'F' \). Find the distance between the point \( S \) and the line \( BC \)  (see the picture).}
{\( 4\sqrt3\,\mathrm{cm} \)}{\( 2\sqrt{13}\,\mathrm{cm} \)}{\( 2\sqrt{21}\,\mathrm{cm} \)}{\( 4\sqrt{21}\,\mathrm{cm} \)}{}{}}
\LANGpl{
\Question{
Dany jest graniastosłup prawidłowy sześciokątny\( ABCDEFA'B'C'D'E'F' \), krawędź podstawy jest równa  \( 4\,\mathrm{cm} \), a wysokość \( 6\,\mathrm{cm} \). \( S \) to środek górnej podstawy  \( A'B'C'D'E'F' \). Oblicz odległość pomiędzy punktem \( S \) a prostą  \( BC \)  (spójrz na rysunek).}
{\( 4\sqrt3\,\mathrm{cm} \)}{\( 2\sqrt{13}\,\mathrm{cm} \)}{\( 2\sqrt{21}\,\mathrm{cm} \)}{\( 4\sqrt{21}\,\mathrm{cm} \)}{}{}}
\LANGcs{
\Question{
Pravidelný šestiboký hranol  \( ABCDEFA'B'C'D'E'F' \) má délku podstavné hrany  \( 4\,\mathrm{cm} \) a výšku \( 6\,\mathrm{cm} \). Vypočítejte vzdálenost bodu \( S \) od přímky \( BC \). Bod  \( S \) je střed horní podstavy \( A'B'C'D'E'F' \)  (viz obrázek).}
{\( 4\sqrt3\,\mathrm{cm} \)}{\( 2\sqrt{13}\,\mathrm{cm} \)}{\( 2\sqrt{21}\,\mathrm{cm} \)}{\( 4\sqrt{21}\,\mathrm{cm} \)}{}{}}
\LANGsk{
\Question{
Pravidelný šesťboký hranol  \( ABCDEFA'B'C'D'E'F' \) má dĺžku podstavnej hrany  \( 4\,\mathrm{cm} \) a výšku \( 6\,\mathrm{cm} \). Vypočítajte vzdialenosť bodu \( S \) od priamky \( BC \). Bod  \( S \) je stred hornej podstavy \( A'B'C'D'E'F' \)  (viď obrázok).}
{\( 4\sqrt3\,\mathrm{cm} \)}{\( 2\sqrt{13}\,\mathrm{cm} \)}{\( 2\sqrt{21}\,\mathrm{cm} \)}{\( 4\sqrt{21}\,\mathrm{cm} \)}{}{}}
\LANGes{
\Question{
Un prisma hexagonal regular \( ABCDEFA'B'C'D'E'F' \) tiene la longitud de la arista básica de  \( 4\,\mathrm{cm} \) y la altura de \( 6\,\mathrm{cm} \). Sea \( S \)  el centro de la cara superior \( A'B'C'D'E'F' \). Determina la distancia entre el punto \( S \) y la recta \( BC \)  (mira la imagen).}
{\( 4\sqrt3\,\mathrm{cm} \)}{\( 2\sqrt{13}\,\mathrm{cm} \)}{\( 2\sqrt{21}\,\mathrm{cm} \)}{\( 4\sqrt{21}\,\mathrm{cm} \)}{}{}}
\End

\Data\ID{1103107009}
\Updated{2020-07-15 03:07:21}
\Level{C}
\SubArea{Lines and planes: distances and angles}
\IMG{1103107009_0.pdf}
\LANGen{
\Question{
Let \( ABCDEFA'B'C'D'E'F' \) be a regular hexagonal prism with the base edge length of \( 4\,\mathrm{cm} \) and the height of \( 8\,\mathrm{cm} \). Find the distance between the points \( F' \) and \( C \) (see the picture).}
{\( 8\sqrt2\,\mathrm{cm} \)}{\( 4\sqrt5\,\mathrm{cm} \)}{\( 4\sqrt7\,\mathrm{cm} \)}{\( 4\sqrt2\,\mathrm{cm} \)}{}{}}
\LANGpl{
\Question{
Dany jest graniastosłup prawidłowy sześciokątny \( ABCDEFA'B'C'D'E'F' \), krawędź podstawy jest równa  \( 4\,\mathrm{cm} \), a wysokość  \( 8\,\mathrm{cm} \). Oblicz odległość pomiędzy punktami \( F' \) i \( C \) (spójrz na rysunek).}
{\( 8\sqrt2\,\mathrm{cm} \)}{\( 4\sqrt5\,\mathrm{cm} \)}{\( 4\sqrt7\,\mathrm{cm} \)}{\( 4\sqrt2\,\mathrm{cm} \)}{}{}}
\LANGcs{
\Question{
Pravidelný šestiboký hranol  \( ABCDEFA'B'C'D'E'F' \) má délku podstavné hrany  \( 4\,\mathrm{cm} \) a výšku \( 8\,\mathrm{cm} \). Vypočítejte vzdálenost bodů \( F' \) a \( C \) (viz obrázek).}
{\( 8\sqrt2\,\mathrm{cm} \)}{\( 4\sqrt5\,\mathrm{cm} \)}{\( 4\sqrt7\,\mathrm{cm} \)}{\( 4\sqrt2\,\mathrm{cm} \)}{}{}}
\LANGsk{
\Question{
Pravidelný šesťboký hranol  \( ABCDEFA'B'C'D'E'F' \) má dĺžku podstavnej hrany  \( 4\,\mathrm{cm} \) a výšku \( 8\,\mathrm{cm} \). Vypočítajte vzdialenosť bodov \( F' \) a \( C \) (viď obrázok).}
{\( 8\sqrt2\,\mathrm{cm} \)}{\( 4\sqrt5\,\mathrm{cm} \)}{\( 4\sqrt7\,\mathrm{cm} \)}{\( 4\sqrt2\,\mathrm{cm} \)}{}{}}
\LANGes{
\Question{
Sea \( ABCDEFA'B'C'D'E'F' \) un prisma hexagonal regular con la longitud de la arista básica de \( 4\,\mathrm{cm} \) y la altura de \( 8\,\mathrm{cm} \). Determina la distancia entre los puntos \( F' \) y \( C \) (mira la imagen).}
{\( 8\sqrt2\,\mathrm{cm} \)}{\( 4\sqrt5\,\mathrm{cm} \)}{\( 4\sqrt7\,\mathrm{cm} \)}{\( 4\sqrt2\,\mathrm{cm} \)}{}{}}
\End

\Data\ID{1103107010}
\Updated{2020-07-15 03:07:21}
\Level{C}
\SubArea{Lines and planes: distances and angles}
\IMG{1103107010_0.pdf}
\LANGen{
\Question{
Let \( ABCDEFA'B'C'D'E'F' \) be a regular hexagonal prism with the base edge length of \( 4\,\mathrm{cm} \) and the height of \( 8\,\mathrm{cm} \). Find the distance between the points \( E' \) and \( C \) (see the picture).}
{\( 4\sqrt7\,\mathrm{cm} \)}{\( 8\sqrt2\,\mathrm{cm} \)}{\( 8+4\sqrt3\,\mathrm{cm} \)}{\( 16\sqrt7\,\mathrm{cm} \)}{}{}}
\LANGpl{
\Question{
Dany jest graniastosłup prawidłowy sześciokątny \( ABCDEFA'B'C'D'E'F' \), krawędź podstawy jest równa  \( 4\,\mathrm{cm} \), a wysokość \( 8\,\mathrm{cm} \). Oblicz odległość pomiędzy punktami  \( E' \) i \( C \) (spójrz na rysunek).}
{\( 4\sqrt7\,\mathrm{cm} \)}{\( 8\sqrt2\,\mathrm{cm} \)}{\( 8+4\sqrt3\,\mathrm{cm} \)}{\( 16\sqrt7\,\mathrm{cm} \)}{}{}}
\LANGcs{
\Question{
Pravidelný šestiboký hranol  \( ABCDEFA'B'C'D'E'F' \) má délku podstavné hrany  \( 4\,\mathrm{cm} \) a výšku \( 8\,\mathrm{cm} \). Vypočítejte vzdálenost bodů  \( E' \) a \( C \) (viz obrázek).}
{\( 4\sqrt7\,\mathrm{cm} \)}{\( 8\sqrt2\,\mathrm{cm} \)}{\( 8+4\sqrt3\,\mathrm{cm} \)}{\( 16\sqrt7\,\mathrm{cm} \)}{}{}}
\LANGsk{
\Question{
Pravidelný šesťboký hranol  \( ABCDEFA'B'C'D'E'F' \) má dĺžku podstavnej hrany  \( 4\,\mathrm{cm} \) a výšku \( 8\,\mathrm{cm} \). Vypočítajte vzdialenosť bodov  \( E' \) a \( C \) (viď obrázok).}
{\( 4\sqrt7\,\mathrm{cm} \)}{\( 8\sqrt2\,\mathrm{cm} \)}{\( 8+4\sqrt3\,\mathrm{cm} \)}{\( 16\sqrt7\,\mathrm{cm} \)}{}{}}
\LANGes{
\Question{
Sea \( ABCDEFA'B'C'D'E'F' \) un prisma hexagonal regular con la longitud de la arista básica de  \( 4\,\mathrm{cm} \)  y la altura de  \( 8\,\mathrm{cm} \).  Determina la distancia entre los puntos \( E' \) y \( C \) (mira la imagen).}
{\( 4\sqrt7\,\mathrm{cm} \)}{\( 8\sqrt2\,\mathrm{cm} \)}{\( 8+4\sqrt3\,\mathrm{cm} \)}{\( 16\sqrt7\,\mathrm{cm} \)}{}{}}
\End

\Data\ID{1103107011}
\Updated{2020-07-15 03:07:21}
\Level{C}
\SubArea{Lines and planes: distances and angles}
\IMG{1103107011.pdf}
\LANGen{
\Question{
Let \( ABCDEFA'B'C'D'E'F' \) be a regular hexagonal prism with the base edge length of \( 4\,\mathrm{cm} \) and the height of \( 8\,\mathrm{cm} \). Find the angle between the line \( FC' \) and the base plane \( ABC \) (see the picture).}
{\( 45^{\circ} \)}{\( 60^{\circ} \)}{\( 30^{\circ} \)}{\( 72^{\circ} \)}{}{}}
\LANGpl{
\Question{
Dany jest graniastosłup prawidłowy sześciokątny \( ABCDEFA'B'C'D'E'F' \), krawędź podstawy jest równa \( 4\,\mathrm{cm} \), a wysokość  \( 8\,\mathrm{cm} \). Oblicz miarę kąta pomiędzy prostą \( FC' \) a płaszczyzną podstawy  \( ABC \) (spójrz na rysunek).}
{\( 45^{\circ} \)}{\( 60^{\circ} \)}{\( 30^{\circ} \)}{\( 72^{\circ} \)}{}{}}
\LANGcs{
\Question{
Pravidelný šestiboký hranol  \( ABCDEFA'B'C'D'E'F' \) má délku podstavné hrany  \( 4\,\mathrm{cm} \) a výšku \( 8\,\mathrm{cm} \). Určete odchylku přímky \( FC' \) a roviny podstavy \( ABC \) (viz obrázek).}
{\( 45^{\circ} \)}{\( 60^{\circ} \)}{\( 30^{\circ} \)}{\( 72^{\circ} \)}{}{}}
\LANGsk{
\Question{
Pravidelný šesťboký hranol  \( ABCDEFA'B'C'D'E'F' \) má dĺžku podstavnej hrany  \( 4\,\mathrm{cm} \) a výšku \( 8\,\mathrm{cm} \). Určte odchýlku priamky \( FC' \) a roviny podstavy \( ABC \) (viď obrázok).}
{\( 45^{\circ} \)}{\( 60^{\circ} \)}{\( 30^{\circ} \)}{\( 72^{\circ} \)}{}{}}
\LANGes{
\Question{
Sea \( ABCDEFA'B'C'D'E'F' \) un prisma hexagonal regular con la longitud de la arista básica de  \( 4\,\mathrm{cm} \)  y la altura de  \( 8\,\mathrm{cm} \). Determina el ángulo entre las rectas \( FC' \) y el plano de la base\( ABC \) (mira la imagen).}
{\( 45^{\circ} \)}{\( 60^{\circ} \)}{\( 30^{\circ} \)}{\( 72^{\circ} \)}{}{}}
\End

\Data\ID{1103107012}
\Updated{2020-07-15 03:07:21}
\Level{C}
\SubArea{Lines and planes: distances and angles}
\IMG{1103107012.pdf}
\LANGen{
\Question{
Let \( ABCDEFA'B'C'D'E'F' \) be a regular hexagonal prism with the base edge length of \( 4\,\mathrm{cm} \) and the height of \( 8\,\mathrm{cm} \). Find the angle between the plane \( ADD' \) and the plane \( CDD' \) (see the picture).}
{\( 60^{\circ} \)}{\( 45^{\circ} \)}{\( 90^{\circ} \)}{\( 72^{\circ} \)}{}{}}
\LANGpl{
\Question{
Dany jest graniastosłup prawidłowy sześciokątny \( ABCDEFA'B'C'D'E'F' \), krawędź podstawy jest równa \( 4\,\mathrm{cm} \), a wysokość \( 8\,\mathrm{cm} \). Oblicz miarę kąta pomiędzy płaszczyzną \( ADD' \) a \( CDD' \) (spójrz na rysunek).}
{\( 60^{\circ} \)}{\( 45^{\circ} \)}{\( 90^{\circ} \)}{\( 72^{\circ} \)}{}{}}
\LANGcs{
\Question{
Pravidelný šestiboký hranol  \( ABCDEFA'B'C'D'E'F' \) má délku podstavné hrany  \( 4\,\mathrm{cm} \) a výšku \( 8\,\mathrm{cm} \). Určete odchylku rovin \( ADD' \) a \( CDD' \) (viz obrázek).}
{\( 60^{\circ} \)}{\( 45^{\circ} \)}{\( 90^{\circ} \)}{\( 72^{\circ} \)}{}{}}
\LANGsk{
\Question{
Pravidelný šesťboký hranol  \( ABCDEFA'B'C'D'E'F' \) má dĺžku podstavnej hrany  \( 4\,\mathrm{cm} \) a výšku \( 8\,\mathrm{cm} \). Určte odchýlku rovín \( ADD' \) a \( CDD' \) (viď obrázok).}
{\( 60^{\circ} \)}{\( 45^{\circ} \)}{\( 90^{\circ} \)}{\( 72^{\circ} \)}{}{}}
\LANGes{
\Question{
Sea \( ABCDEFA'B'C'D'E'F' \) un prisma hexagonal regular con la longitud de la arista básica de  \( 4\,\mathrm{cm} \)  y la altura de  \( 8\,\mathrm{cm} \). Determina el ángulo entre el plano \( ADD' \) y el plano \( CDD' \) (mira la imagen).}
{\( 60^{\circ} \)}{\( 45^{\circ} \)}{\( 90^{\circ} \)}{\( 72^{\circ} \)}{}{}}
\End

\Data\ID{1103107013}
\Updated{2020-07-15 03:07:21}
\Level{C}
\SubArea{Lines and planes: distances and angles}
\IMG{1103107013.pdf}
\LANGen{
\Question{
Let \( ABCDEFA'B'C'D'E'F' \) be a regular hexagonal prism with the base edge length of \( 4\,\mathrm{cm} \) and the height of \( 8\,\mathrm{cm} \). Find the angle between the plane \( BCC' \) and the plane \( CDD' \) (see the picture).}
{\( 60^{\circ} \)}{\( 120^{\circ} \)}{\( 90^{\circ} \)}{\( 72^{\circ} \)}{}{}}
\LANGpl{
\Question{
Dany jest graniastosłup prawidłowy sześciokątny \( ABCDEFA'B'C'D'E'F' \), krawędź podstawy jest równa  \( 4\,\mathrm{cm} \), a wysokość \( 8\,\mathrm{cm} \). Oblicz miarę kąta pomiędzy płaszczyzną \( BCC' \) a \( CDD' \) (spójrz na rysunek).}
{\( 60^{\circ} \)}{\( 120^{\circ} \)}{\( 90^{\circ} \)}{\( 72^{\circ} \)}{}{}}
\LANGcs{
\Question{
Pravidelný šestiboký hranol  \( ABCDEFA'B'C'D'E'F' \) má délku podstavné hrany  \( 4\,\mathrm{cm} \) a výšku \( 8\,\mathrm{cm} \). Určete odchylku rovin \( BCC' \) a \( CDD' \) (viz obrázek).}
{\( 60^{\circ} \)}{\( 120^{\circ} \)}{\( 90^{\circ} \)}{\( 72^{\circ} \)}{}{}}
\LANGsk{
\Question{
Pravidelný šesťboký hranol  \( ABCDEFA'B'C'D'E'F' \) má dĺžku podstavnej hrany  \( 4\,\mathrm{cm} \) a výšku \( 8\,\mathrm{cm} \). Určte odchýlku rovín \( BCC' \) a \( CDD' \) (viď obrázok).}
{\( 60^{\circ} \)}{\( 120^{\circ} \)}{\( 90^{\circ} \)}{\( 72^{\circ} \)}{}{}}
\LANGes{
\Question{
Sea \( ABCDEFA'B'C'D'E'F' \) un prisma hexagonal regular con la longitud de la arista básica de  \( 4\,\mathrm{cm} \)  y la altura de  \( 8\,\mathrm{cm} \). Determina el ángulo entre el plano \( BCC' \) y el plano \( CDD' \) (mira la imagen).}
{\( 60^{\circ} \)}{\( 120^{\circ} \)}{\( 90^{\circ} \)}{\( 72^{\circ} \)}{}{}}
\End

\Data\ID{1103107014}
\Updated{2020-07-15 03:07:21}
\Level{C}
\SubArea{Lines and planes: distances and angles}
\IMG{1103107014.pdf}
\LANGen{
\Question{
Let \( ABCDEFA'B'C'D'E'F' \) be a regular hexagonal prism with the base edge length of \( 4\,\mathrm{cm} \) and the height of \( 8\,\mathrm{cm} \). Find the angle between the line \( BA’ \) and the plane \( AEE’ \) (see the picture). Round the result to two decimal places.}
{\( 26.57^{\circ} \)}{\( 63.43^{\circ} \)}{\( 30^{\circ} \)}{\( 22.5^{\circ} \)}{}{}}
\LANGpl{
\Question{
Dany jest graniastosłup prawidłowy sześciokątny \( ABCDEFA'B'C'D'E'F' \), krawędź podstawy jest równa \( 4\,\mathrm{cm} \), a wysokość  \( 8\,\mathrm{cm} \). Oblicz miarę kąta pomiędzy prostą \( BA’ \) a płaszczyzną  \( AEE’ \) (spójrz na rysunek). Zaokrągli wynik do dwóch miejsc po przecinku.}
{\( 26{,}57^{\circ} \)}{\( 63{,}43^{\circ} \)}{\( 30^{\circ} \)}{\( 22{,}5^{\circ} \)}{}{}}
\LANGcs{
\Question{
Pravidelný šestiboký hranol  \( ABCDEFA'B'C'D'E'F' \) má délku podstavné hrany  \( 4\,\mathrm{cm} \) a výšku \( 8\,\mathrm{cm} \). Určete odchylku přímky \( BA’ \) a roviny \( AEE’ \) (viz obrázek). Výsledek zaokrouhlete na dvě desetinná místa.}
{\( 26{,}57^{\circ} \)}{\( 63{,}43^{\circ} \)}{\( 30^{\circ} \)}{\( 22{,}5^{\circ} \)}{}{}}
\LANGsk{
\Question{
Pravidelný šesťboký hranol  \( ABCDEFA'B'C'D'E'F' \) má dĺžku podstavnej hrany  \( 4\,\mathrm{cm} \) a výšku \( 8\,\mathrm{cm} \). Určte odchýlku priamky \( BA’ \) a roviny \( AEE’ \) (viď obrázok). Výsledok zaokrúhlite na dve desatinné miesta.}
{\( 26{,}57^{\circ} \)}{\( 63{,}43^{\circ} \)}{\( 30^{\circ} \)}{\( 22{,}5^{\circ} \)}{}{}}
\LANGes{
\Question{
Sea \( ABCDEFA'B'C'D'E'F' \) un prisma hexagonal regular con la longitud de la arista básica de  \( 4\,\mathrm{cm} \)  y la altura de  \( 8\,\mathrm{cm} \). Determina el ángulo entre la recta \( BA’ \) y el plano \( AEE’ \) (mira la imagen). Redondea a dos posiciones decimales.}
{\( 26.57^{\circ} \)}{\( 63.43^{\circ} \)}{\( 30^{\circ} \)}{\( 22.5^{\circ} \)}{}{}}
\End

\Data\ID{9000120304}
\Updated{2020-07-15 03:07:37}
\Level{C}
\SubArea{Lines and planes: distances and angles}
\IMG{001203_04-figure0.pdf}
\LANGen{
\Question{
The side of a regular hexagonal prism
\(ABCDEFA'B'C'D'E'F'\) is
\(a = 3\, \mathrm{cm}\) and the height
\(v = 8\, \mathrm{cm}\). Find the length
of the diagonal \(AD'\).}
{\(10\, \mathrm{cm}\)}{\(\sqrt{73}\, \mathrm{cm}\)}{\(\sqrt{82}\, \mathrm{cm}\)}{\(2\sqrt{8}\, \mathrm{cm}\)}{\(2\sqrt{6}\, \mathrm{cm}\)}{}}
\LANGpl{
\Question{
Bok graniastosłupa prawidłowego sześciokątnego 
\(ABCDEFA'B'C'D'E'F'\) jest równy
\(a = 3\, \mathrm{cm}\), a jego wysokość wynosi
\(v = 8\, \mathrm{cm}\). Oblicz długość przekątnej 
 \(AD'\).}
{\(10\, \mathrm{cm}\)}{\(\sqrt{73}\, \mathrm{cm}\)}{\(\sqrt{82}\, \mathrm{cm}\)}{\(2\sqrt{8}\, \mathrm{cm}\)}{\(2\sqrt{6}\, \mathrm{cm}\)}{}}
\LANGcs{
\Question{
V pravidelném šestibokém hranolu
\(ABCDEFA'B'C'D'E'F'\) je délka
podstavné hrany \(a = 3\, \mathrm{cm}\),
výška \(v = 8\, \mathrm{cm}\). Délka úhlopříčky \(AD'\) 
je rovna:}
{\(10\, \mathrm{cm}\)}{\(\sqrt{73}\, \mathrm{cm}\)}{\(\sqrt{82}\, \mathrm{cm}\)}{\(2\sqrt{8}\, \mathrm{cm}\)}{\(2\sqrt{6}\, \mathrm{cm}\)}{}}
\LANGsk{
\Question{
V pravidelnom šesťbokom hranole
\(ABCDEFA'B'C'D'E'F'\) je dĺžka
hrany podstavy \(a = 3\, \mathrm{cm}\),
výška \(v = 8\, \mathrm{cm}\). Dĺžka uhlopriečky \(AD'\) 
je rovná:}
{\(10\, \mathrm{cm}\)}{\(\sqrt{73}\, \mathrm{cm}\)}{\(\sqrt{82}\, \mathrm{cm}\)}{\(2\sqrt{8}\, \mathrm{cm}\)}{\(2\sqrt{6}\, \mathrm{cm}\)}{}}
\LANGes{
\Question{
La arista de un prisma hexagonal regular 
\(ABCDEFA'B'C'D'E'F'\) es
\(a = 3\, \mathrm{cm}\) y la altura es
\(v = 8\, \mathrm{cm}\).  Determina la longitud de la diagonal \(AD'\).}
{\(10\, \mathrm{cm}\)}{\(\sqrt{73}\, \mathrm{cm}\)}{\(\sqrt{82}\, \mathrm{cm}\)}{\(2\sqrt{8}\, \mathrm{cm}\)}{\(2\sqrt{6}\, \mathrm{cm}\)}{}}
\End

\Data\ID{9000120305}
\Updated{2021-07-24 10:07:04}
\Level{C}
\SubArea{Lines and planes: distances and angles}
\IMG{001203_05-figure0.pdf}
\LANGen{
\Question{
The side of a regular hexagonal prism
\(ABCDEFA'B'C'D'E'F'\) is
\(a = 3\, \mathrm{cm}\) and the height
\(v = 8\, \mathrm{cm}\). Find the angle
between the diagonal \(AD'\) 
and the base plane \(ABC\) 
(round your result to the nearest degree).}
{\(53^{\circ }\)}{\(37^{\circ }\)}{\(45^{\circ }\)}{\(61^{\circ }\)}{\(72^{\circ }\)}{}}
\LANGpl{
\Question{
Dany jest graniastosłup prawidłowy sześciokątny
\(ABCDEFA'B'C'D'E'F'\) o boku
\(a = 3\, \mathrm{cm}\) i wysokości
\(v = 8\, \mathrm{cm}\). Oblicz kąt 
między przekątną  \(AD'\) 
a płaszczyzną podstawy \(ABC\) 
(zaokrągli wynik do pełnych stopni).}
{\(53^{\circ }\)}{\(37^{\circ }\)}{\(45^{\circ }\)}{\(61^{\circ }\)}{\(72^{\circ }\)}{}}
\LANGcs{
\Question{
V pravidelném šestibokém hranolu
\(ABCDEFA'B'C'D'E'F'\) je délka
podstavné hrany \(a = 3\, \mathrm{cm}\),
výška \(v = 8\, \mathrm{cm}\). Najděte odchylku úhlopříčky \(AD'\) 
od roviny podstavy \(ABC\) . Výsledek zaokrouhlete na celé stupně.}
{\(53^{\circ }\)}{\(37^{\circ }\)}{\(45^{\circ }\)}{\(61^{\circ }\)}{\(72^{\circ }\)}{}}
\LANGsk{
\Question{
V pravidelnom šesťbokom hranole
\(ABCDEFA'B'C'D'E'F'\) je dĺžka
hrany podstavy \(a = 3\, \mathrm{cm}\),
výška \(v = 8\, \mathrm{cm}\). Odchýlka uhlopriečky \(AD'\) 
od roviny podstavy \(ABC\) 
je rovná (výsledok zaokrúhlite na celé stupne):}
{\(53^{\circ }\)}{\(37^{\circ }\)}{\(45^{\circ }\)}{\(61^{\circ }\)}{\(72^{\circ }\)}{}}
\LANGes{
\Question{
La arista de un prisma hexagonal regular
\(ABCDEFA'B'C'D'E'F'\) es
\(a = 3\, \mathrm{cm}\) y la altura es
\(v = 8\, \mathrm{cm}\). Determina el ángulo
entre la diagonal \(AD'\) 
y el plano de la base \(ABC\). 
Redondea la respuesta final al grado más cercano.}
{\(53^{\circ }\)}{\(37^{\circ }\)}{\(45^{\circ }\)}{\(61^{\circ }\)}{\(72^{\circ }\)}{}}
\End
